\documentclass[doublespacing]{utdthesis}
% For one-and-a-half spacing, use: \documentclass[halfspacing]{utdthesis}

%%% Load any desired packages in the space below.
%%% Warning: Do not load packages that change the margins, headers, or footers!
%%%
% Optional: If you want to use Times as your font, load it here.  Note that
% although package "times" should work, it may not be the best choice.  Newer
% LaTeX distributions offer "mathptmx" and "newtxtext,newtxmath" as superior
% replacements.  You should find out which is best for your LaTeX.  (If this
% sounds confusing, you probably shouldn't try to change the font to Times.)
%\usepackage{times}
%
% Optional: If your LaTeX has microtype, use it to improve text quality:
\usepackage{microtype}
%
% Recommended: If your dissertation contains math, use the AMS packages:
\usepackage{amsmath,amssymb,amsthm}
%
% Recommended: If your dissertation needs embedded graphics, use graphicx:
\usepackage{graphicx}
%
% Recommended: If your bibliography contains web page URLs, the url package
% improves their appearance (e.g., better line breaking):
\usepackage{url}

\usepackage{booktabs,threeparttable}


\usepackage{subcaption}
% "natbib" package.  (Use other citation packages at your own risk; not all
% are flexible enough to meet UTD's requirements.)  If you wish to use numeric
% citations, change "authoryear" to "numbers" below.  Use the "chicago" BibTeX
% style, which most closely matches the Turabian formatting required by UTD.
% UTD mandates a blank line between each pair of bibliography entries, so set
% \bibsep as shown below.  Finally, if you are accustomed to using \cite as
% your citation macro, point it to natbib's \citep macro as shown.

% -------------------------------
% Bibliography setup (UTD template default)
% -------------------------------
\usepackage[authoryear]{natbib}
\bibliographystyle{chicago}
\setlength{\bibsep}{12pt plus 1pt minus 1pt}
\let\cite=\citep
%
% Required: If you have any wide tables or figures that need to be typeset
% in landscape, use the rotating package:
\usepackage{rotating}
%
% Optional: line spacing controls (used here to force double spacing in the
% General Introduction chapter body)
\usepackage{setspace}
%
% Optional: If you use hyperref to auto-generate hyperlinks, always load it
% LAST since it modifies everything else.  In addition, only load hyperref if
% you use pdftex or pdflatex to generate PDFs directly.  Do NOT use it if you
% use plain tex or latex to generate a DVI file.  (If you are generating DVI
% files which you then convert to PDF, you should seriously consider switching
% to pdflatex.  The DVI format loses information because it cannot support
% modern PDF document features.  Using pdflatex to generate PDFs directly
% therefore results in documents of significantly higher quality.)
\usepackage{ifpdf}
\ifpdf
\usepackage[hidelinks]{hyperref}
\hypersetup{pdfborder={0 0 0}}
\fi





%%% End of packages.

%%% Define all your personal macros here (if you have any).
%
\providecommand{\hyperref}[2][]{#2}



\newenvironment{exampleclasscode}
 {\parindent=1cm\vskip0pt plus2pt minus0pt\begin{verse}}
 {\end{verse}\vskip0pt plus2pt minus0pt}
%
%%% End of personal macro definitions.


%%% The following definitions MUST come before the document begins.
%
\author{Yichao Jin}
\title{Intertemporal Preferences and Policy Design for Vaccination: A Behavioral Public Health Approach}
\thesistype{Dissertation}
\degreefull{Doctor of Philosophy}
\degreeabbr{Ph.D.}
\subject{Public Policy and Political Economy}
\graduationmonth{March}
\graduationyear{2026}
\prevdegrees{B.A., M.A.}


% List committee members in order.  Mark chairpersons with a "*":
\committeemember*{Dohyeong Kim}
\committeemember{Richard Scotch}
\committeemember{Soyoung Kwon}
\committeemember{Khai Chiong}
%
%%% End of definitions.


%%% Beginning of actual thesis document.

\begin{document}

\frontmatter

\signaturepage
\copyrightpage{2026} % optional

\begin{dedication} % optional
This dissertation \\
is dedicated to my family and mentors, \\
whose support made this research possible.
\end{dedication}

\maketitle

\begin{acks}{April 2026} % date when thesis first submitted to committee

\end{acks}

\begin{abstract}
Timely vaccination is a cornerstone of global public health and epidemic control, yet real-world vaccination behavior is frequently characterized by significant delays rather than simple refusal. Many individuals postpone vaccination even when vaccines are available, remaining exposed to infection risk and weakening the effective coverage of immunization programs. This dissertation conceptualizes vaccination timing as a critical intertemporal choice problem, where individuals trade off immediate behavioral costs against future health protection. Using a large-scale discrete choice experiment (DCE) fielded in Wuhan, China ($N=1,027$), I quantify the behavioral cost of waiting and develop a suite of policy-impact evaluation tools to optimize vaccination program design.

Essay~1 establishes the behavioral-economic foundation by estimating structural intertemporal preferences. I test for present-biased (hyperbolic) discounting in the valuation of delayed protection and introduce the Behavioral Delay Multiplier (BDM) as a novel metric for quantifying the "psychological time tax" of waiting. Essay~2 translates these behavioral preferences into a policy-ready monetary metric, deriving the Marginal Willingness to Accept (MWTA) for waiting time—defined as the "shadow price of waiting." The results reveal substantial heterogeneity in delay valuation across socio-demographics, biological vulnerability, and institutional trust. I operationalize these estimates through counterfactual simulations and iso-acceptance frontiers that guide incentive design. Essay~3 integrates structural parameters and monetary valuations into a unified policy framework. I derive equity-sensitive decision rules for choosing between logistical speed-ups (capacity investment) and monetary compensation (incentive provision) under explicit budget constraints.

Together, this dissertation makes three core contributions to behavioral public health: (i) providing structural evidence on how populations discount delayed health investments, (ii) introducing the BDM and MWTA as intuitive tools for policy-impact evaluation, and (iii) offering a transparent decision framework for designing immunization campaigns that account for behavioral heterogeneity and distributional equity. The findings demonstrate that time sensitivity is a first-order determinant of vaccine uptake, and that targeting interventions based on structural intertemporal preferences can significantly improve the efficiency and equity of public health responses.

\bigskip
\noindent\textbf{Keywords:} Vaccination timing; Behavioral public health; Intertemporal choice; Present bias; Discrete choice experiment; Policy impact evaluation; Behavioral Delay Multiplier (BDM); Shadow price of waiting; Health economics.
\end{abstract}

\tableofcontents
\listoffigures % required if you have any figures
\listoftables % required if you have any tables

\mainmatter

\chapter{General Introduction}
\doublespacing 

Vaccination is widely recognized as a cornerstone of global public health and infectious-disease prevention, yet real-world vaccination behavior is complex, characterized not merely by binary uptake but by significant heterogeneity in timing. While outright refusal is often the focus of public discourse, \textbf{delayed uptake} represents a critical but distinct challenge for behavioral public health: it leaves individuals exposed to infection risk for longer periods, thereby reducing the effective coverage of immunization programs and delaying herd immunity \cite{MacDonald2015_SAGE, Larson2014, Dai2021_Nature}. Unlike permanent refusal, delay often reflects a sophisticated—albeit potentially biased—intertemporal trade-off.

To explain why delays in vaccination occur—and why the same delay may impose very different behavioral burdens across individuals—this dissertation draws on three complementary strands of scholarship. The first concerns \textbf{behavioral public health and intertemporal choice}, which predicts that immediate behavioral costs (e.g., logistical friction, time, side effects) can loom disproportionately large relative to delayed and uncertain health benefits \cite{Frederick2002, Laibson1997, ODonoghueRabin1999}. The second examines the \textbf{structural determinants of vaccination behavior}, highlighting how trust, perceived risk, and structural barriers shape uptake in practice \cite{Brewer2017, Milkman2021_PNAS, Schmid2017}. The third provides the \textbf{stated-preference and discrete choice experiment (DCE)} toolkit for quantifying these trade-offs and performing policy impact evaluations in a rigorous and internally consistent framework \cite{Louviere2000, Train2009, Soekhai2019_Pharmacoeconomics}. The next section briefly synthesizes these literatures and clarifies how they motivate the dissertation’s measurement strategy and research questions.

\section{Literature and Conceptual Foundations}

Delays in vaccination are increasingly recognized as conceptually distinct from outright refusal. Whereas refusal implies a stable decision not to vaccinate, delayed uptake reflects a timing choice that extends the period of biological vulnerability and can weaken the effective coverage of immunization programs \cite{MacDonald2015_SAGE, Larson2014}. Explaining why individuals postpone protection therefore requires integrating insights from (i) intertemporal choice and discounting, (ii) behavioral determinants of vaccination, and (iii) policy impact evaluation methods that can credibly quantify trade-offs involving waiting time.

\subsubsection{Intertemporal choice, discounting, and preventive health}

A central premise of this dissertation is that vaccination timing can be understood as an \textbf{intertemporal health investment}. In the health-capital tradition, individuals allocate effort and resources today to improve future health states, implying that preventive actions are shaped by the valuation of future outcomes relative to present costs \cite{Grossman1972, Cropper1977}. Yet a broad body of work in behavioral public health shows that individuals often discount delayed outcomes and may do so in a time-inconsistent manner, placing disproportionate weight on immediate costs relative to future benefits \cite{Frederick2002, Attema2018}. In particular, \textbf{present bias}—commonly formalized in quasi-hyperbolic discounting—predicts that immediate costs are treated as disproportionately burdensome, leading to suboptimal delays even when the long-run value of protection is high \cite{Laibson1997, ODonoghueRabin1999}. This behavioral feature is especially relevant to vaccination decisions because many perceived costs are immediate, while benefits (risk reduction) are delayed and probabilistic. 

For policy impact evaluation, the discounting literature also emphasizes that the choice of discounting assumptions—and the possibility of departures from exponential discounting—has meaningful implications for how individuals value delayed health outcomes and how they respond to different policy interventions \cite{Attema2018, Baicker2011}. Together, these strands motivate the dissertation’s first measurement objective: to quantify whether delayed vaccine protection is discounted in a manner consistent with present bias and time inconsistency, and to identify systematic heterogeneity in such discounting.

\textit{Gap and link to this dissertation.} While discounting and present bias are widely studied in intertemporal choice, there is limited preference-based evidence that directly quantifies present bias specifically in the valuation of \emph{delayed vaccine protection} as an exposure-relevant health investment \cite{Frederick2002, Laibson1997, Attema2018}. This motivates Essay 1’s focus on estimating structural discounting parameters for vaccination timing.

\subsubsection{Vaccination behavior and psychological determinants}

A second literature emphasizes that vaccination uptake is shaped by psychological and contextual determinants that extend beyond simple cost–benefit calculations. Frameworks on vaccine hesitancy highlight that confidence in vaccines and institutions, perceived risks and benefits, and perceived barriers to access jointly affect vaccination behavior \cite{MacDonald2015_SAGE, Brewer2017}. Systematic evidence further suggests that hesitancy is heterogeneous across populations and settings, with determinants that include trust, information environments, and healthcare-system experiences \cite{Larson2014, Callaghan2021_SSM}.

Importantly, these determinants are plausibly intertemporal in nature: low confidence or high perceived barriers can amplify the subjective burden of waiting, while stronger preventive-health routines may reduce the immediate friction associated with acting today \cite{Milkman2021_PNAS}. In parallel, behavioral research on vaccination identifies a set of proximate mechanisms—risk perception, anticipated regret, social norms, and perceived constraints—that help explain observed uptake patterns \cite{Brewer2017, Betsch2018_PLOS}. These mechanisms can naturally interact with timing: when protection is delayed, the decision becomes not only “vaccinate or not,” but also “is it worth bearing immediate costs now for benefits that arrive later?” This interaction is central to the dissertation’s conceptual framing of delayed protection as a biologically meaningful interval during which individuals remain exposed.

\textit{Gap and link to this dissertation.} Although vaccination research has richly documented determinants of uptake and hesitancy, less is known about how these determinants map into \emph{quantified} willingness to tolerate delays, and how they translate into policy-relevant metrics for targeting interventions. This motivates the dissertation’s emphasis on heterogeneous delay valuation and its policy implications across subgroups.

\subsubsection{Stated-preference and DCE methods for valuing health-related trade-offs}A third foundation is methodological: quantifying the value of waiting requires an empirical design that can cleanly vary waiting time while holding other attributes constant. **Discrete choice experiments (DCEs)** provide a widely used stated-preference approach for eliciting preferences over multi-attribute alternatives and estimating the trade-offs individuals are willing to make among attributes \cite{Louviere2000, Train2009, deBekkerGrob2012}. In health applications, best-practice guidance emphasizes careful attribute construction, experimental design, and interpretation of trade-offs to ensure policy relevance and internal validity \cite{Johnson2013ISPOR, LancsarLouviere2008}.

Embedding waiting time directly as an attribute of vaccination alternatives allows the analyst to estimate not only the direction and magnitude of delay aversion, but also to translate it into interpretable metrics such as willingness-to-accept compensation for delayed protection and predicted uptake under counterfactual policy bundles.

\textit{Gap and link to this dissertation.} While DCE methods are well established for valuing health-related trade-offs, they have been less frequently used to treat \emph{waiting time to protection} as a first-order attribute and to jointly report (i) structural discounting parameters and (ii) monetary shadow prices in an integrated framework. This motivates the dissertation’s multi-metric approach: estimating $\kappa/\delta$ (Essay 1), translating delay aversion into MWTA (Essay 2), and integrating both into equity-sensitive policy rules (Essay 3).



This dissertation conceptualizes vaccination decisions not as a one-time static choice, but as an intertemporal health investment \cite{Grossman1972, Cropper1977}. From this perspective, individuals face a fundamental tension: the costs of vaccination (e.g., time, inconvenience, immediate side effects, anxiety) are incurred in the present, while the benefits (protection against uncertain future infection) are realized only with a delay. A robust literature in behavioral economics suggests that individuals frequently discount future health benefits heavily and exhibit present bias, placing disproportionate weight on immediate costs relative to delayed rewards \cite{Laibson1997, ODonoghueRabin1999, Frederick2002}. In the context of preventive health, this "behavioral hazard" can lead to suboptimal delays even when long-term intention to vaccinate exists \cite{Baicker2011, Chapman2005}.


Empirically, quantifying these intertemporal valuations is challenging because vaccination timing is often jointly determined by preferences, constraints, and policy environments. Discrete choice experiments (DCEs) offer a structural way to elicit and quantify these preferences by embedding delays and other vaccine attributes directly into choice tasks, enabling estimation of marginal valuations and trade-offs under controlled, policy-relevant counterfactuals \cite{Louviere2000,Train2009}. In health applications, best-practice guidance emphasizes careful experimental design, attribute framing, and interpretation of trade-offs for policy translation \cite{Johnson2013ISPOR}. Building on this toolkit, this dissertation treats “waiting time to protection” as a first-order attribute of vaccination alternatives and measures how individuals trade off waiting against monetary and non-monetary attributes.

Overall research objective. This dissertation quantifies the behavioral cost of delayed vaccine protection and its heterogeneity—by estimating both monetary trade-offs (MWTA) and structural intertemporal preference parameters (e.g., $\kappa$, $\delta$)—and translates these estimates into actionable, equity-sensitive policy guidance on when to prioritize speeding up access versus offering compensation and targeted prioritization \cite{Attema2018,Train2009,Louviere2000,Johnson2013ISPOR}.

To achieve this objective, the dissertation addresses three research questions:

RQ1 (Essay 1 — Time Discounting).
Do individuals exhibit present bias when valuing delayed vaccine protection, and what factors explain heterogeneity in the intertemporal preference parameters $\kappa$ and $\delta$?

RQ2 (Essay 2 — MWTA).
What is the monetary value individuals assign to delayed vaccine protection (MWTA), and how does this valuation vary across population subgroups?

RQ3 (Essay 3 — Integration and policy rules).
How can intertemporal preference parameters $(\kappa, \delta)$ and MWTA be jointly integrated to derive equity-sensitive policy rules that guide when to prioritize faster access versus monetary compensation?

This dissertation is organized as a three-essay study rooted in a discrete choice experiment (DCE) that treats waiting time to protection as a primary attribute of vaccination alternatives. This design enables a rigorous, policy-relevant valuation of delayed protection beyond simple binary uptake \cite{Louviere2000, Train2009, Johnson2013ISPOR}.

Essay 1 establishes the intertemporal-behavioral foundation. It quantifies how individuals discount delayed vaccine protection by estimating structural discounting parameters (e.g., $\kappa$ and $\delta$). This essay tests for the presence of present bias (time inconsistency) and examines how delay sensitivity varies with biological vulnerability, institutional trust, and preventive routines \cite{Frederick2002, Laibson1997, Attema2018}.Essay 2 translates this delay aversion into a policy-facing monetary metric. By estimating the Marginal Willingness to Accept (MWTA) compensation for delayed protection, it characterizes the "shadow price" of waiting across population subgroups and simulates how incentive–delay bundles can shift uptake under counterfactual access scenarios \cite{Louviere2000, Train2009}.Essay 3 integrates these monetary (MWTA) and structural ($\kappa/\delta$) measures into a unified "time-tax" policy framework. It derives equity-sensitive design rules for immunization programs, identifying conditions under which policymakers should prioritize accelerating access versus offering compensation or targeting high-burden groups \cite{Attema2018, Brewer2017}.


This dissertation makes five specific contributions to the literature on vaccination timing and intertemporal health decision-making:

1. Reframing Vaccination as Intertemporal Investment: It conceptualizes vaccination timing not merely as compliance, but as an intertemporal health investment where delayed protection extends the exposure window. This links classic health-capital models to preventive behavior, extending the theoretical frameworks of \citet{Grossman1972} and \citet{Cropper1977}.
2. Methodological Advancement in Health DCEs: It provides a novel measurement strategy by embedding waiting time directly into vaccine choice sets. This approach allows for the valuation of delay distinct from refusal, consistent with the specific nature of vaccine hesitancy \cite{MacDonald2015_SAGE} and best practices in stated-preference methods \cite{Louviere2000, Johnson2013ISPOR}.
3. Dual Quantification of Delay Valuation: It characterizes delay cost through two complementary metrics: (i) structural parameters ($\kappa$, $\delta$) that detect behavioral biases like time inconsistency \cite{Frederick2002, ODonoghueRabin1999}, and (ii) a monetary measure (MWTA) that translates delay aversion into interpretable economic units for program evaluation \cite{Train2009}.
4. Documenting Heterogeneity: It systematically documents how delay sensitivity varies across individuals, providing empirical evidence that time preferences are central determinants of preventive health decisions, alongside traditional psychological factors \cite{Chapman2005, Brewer2017}.
5. Connecting Behavior to Policy Design: Finally, by integrating structural and monetary measures, it derives actionable, equity-sensitive policy rules. This bridges the gap between behavioral valuation and the design of immunization campaigns, offering a framework to mitigate the "time tax" imposed on vulnerable populations \cite{Attema2018, Baicker2011}.




\chapter{Common Data, DCE Design, and Econometric Backbone}
\label{ch:common_methods}

\section{Study Context and Sampling}

Vaccination is widely recognized as a cornerstone of infectious-disease prevention, yet real-world vaccination behavior is not simply a binary choice between uptake and refusal. A central policy challenge is heterogeneity in timing: many individuals delay vaccination even when vaccines are available, thereby remaining exposed to infection risk for longer periods and weakening the effective coverage of immunization programs \citep{Halilova2022, Halilova2023}.

Empirically, this dissertation studies these timing trade-offs in a Chinese city---Wuhan---during the broader COVID-era immunization environment, where vaccination decisions were highly salient and program logistics could generate meaningful variation in waiting time. This setting provides a salient context for measuring how individuals value timeliness when delayed protection implies a longer exposure window. Because delayed protection is evaluated in a high-salience context, the estimated valuation of waiting provides an informative benchmark for other immunization settings, where perceived risk and salience may differ.

From a policy perspective, these delays matter because they shift protection into the future while infection risk persists in the present, slowing the epidemiological and social benefits of immunization. Consequently, vaccination timing can be understood as an intertemporal health investment decision in which immediate costs (time, inconvenience, perceived side effects) are weighed against delayed and uncertain health benefits \citep{Grossman1972, Meissner2025, Strickland2022}. This intertemporal framing is especially relevant in infectious-disease settings where ``waiting'' has biological meaning: delaying vaccination extends the period of vulnerability and can amplify perceived uncertainty and behavioral burdens.

At the same time, behavioral research emphasizes that vaccination decisions are shaped by psychological and contextual determinants---such as confidence, perceived risks, and perceived barriers---which can interact with intertemporal trade-offs when protection is delayed \citep{Halilova2024, Grytten2024}. These insights imply that policies aimed at improving vaccination timeliness may require more than information provision or supply expansion alone; they may also involve reducing waiting time, offering compensation or incentives, and targeting high-burden groups whose delay sensitivity is systematically greater.

Empirically, identifying the valuation of delayed protection is challenging using observational timing data alone because realized waiting time is jointly determined by preferences, constraints, and program rules. To isolate preferences for timeliness in a policy-relevant way, this dissertation uses a discrete choice experiment (DCE) that embeds waiting time to protection directly as an attribute of vaccination alternatives, allowing estimation of trade-offs between timeliness and other vaccine attributes under controlled counterfactual scenarios \citep{Louviere2000, Train2009, DCEreview2025}. This approach supports two complementary policy outputs: structural estimates of how individuals discount delayed protection and monetary translations of delay aversion that can be used to compare ``speeding up access'' versus ``compensation'' strategies.


\subsection{Participant recruitment and eligibility}
Participants were recruited from the general adult population in Wuhan, China, as part of a survey-based discrete choice experiment (DCE) on vaccination preferences and timing. Recruitment was conducted during the broader COVID-era immunization environment, using a online panel survey platform to reach adult respondents residing in Wuhan. Eligibility required that participants (i) were at least 18 years old, (ii) resided in Wuhan at the time of the survey, and (iii) provided informed consent prior to beginning the questionnaire. Respondents who failed key data-quality checks (e.g., incomplete survey submissions, implausibly short completion times, or comprehension checks where applicable) were excluded from the analytic sample.

The study instrument consisted of two components. First, respondents completed a set of background questions capturing socio-demographic characteristics, health status and indicators of biological vulnerability, as well as measures related to institutional trust and preventive-health routines. Second, respondents completed the discrete choice experiment (DCE). To reduce respondent burden while maintaining statistical efficiency, each respondent completed \textbf{six} choice tasks drawn from the final experimental design; each task presented \textbf{two} hypothetical vaccine alternatives and an \textbf{opt-out} option. After applying standard attention and data-quality checks, the final analytic sample consisted of \textbf{1{,}027} respondents who completed all required DCE tasks and met the eligibility and quality criteria described above.


\subsection{Sample profile}
\label{subsec:sample_profile}

The final analytic sample consists of 1{,}027 adult respondents residing in Wuhan, China. The sample is broadly balanced by gender (53.0\% female; 47.0\% male). Age composition spans the adult life course: 21.0\% are aged 18--34, 31.2\% are 35--44, 22.5\% are 45--54, 13.2\% are 55--65, and 12.0\% are 65 or older. Respondents include both urban and rural residents (47.0\% urban; 53.0\% rural), enabling analyses of place-based heterogeneity in delay sensitivity and valuation. Educational attainment is diverse: 4.0\% report no formal schooling, 22.0\% primary school, 25.7\% middle school, 20.4\% high school, 15.6\% college/university, and 11.0\% a graduate degree. Overall, these distributions provide substantial demographic variation for the subgroup analyses developed in Essays 1--3.

To contextualize the sample composition, Table~\ref{tab:table1_census_benchmark} benchmarks the analytic sample against the 2020 Wuhan Census for gender, age group, and urban--rural residence. The sample closely approximates the city-level population distributions along these dimensions, with minor deviations that follow typical patterns in online survey panels, including a modest upward shift in higher educational attainment. Pearson $\chi^2$ tests for gender, age group, and residence are statistically non-significant at the 5\% level, indicating no systematic deviation from census distributions. In addition, post-stratification weighted mixed-logit models yield parameter estimates substantively similar to those from the unweighted sample, supporting the robustness of the main findings to demographic composition (see Appendix~C).

\begin{table}[htbp]
\centering
\caption{Comparison of analytic sample with 2020 Wuhan Census benchmarks}
\label{tab:table1_census_benchmark}
\begin{tabular}{lcc}
\toprule
Characteristic & Sample (\%) & Wuhan Census (\%) \\
\midrule
\textbf{Gender} \\
\hspace{3mm}Female & 53.0 & 51.1 \\
\hspace{3mm}Male   & 47.0 & 48.9 \\
\addlinespace

\textbf{Age distribution} \\
\hspace{3mm}18--34 & 21.0 & 23.9 \\
\hspace{3mm}35--44 & 31.2 & 28.6 \\
\hspace{3mm}45--54 & 22.5 & 21.7 \\
\hspace{3mm}55--65 & 13.2 & 14.3 \\
\hspace{3mm}65+    & 12.0 & 11.5 \\
\addlinespace

\textbf{Residence} \\
\hspace{3mm}Urban & 47.0 & 48.3 \\
\hspace{3mm}Rural & 53.0 & 51.7 \\
\addlinespace

\textbf{Education level} \\
\hspace{3mm}No formal schooling & 4.0  & 5.6 \\
\hspace{3mm}Primary school      & 22.0 & 24.3 \\
\hspace{3mm}Middle school       & 25.7 & 27.1 \\
\hspace{3mm}High school         & 20.4 & 21.7 \\
\hspace{3mm}College/University  & 15.6 & 13.4 \\
\hspace{3mm}Graduate degree     & 11.0 & 7.9 \\
\bottomrule
\end{tabular}

\vspace{2mm}
\begin{minipage}{0.92\linewidth}
\footnotesize \textit{Note.} Census benchmarks are drawn from the Wuhan Bureau of Statistics \cite{WuhanStats2021}
 Percentages may not sum to 100 due to rounding. Post-stratification weighting yields mixed-logit and discount-rate estimates substantively similar to unweighted results (see Appendix~C).
\end{minipage}
\end{table}

\section{Econometric Backbone and Structural Metrics}
\label{sec:backbone_econometrics}

This section formalizes the econometric framework used throughout the dissertation to recover behavioral parameters from the DCE. The primary objective is to move beyond descriptive choice patterns toward structural metrics that quantify the intertemporal trade-offs individuals make during vaccination decisions.

\subsection{The Random Utility Framework (RUM)}
Consistent with established discrete choice theory \citep{Louviere2000, Train2009}, I model vaccination choice within a Random Utility Framework. The utility $U_{nit}$ that individual $n$ derives from vaccine alternative $i$ in choice task $t$ is decomposed into a systematic component $V_{nit}$ and a stochastic error term $\epsilon_{nit}$:
\begin{equation}
U_{nit} = V_{nit} + \epsilon_{nit} = \beta_{0} \text{OptOut} + \beta_{T} \cdot f(T_{nit}, \delta, \kappa) + \sum_{k} \beta_{k} X_{knit} + \epsilon_{nit},
\label{eq:rum_base}
\end{equation}
where:
\begin{itemize}
    \item \textbf{OptOut}: A binary indicator for the no-vaccination alternative.
    \item \textbf{T}: Waiting time to protection (0, 1, 3, or 6 months).
    \item \textbf{f(T)}: The intertemporal transformation function (Exponential or Hyperbolic).
    \item \textbf{X}: A vector of non-temporal attributes (Efficacy, Origin, Side Effects, Incentives).
    \item \textbf{$\epsilon$}: An i.i.d.\ extreme value (Type I) error term.
\end{itemize}

\subsection{Structural Metrics: BDM and MWTA}
A core contribution of this dissertation is the development and application of two structural metrics that translate raw coefficients into policy-relevant units: the Behavioral Delay Multiplier (BDM) and the Marginal Willingness to Accept (MWTA).

\paragraph{The Behavioral Delay Multiplier (BDM).} 
The BDM is introduced as a novel metric to quantify the "psychological time tax" imposed by delayed protection. While discount parameters ($\delta, \kappa$) are structurally informative, they often lack intuitive clarity for policymakers. The BDM maps the hyperbolic delay parameter into a subjective time-inflation factor:
\begin{equation}
\text{BDM} = \frac{1}{1-\hat{\kappa}}.
\label{eq:bdm_def_backbone}
\end{equation}
The BDM represents the ratio of \textit{perceived} to \textit{objective} waiting time. For example, a BDM of 1.25 implies that a 30-day administrative delay is experienced as a 37.5-day behavioral burden. This metric allows for a direct comparison of the experiential burden of delay across different population subgroups, such as those with varying levels of institutional trust.

\paragraph{Marginal Willingness to Accept (MWTA).} 
To perform rigorous policy impact evaluation, I translate delay aversion into monetary units using the MWTA, defined as the ratio of the marginal utility of time to the marginal utility of money:
\begin{equation}
\text{MWTA}_{\text{delay}} = - \frac{\partial V / \partial T}{\partial V / \partial \text{Incentive}} = - \frac{\beta_{T}}{\beta_{\text{money}}}.
\label{eq:mwta_def_backbone}
\end{equation}
This "shadow price of waiting" provides a common metric for evaluating the efficiency and equity of speed-ups versus monetary compensation strategies, which is the focus of Essay 2 and Essay 3.

\subsection{Capturing Heterogeneity: Mixed Logit (MXL)}
To account for unobserved preference heterogeneity and the panel nature of the DCE (six tasks per respondent), I utilize Mixed Logit models. The coefficients $\beta_n$ are assumed to follow a multivariate normal distribution $f(\beta|\theta)$, where $\theta$ represents the mean and covariance parameters. Models are estimated via Simulated Maximum Likelihood (SML) using 1,000 Halton draws to ensure robust identification of structural parameters and their distributions across the Wuhan population.

\section{DCE design overview}

\subsection{Choice task structure}
Each DCE choice task asked respondents to choose among \textbf{two} hypothetical vaccine alternatives (A vs.\ B) and an \textbf{opt-out} option. Each respondent completed \textbf{six} choice tasks, a task load selected to balance respondent burden with statistical efficiency. Choice tasks were administered in an online survey environment, and the ordering of tasks and the display of alternatives were randomized by the survey platform to minimize ordering and labeling effects. 

\subsection{Development of Attributes and Levels}
\label{subsec:attr_level_dev}

Attribute development aimed to capture the core trade-offs individuals face in real vaccination decisions and followed established good-practice principles for stated-preference research in health economics. Candidate attributes were identified through a review of vaccination-behavior studies, policy reports, and behavioral-economic frameworks emphasizing delay, risk, and incentives. Consultations with public-health practitioners in China and a 60-respondent pilot survey refined attribute wording, ensured cognitive clarity, and confirmed that the attributes reflected realistic decision environments.



The final attributes included vaccine delays (0, 1, 3, 6 months), vaccine efficacy (50\%, 70\%, 95\%), side-effect severity (mild, moderate, severe), monetary incentives (0, 50, 200, 800 RMB), and vaccine origin (domestic, imported). Vaccine-delay levels were chosen to reflect typical postponements generated by supply fluctuations, phased eligibility, and appointment bottlenecks in large-scale rollouts. Pretest findings indicated that most respondents viewed delays of 1--3 months as realistic and acceptable, confirming that monthly increments represent a salient and credible unit over which individuals evaluate delayed protection.

Monthly increments also align with the practical realities of vaccination planning, where delays of one to several months---not days---are the operationally relevant horizon over which programmes adjust supply distribution, staffing, and booking systems. Using monthly delay levels therefore enables estimation of the behaviourally meaningful compensation required to tolerate the types of postponements that actually occur in real programmes.

Methodologically, this implies that the DCE is not intended to recover very short-run discounting curvature or identify structural $\beta$--$\delta$ parameters, which would require day-level variation. Instead, the experiment yields reduced-form valuations of vaccine delay over a mid-range horizon, providing estimates that align directly with the operational decisions policymakers face when managing congestion and scheduling in vaccination programmes.

Financial-incentive levels were informed by magnitudes used in recent field experiments
and policy pilots in China and neighbouring East Asian jurisdictions, ensuring that respondents
evaluated realistic and policy-relevant amounts \cite{CamposMercade2021,Zhang2023,Thirumurthy2022}.


\begin{table}[htbp]
\centering
\caption{Experimental attributes, levels, and coding used in the discrete choice experiment}

\label{tab:attributes_levels}
\begin{tabular}{p{3.3cm}p{4.2cm}p{3.2cm}p{5.0cm}}
\toprule
\textbf{Attribute} & \textbf{Levels} & \textbf{Coding} & \textbf{Description} \\
\midrule
Vaccine delays & 0, 1, 3, 6 months (immediate/walk-in to 6 months) & 0, 1, 3, 6 &
Waiting time until effective protection becomes available. \\
\addlinespace

Vaccine origin & Domestic; Imported & 0 = Domestic, 1 = Imported &
Indicates whether the vaccine was produced locally or overseas. \\
\addlinespace

Vaccine efficacy & 50\%, 70\%, 95\% & 0.50, 0.70, 0.95 &
Probability of protection against infection. \\
\addlinespace

Side effects & Mild; Moderate; Severe & 0, 1, 2 &
Ordered severity of expected vaccine side effects; higher values indicate more severe reactions. \\
\addlinespace

Monetary incentives (RMB) & 0, 50, 200, 800 & 0, 50, 200, 800 &
Monetary compensation offered upon receiving the vaccine. \\
\bottomrule
\end{tabular}

\vspace{2mm}
\begin{minipage}{0.92\linewidth}
\footnotesize \textit{Notes.} Attribute codings reflect the numerical values entered into the econometric model. Side-effect severity is treated as an ordinal variable.
\end{minipage}
\end{table}




\clearpage

\subsection{Experimental design generation}
The experimental design was constructed using a \textit{D-efficient} design strategy to maximize statistical efficiency given the selected attributes and levels. Initial parameter priors were obtained from a pilot study (n=60) and used to generate the final D-efficient design. The design was evaluated to ensure adequate balance and overlap across attribute levels and to preserve plausibility of vaccine profiles. The final design was administered \textbf{without blocking}.

\subsection{Survey instrument and framing}
The survey instrument consisted of two components. First, respondents completed a background module collecting socio-demographic characteristics, health status and indicators of biological vulnerability, as well as measures related to institutional trust and preventive-health routines. Second, respondents completed the DCE module. Prior to the choice tasks, the survey provided standardized instructions defining ``waiting time to protection'' as the time until vaccine-induced protection becomes effective, and emphasized that longer waiting time implies remaining exposed to infection risk for a longer period. Attribute descriptions were presented in simple, concrete language to support comprehension, and the instrument was administered in Chinese.

To support data quality, the survey incorporated clear task instructions and standardized attribute definitions across respondents. The pilot study (n=60) was used to assess comprehension and cognitive burden and to refine wording and presentation; minor revisions were implemented to improve clarity and reduce ambiguity prior to fielding the main survey. Each task also included an opt-out option to reflect realistic choice settings in which respondents may decline vaccination under certain attribute combinations.

\subsection{Pre-test and instrument refinement}
\label{subsec:pretest_refinement}

Prior to fielding the main survey, we conducted a pilot study (n=60) to assess comprehension, cognitive burden, and the clarity of attribute descriptions and choice-task wording. The pilot served two purposes. First, it provided a structured cognitive refinement step to verify that respondents understood the key concepts---especially ``waiting time to protection''---and that the implied trade-offs were interpreted as intended. Second, pilot responses were used to inform parameter priors for the D-efficient experimental design used in the main study.

Figure~\ref{fig:attribute-development-flow} summarizes the attribute-development and refinement workflow. Candidate attributes were first identified from the evidence base (literature on vaccination behavior and behavioral economics, relevant policy reports, and expert consultation). The pilot study then supported cognitive checks and wording refinements, including removal of redundant or confusing elements where necessary. In addition, attribute levels were justified to span realistic and policy-relevant ranges (e.g., delay horizons reflecting operational roll-out constraints and incentive magnitudes consistent with policy pilots). The pilot study confirmed that respondents understood the tasks, that no single attribute dominated choices, and that minor wording revisions improved clarity and reduced cognitive burden prior to fielding the main survey.

\begin{figure}[htbp]
  \centering
  \includegraphics[width=0.85\textwidth]{Polit plot.png}
  \caption{Attribute development, pilot-based cognitive refinement (n=60), and finalization of attributes and levels for the DCE.}
  \label{fig:attribute-development-flow}
\end{figure}

\clearpage

\section{Data construction and variable definitions}

\subsection{Choice data structure}
The DCE data are organized in long format, with one row per respondent--choice set--alternative. The analytic dataset contains $1{,}027$ respondents, each completing 6 choice sets, and each set containing three alternatives (A, B, and an opt-out alternative C), yielding $1{,}027 \times 6 \times 3 = 18{,}486$ observations. The key identifiers are \texttt{RespondentID} (respondent), \texttt{Choiceset} (choice-task index), and \texttt{Alt} (alternative label). The dependent variable \texttt{Choice} is coded as 1 for the selected alternative within each choice set and 0 otherwise. The opt-out alternative is explicitly represented as \texttt{Alt = C}, and an opt-out alternative-specific constant is provided as \texttt{ASC\_optout} for model estimation. For compatibility with standard discrete-choice software, the dataset also includes \texttt{chid} (choice-set identifier) and \texttt{alt} (numeric alternative index).

\subsection{Attribute coding}
All DCE attributes are coded using the numerical values entered into the econometric model. Waiting time to protection is captured by \texttt{Vaccine Delay} (in months), taking values 0, 1, 3, and 6 for vaccine alternatives. Vaccine origin is coded as \texttt{VaccineOrigin}$=0$ for domestic and $=1$ for imported. Vaccine efficacy is coded as \texttt{VaccineEfficacy}\,$\in\{0.50, 0.70, 0.95\}$. Side-effect severity is coded as an ordered variable \texttt{SideEffects}$\in\{1,2,3\}$ corresponding to mild, moderate, and severe, respectively, and is treated as ordinal in the main specifications. Monetary incentives are coded in RMB as \texttt{CashIncentives}$\in\{0,50,200,800\}$. For the opt-out alternative (C), attribute values are set to zero by construction (e.g., \texttt{WaitTime}=0, \texttt{VaccineEfficacy}=0, \texttt{SideEffects}=0, \texttt{CashIncentives}=0), and opt-out behavior is captured via \texttt{ASC\_optout}. In addition to raw attribute values, standardized versions (\texttt{Vaccine Delay\_std}, \texttt{VaccineEfficacy\_std}, \texttt{SideEffects\_std}, \texttt{CashIncentives\_std}) are provided for robustness and comparability across specifications.

\subsection{Key covariates for heterogeneity}
\label{subsec:covariates_heterogeneity}

To examine systematic heterogeneity in preferences for vaccination delay and other vaccine attributes, the analysis incorporates respondent-level covariates measured in the background module of the survey. These covariates are used either as interaction terms with key DCE attributes (to test for systematic differences in marginal utilities) or to define pre-specified subgroups reported throughout Essays~1--3.

First, \textit{biological vulnerability} captures health-related characteristics that may increase the perceived cost of remaining unprotected while waiting for vaccination. This construct is proxied using self-reported health measures collected in the survey, including [chronic-condition indicator / self-rated health / other vulnerability measures used in Essays~1--3]. Second, \textit{institutional trust} captures confidence in public-health authorities and vaccine-related information sources. Trust is measured using [scale name or item battery], and the analysis uses either the continuous trust score or threshold-based groups (e.g., high vs.\ low trust) consistent with the specifications in the essays. Third, \textit{preventive-health routines} capture stable protective behaviors and habits (e.g., mask use, hygiene practices, or other routine prevention measures) that may shape perceived urgency and the salience of delayed protection. These routines are summarized using a composite score constructed from survey items.

In addition, the analysis includes standard socio-demographic controls such as age, gender, education, income, and urban--rural residence, as well as other behavioral or risk-related measures where available (e.g., perceived infection risk). Detailed item wording, coding, and any reliability statistics for constructed indices are documented in Appendix~[X], and the main text focuses on how these covariates explain heterogeneity in delay sensitivity, monetary valuation of delay, and model-implied uptake under counter.



\section{Econometric backbone: baseline preference model}

\subsection{Random utility framework and baseline mixed logit}
\label{subsec:ruf_mxl}

Respondents are assumed to maximize utility when choosing among alternatives in each DCE task. Let $n=1,\dots,N$ index respondents, $t=1,\dots,T$ index choice tasks, and $i \in \{A,B,C\}$ index alternatives (with $C$ denoting the opt-out). Utility from alternative $i$ in task $t$ for respondent $n$ is specified as
\begin{equation}
U_{nit} = V_{nit} + \varepsilon_{nit}.
\end{equation}
The systematic component $V_{nit}$ is linear in attributes:
\begin{equation}
V_{nit} = \alpha \cdot \text{ASC}^{\text{opt}}_{nit}
+ \beta_{n,d}\,\text{Delay}_{nit}
+ \beta_{n,e}\,\text{Efficacy}_{nit}
+ \beta_{n,s}\,\text{SideEffects}_{nit}
+ \beta_{n,o}\,\text{Origin}_{nit}
+ \beta_{n,c}\,\text{Incentive}_{nit},
\label{eq:mxl_utility}
\end{equation}
where $\text{ASC}^{\text{opt}}_{nit}=1$ for the opt-out alternative and 0 otherwise. $\text{Delay}_{nit}$ denotes vaccination delay (waiting time to protection, in months), $\text{Efficacy}_{nit}$ is vaccine efficacy (coded as probabilities), $\text{SideEffects}_{nit}$ is side-effect severity (coded ordinally), $\text{Origin}_{nit}$ indicates imported versus domestic origin, and $\text{Incentive}_{nit}$ is the cash incentive in RMB.

The baseline model is a \textit{mixed logit} (random-parameters logit), which allows preference parameters to vary across individuals:
\begin{equation}
\boldsymbol{\beta}_n = \boldsymbol{\beta} + \boldsymbol{\eta}_n,
\end{equation}
where $\boldsymbol{\beta}$ denotes the population mean coefficients and $\boldsymbol{\eta}_n$ captures individual-specific deviations drawn from a parametric distribution (e.g., normal) to represent unobserved preference heterogeneity. Conditional on $\boldsymbol{\beta}_n$ and assuming $\varepsilon_{nit}$ is i.i.d.\ Type-I extreme value, the choice probability takes the standard logit form
\begin{equation}
P_{nit}(\boldsymbol{\beta}_n)=
\frac{\exp(V_{nit})}{\sum_{j \in \{A,B,C\}} \exp(V_{njt})}.
\end{equation}
The unconditional probability integrates over the distribution of random coefficients:
\begin{equation}
P_{nit}=\int P_{nit}(\boldsymbol{\beta}_n)\, f(\boldsymbol{\beta}_n \mid \theta)\, d\boldsymbol{\beta}_n,
\label{eq:mxl_prob}
\end{equation}
where $f(\cdot)$ is parameterized by $\theta$ (means and dispersion parameters). The integral in \eqref{eq:mxl_prob} is evaluated by simulation, and parameters are estimated by maximum simulated likelihood using the panel structure of the data (choices repeated within respondent).


\subsection{Benchmark models: multinomial logit and conditional logit}
\label{subsec:benchmark_mnl_cl}

As benchmark specifications, we estimate a multinomial logit (MNL) model and a conditional logit (CL) model. These benchmarks provide transparent reference points for the baseline mixed logit results and facilitate comparison of implied trade-offs and predicted uptake.

\subsubsection{Multinomial logit (MNL)}
\label{subsec:mnl}
Under the assumption that the idiosyncratic errors $\varepsilon_{nit}$ are independently and identically distributed Type-I extreme value, the probability that respondent $n$ chooses alternative $i$ in task $t$ is
\begin{equation}
P_{nit} =
\frac{\exp(V_{nit})}{\sum_{j \in \{A,B,C\}} \exp(V_{njt})},
\end{equation}
where $V_{nit}$ is specified as in Equation~\eqref{eq:mxl_utility} but with homogeneous (fixed) coefficients across individuals. In this benchmark, the opt-out alternative is accommodated using an alternative-specific constant $\text{ASC}^{\text{opt}}_{nit}$ that equals 1 for the opt-out alternative and 0 otherwise.

\subsubsection{Conditional logit (CL)}
\label{subsec:cl}
The conditional logit model specifies utility as a function of alternative-specific attributes and assumes the same i.i.d.\ Type-I extreme value distribution for $\varepsilon_{nit}$. When utility depends only on alternative-specific attributes (as in this DCE) and coefficients are fixed, the CL and MNL formulations yield the same logit choice probabilities. We report CL estimates as an additional benchmark because it is the canonical specification for DCE data and provides a direct comparison to the mixed logit framework when preference heterogeneity is not modeled.

\subsection{Derived metrics and policy-relevant quantities}
\label{subsec:derived_metrics}

To translate estimated preference parameters into policy-relevant measures, we report three classes of derived quantities: (i) monetary valuations of vaccination delay, (ii) implied discounting measures for delayed protection, and (iii) model-based counterfactual predictions of uptake under alternative policy bundles.

\subsubsection{Monetary valuation of delay (MWTA)}
A central object is the monetary compensation required to offset an additional month of vaccination delay. Let $\beta_{d}$ denote the marginal utility of vaccination delay (waiting time to protection) and $\beta_{c}$ denote the marginal utility of cash incentives. The marginal willingness-to-accept compensation (MWTA) for a one-month increase in delay is defined as
\begin{equation}
\text{MWTA} = -\frac{\beta_{d}}{\beta_{c}},
\label{eq:mwta}
\end{equation}
expressed in RMB per month. In mixed logit models, MWTA can be computed using mean coefficients and, where appropriate, using individual-specific posterior means of random coefficients to characterize heterogeneity in delay valuation. Uncertainty intervals are obtained using [delta method / bootstrap], consistent with the estimation approach reported in each essay.

\subsubsection{Discounting measures for delayed protection}
The estimated delay sensitivity can also be summarized in terms of implied discounting over delayed protection. In particular, Essays~[1/2] map the estimated delay coefficient(s) to reduced-form discounting metrics that summarize how the value of protection declines with waiting time over the 0--6 month horizon. These measures are reported as [e.g., implied discount rate / behavioral delay multiplier / relative value of protection at $t$ months compared to immediate protection], and are interpreted as mid-range valuations of delay rather than structural short-run discounting curvature.

\subsubsection{Counterfactual policy simulations}
Finally, the estimated models are used to simulate predicted uptake under counterfactual policy scenarios that vary (i) waiting time to protection, (ii) incentive amounts, and (iii) vaccine quality attributes. For a given policy bundle, predicted uptake is computed as the model-implied probability of choosing either vaccination alternative (A or B) rather than opting out, averaged over respondents and choice tasks:
\begin{equation}
\widehat{\Pr}(\text{Uptake}) = \frac{1}{N}\sum_{n=1}^{N}\frac{1}{T}\sum_{t=1}^{T} \left( \hat{P}_{nAt} + \hat{P}_{nBt} \right),
\label{eq:uptake}
\end{equation}
where $\hat{P}_{nit}$ denotes the estimated choice probability under the counterfactual attribute levels. These simulations allow direct comparison of ``speed-up'' policies (reducing delay) versus ``compensation'' policies (increasing incentives), as well as combined strategies, and they form the basis for the policy trade-off menus developed in Essays~1--3.


\section{Validity checks and robustness summary}
\label{sec:validity_robustness}

\subsection{Internal validity and data quality}
\label{subsec:internal_validity}

A set of internal-validity and data-quality diagnostics was implemented to assess whether respondents understood the choice tasks and whether estimated preference parameters reflect stable trade-offs rather than random responding. These checks focus on (i) dominance behavior, (ii) attribute non-attendance (ANA), and (iii) within-respondent choice consistency across repeated tasks.Additional diagnostics on transitivity/choice consistency and correlations among individual-level covariates indicate no systematic preference reversals and negligible multicollinearity (Appendix~A).


\paragraph{Dominance and coherence.}
Dominance tests were conducted to evaluate whether respondents made coherent choices when presented with alternatives that were strictly inferior across attributes. Results indicate a very low dominated-choice rate (1.8\%), consistent with high internal coherence of the DCE responses (Table~\ref{tab:dominance_coherence}). In complementary summary diagnostics, more than 85\% of respondents selected the alternative offering superior attributes in dominance-style comparisons.

\begin{table}[htbp]
\centering
\caption{\textbf{Dominance and coherence diagnostics}}
\label{tab:dominance_coherence}
\begin{tabular}{p{8.8cm}p{3.2cm}}
\toprule
\textbf{Diagnostic} & \textbf{Result} \\
\midrule
Dominated-choice rate (share of dominance-test choices that violate dominance) & 1.8\% \\
Share choosing the alternative with superior attributes in dominance-style comparisons & $>$85\% \\
\bottomrule
\end{tabular}

\vspace{2mm}
\begin{minipage}{0.92\linewidth}
\footnotesize \textit{Notes.} Dominance tests evaluate whether respondents avoid selecting an alternative that is strictly inferior across attributes. A low dominated-choice rate and a high share selecting the superior alternative support internal coherence of stated preferences.
\end{minipage}
\end{table}



\paragraph{Attribute non-attendance (ANA).}
The extent to which respondents might ignore specific attributes was assessed using stated ANA reports (where available) and sensitivity diagnostics based on conditional-logit choice patterns. The diagnostics show no evidence of systematic ignoring for most attributes: choices invariant to delay levels, efficacy levels, side-effect levels, and cash-incentive levels were each 0.0\%. The only non-zero ANA signal is for vaccine origin (5.4\%), suggesting that a small subset of respondents treated origin as less salient than other attributes (Table~\ref{tab:ana_diagnostics}). Re-estimating models after excluding respondents who potentially ignored certain attributes yields coefficients nearly identical to the main specification, indicating minimal sensitivity of the results to ANA.

\begin{table}[htbp]
\centering
\caption{\textbf{Attribute non-attendance (ANA) diagnostics}}
\label{tab:ana_diagnostics}
\begin{tabular}{p{6.2cm}p{2.4cm}p{6.6cm}}
\toprule
\textbf{Attribute} & \textbf{ANA signal} & \textbf{Operational definition} \\
\midrule
Vaccination delay (waiting time to protection) & 0.0\% &
Choices invariant to delay levels across tasks (no sensitivity to delay variation). \\
\addlinespace

Vaccine efficacy & 0.0\% &
Choices invariant to efficacy levels across tasks (no sensitivity to efficacy variation). \\
\addlinespace

Side-effect severity & 0.0\% &
Choices invariant to side-effect levels across tasks (no sensitivity to side-effect variation). \\
\addlinespace

Monetary incentives & 0.0\% &
Choices invariant to incentive levels across tasks (no sensitivity to cash variation). \\
\addlinespace

Vaccine origin (domestic vs.\ imported) & 5.4\% &
Choices invariant to origin levels across tasks (origin treated as non-salient). \\
\bottomrule
\end{tabular}

\vspace{2mm}
\begin{minipage}{0.92\linewidth}
\footnotesize \textit{Notes.} ANA diagnostics report the share of respondents whose observed choices show no variation with respect to a given attribute (i.e., choices appear invariant to that attribute’s levels). Percentages are reported for the analytic sample.
\end{minipage}
\end{table}


\clearpage

\paragraph{Within-respondent stability.}
Within-respondent stability was assessed by examining the consistency of choices across repeated tasks. Intra-individual agreement across the six tasks is approximately 70\%, indicating meaningful stability in preferences rather than purely random responding. This level of within-respondent consistency supports the interpretation that estimated trade-offs reflect systematic valuation of vaccination delay, vaccine attributes, and incentives over repeated choice occasions.

\paragraph{Additional robustness checks.}
Additional robustness analyses were conducted to assess whether the main findings are sensitive to resampling variability, influential observations, or alternative sample definitions. Clustered bootstrap procedures at the respondent level (1{,}000 replications) yield coefficient magnitudes and uncertainty intervals that are consistent with the baseline estimates. Sensitivity analyses that trim the upper and lower tails of the MWTA distribution (e.g., excluding the top and bottom 5\%) produce nearly identical coefficient patterns and substantive conclusions, indicating that results are not driven by outliers. Further diagnostics on transitivity/choice consistency and correlations among individual-level covariates are reported in the appendix and indicate no evidence of systematic preference reversals and negligible multicollinearity.

% --- Placeholders for the appendix items you already have ---
% Figure: example choice task screenshot (your "Figure 4" in the paper appendix)
% Table: internal validity diagnostics (your "Table 7")
% Table: bootstrap CIs (your "Table 9")
% Table: trimmed-sample sensitivity (your "Table 10")
\chapter{Essay 1: Time Discounting in Vaccination Timing Preferences}
\label{ch:essay1_discounting}

\section{Introduction}

Vaccination is widely recognized as a cornerstone of infectious-disease prevention, yet vaccination behavior in practice is not merely a binary choice between uptake and refusal. A central policy challenge is heterogeneity in \emph{timing}: even when vaccines are available, many individuals postpone uptake, remaining exposed to infection risk for longer periods and weakening the effective coverage of immunization programs \cite{MacDonald2015,Larson2014}. From a program-design perspective, these delays matter because they shift protection into the future while infection risk persists in the present, slowing the epidemiological and social benefits of immunization. Vaccination timing can therefore be understood as an intertemporal health-investment decision in which immediate costs—time, inconvenience, anxiety, and perceived side effects—are traded off against delayed and uncertain health benefits \cite{Grossman1972,Frederick2002,Laibson1997}.

Despite extensive research on vaccine acceptance and hesitancy, evidence remains limited on how individuals value \emph{delayed protection} in a way that is behaviorally interpretable and directly usable for operational policy design. Logistical frictions such as supply fluctuations, phased eligibility, appointment bottlenecks, and administrative constraints routinely generate nontrivial postponements during large-scale rollouts. Yet much empirical work treats these delays as administrative noise rather than as a meaningful interval of continued vulnerability. Measuring how sharply the value of protection declines with waiting time is thus essential for designing policies that accelerate access, allocate limited capacity, or compensate individuals for delays.

This essay quantifies \emph{time discounting in vaccination timing preferences} using a discrete choice experiment (DCE) that varies waiting time to protection over a policy-relevant mid-range horizon (0--6 months), alongside vaccine efficacy, side-effect severity, vaccine origin, and monetary incentives. DCEs are well suited for recovering marginal valuations of policy-relevant attributes under controlled, experimentally varied choice environments \cite{Louviere2000,Train2009,Johnson2013ISPOR}. The empirical setting is Wuhan, China, during the broader COVID-era immunization environment, where vaccination decisions were salient and program logistics plausibly generated meaningful variation in time-to-protection. The design is intentionally aligned with operational planning horizons—months rather than days—so that the resulting estimates capture the valuations relevant to rollout management rather than short-lived day-to-day fluctuations.

Methodologically, the objective is not to identify fine-grained short-run curvature in discounting or to recover structural $\beta$--$\delta$ parameters that require day-level experimental variation \cite{Frederick2002,Laibson1997}. Instead, the analysis derives reduced-form discounting measures over the mid-range horizon that vaccination programs commonly confront. The estimated sensitivity to delay is used to characterize both average discounting and heterogeneity across individuals and subgroups, with particular attention to determinants that plausibly shape the urgency of protection: biological vulnerability, institutional trust, and preventive-health routines.

The contributions of this essay are threefold. First, it provides an empirical characterization of how delayed vaccine protection is discounted within a policy-relevant time window, translating operational delays into interpretable preference parameters. Second, it documents systematic heterogeneity in delay discounting associated with vulnerability and with institutional and behavioral factors, informing equity-relevant policy considerations. Third, it establishes a foundation for subsequent essays that translate delay valuations into monetary metrics and policy trade-off menus, enabling explicit comparisons between speeding up access and offering compensation.

The remainder of this essay is organized as follows. Section~\ref{sec:essay1_rq} states the research question and the estimands of interest. Section~\ref{sec:essay1_data} briefly summarizes the data and study design, with full details provided in Chapter~2. Section~\ref{sec:essay1_model} presents the econometric specification, identification strategy, and the mapping from estimated delay sensitivity to the exponential and hyperbolic discount parameters. Section~\ref{sec:essay1_results} reports the baseline discounting estimates, the Behavioral Delay Multiplier (BDM), and heterogeneity patterns across subgroups. Robustness and validity checks are presented in Section~\ref{sec:essay1_robust}. Section~\ref{sec:essay1_discussion} discusses implications for vaccination program design and concludes by linking to Essays~2 and~3.


\section{Research question and estimands}
\label{sec:essay1_rq}

This essay examines whether individuals discount delayed vaccine protection and how the strength of discounting varies across individuals and subgroups. The central research question is:

\begin{quote}
\textit{Do individuals exhibit time discounting when valuing delayed vaccine protection, and how does discounting vary with biological vulnerability, institutional trust, and preventive-health routines?}
\end{quote}

The analysis focuses on two primary estimands. First, the marginal utility of vaccination delay (waiting time to protection), denoted $\beta_{d}$ in the mixed logit specification, captures the average disutility associated with an additional month of delay relative to immediate protection. A negative and statistically significant $\beta_{d}$ indicates that, holding other vaccine attributes constant, individuals place lower value on protection that begins later—consistent with time discounting.

Second, the essay derives reduced-form discounting measures over the experimentally varied 0--6 month horizon. Using the estimated delay coefficient, the analysis characterizes how the utility value of protection declines with waiting time, conditional on vaccine efficacy, side effects, origin, and monetary incentives. These derived measures summarize the behavioral rate at which protection is discounted within the operationally relevant time window faced by vaccination programs.

Heterogeneity in discounting is assessed in two complementary ways. Unobserved heterogeneity is captured via random coefficients on the delay attribute in the mixed logit framework, allowing the marginal disutility of delay to vary across individuals. Observed heterogeneity is examined through pre-specified subgroup comparisons and interaction specifications linking delay sensitivity to biological vulnerability, institutional trust, and preventive-health routines.

Together, these estimands provide a structured empirical characterization of time discounting in vaccination timing decisions, forming the basis for subsequent essays that translate delay sensitivity into monetary valuations and policy trade-off comparisons.

\section{Theoretical Framework}
\label{sec:essay1_theory}

Vaccination timing can be conceptualized as an investment in health capital, where individuals evaluate the immediate utility costs of uptake---such as effort, inconvenience, or anxiety---against the discounted preservation of future health stocks \cite{Frederick2002}. In this framework, vaccination is an intertemporal choice in which discount parameters $\kappa$ and $\delta$ are not merely abstract preferences but can be shaped by an individual’s biological state and institutional environment. When individuals heavily discount future health returns, even modest delays impose substantial behavioral disutility, leading to underinvestment in preventive health.

Discrete choice experiments (DCEs) provide a structural method to quantify these intertemporal valuations by embedding temporal delay directly into the latent utility function for each vaccine alternative. Under an exponential specification, the utility of delayed protection is discounted at a constant rate $\delta$, reflecting time-consistent preferences. Conversely, a quasi-hyperbolic model allows discounting to be disproportionately steep in the short run, capturing present-biased preferences summarized by $\kappa$ \cite{Laibson1997}. These parameters translate the marginal disutility of waiting into behaviorally interpretable measures of how individuals trade off current utility costs against delayed biological protection.

Central to this framework is the Behavioral Delay Multiplier (BDM), which links formal discounting to subjective time perception. The BDM is defined as the ratio between perceived and objective waiting duration, representing a behavioral measure of the ``psychological time tax'' associated with delay. Within this health-economic context, the delay coefficient recovered from the DCE represents sensitivity to temporal friction, and the implied BDM provides an interpretable transformation of that sensitivity. This sensitivity is hypothesized to vary systematically across biological and health-related characteristics---such as age groups and chronic medical conditions---and levels of institutional trust, which can be interpreted as intertemporal cognitive capital by reducing subjective uncertainty about future health returns.


\section{Data and study design}
\label{sec:essay1_data}

This essay uses the common data and experimental framework described in Chapter~\ref{ch:common_methods}. The empirical setting is Wuhan, China, during the broader COVID-era immunization environment, where vaccination decisions were highly salient and logistical constraints generated meaningful variation in waiting times. The final analytic sample consists of 1{,}027 adult respondents who completed all required discrete choice tasks and satisfied the pre-specified eligibility and data-quality criteria detailed in Chapter~\ref{ch:common_methods}.

The survey instrument comprised two components. The first collected background information on socio-demographic characteristics, health status and indicators of biological vulnerability, institutional trust, and preventive-health routines. The second component consisted of a discrete choice experiment (DCE) designed to elicit preferences over vaccination timing and related attributes. Each respondent completed six choice tasks, with each task presenting two hypothetical vaccine alternatives and an opt-out option.

Vaccination alternatives varied along five experimentally manipulated attributes: (i) vaccination delay (waiting time to protection: 0, 1, 3, or 6 months), (ii) vaccine efficacy (50\%, 70\%, 95\%), (iii) side-effect severity (mild, moderate, severe), (iv) vaccine origin (domestic or imported), and (v) monetary incentives (50, 200, 800 RMB). The design was generated using a D-efficient algorithm informed by a 60-respondent pilot study. Task ordering and alternative presentation were randomized within the online survey platform to reduce ordering and labeling effects.

The DCE yields panel data with six repeated choice observations per respondent. Identification of delay sensitivity in this essay relies on randomized within-task variation in vaccination delay, holding other attributes constant across alternatives. Because the experimental design, coding decisions, and internal-validity diagnostics are common across the dissertation, those elements are documented in detail in Chapter~\ref{ch:common_methods}. The present section summarizes only the features necessary for interpreting the time-discounting estimates reported below.

\section{Econometric specification and identification}
\label{sec:essay1_model}
This section formalizes the econometric model used to recover intertemporal discount parameters and clarifies the sources of identification. Essay~1 focuses on identifying delay sensitivity and mapping it into structurally interpretable discount parameters $(\hat{\kappa},\hat{\delta})$. The full DCE instrument, attribute definitions, experimental design construction (pilot priors and D-efficient design), and general internal-validity diagnostics (e.g., dominance and attribute non-attendance) are documented in Chapter~2. Full details on estimation implementation details—including simulation settings for the mixed logit, coefficient standardization, and the full derivations mapping $\beta_{\text{delay}}$ to $(\hat{\kappa},\hat{\delta})$ —are reported in Chapter Appendix.

\subsection*{Utility specification}

Vaccination choices are modeled within a random utility framework. Individual $i$ chooses among alternatives $j \in \{A,B,\text{opt-out}\}$ in choice task $t$. Indirect utility is specified as:
\begin{align}
U_{ijt} \;=\;&
\beta_{\text{delay}}\, \text{Delay}_{ijt}
+ \beta_{\text{eff}}\, \text{Efficacy}_{ijt}
+ \beta_{\text{side}}\, \text{SideEffects}_{ijt}
+ \beta_{\text{origin}}\, \text{Origin}_{ijt}
+ \beta_{\text{cash}}\, \text{Cash}_{ijt}
\nonumber\\
&+ \text{ASC}_{\text{opt-out}} \cdot \mathbbm{1}\{j=\text{opt-out}\}
+ \varepsilon_{ijt},
\label{eq:utility_main}
\end{align}
where $\varepsilon_{ijt}$ is assumed i.i.d. Type~I extreme value. The indicator $\mathbbm{1}\{\cdot\}$ equals one for the opt-out alternative and zero otherwise.%
\footnote{If you do not use $\mathbbm{1}\{\cdot\}$ elsewhere, replace it with $\mathbf{1}\{\cdot\}$ and remove the need for the \texttt{bbm} package.}

The coefficient of primary interest is $\beta_{\text{delay}}$, which measures the marginal utility loss associated with a one-month increase in vaccination delay. In this essay, delay is interpreted as a biologically meaningful interval of continued exposure to infection risk, rather than a purely administrative inconvenience.

To account for unobserved preference heterogeneity and within-respondent correlation across repeated tasks, we estimate a mixed logit (random-parameters logit) model. Let $\boldsymbol{\beta}_i$ denote individual-specific taste parameters. We allow the delay coefficient (and, where appropriate, additional attributes) to vary across individuals:
\begin{equation}
\beta_{\text{delay},i} = \bar{\beta}_{\text{delay}} + \eta^{(\text{delay})}_i,
\qquad
\eta^{(\text{delay})}_i \sim \mathcal{N}(0,\sigma^2_{\text{delay}}),
\label{eq:random_delay}
\end{equation}
with analogous expressions for other random coefficients when included. Under mixed logit, the probability that respondent $i$ chooses alternative $j$ in task $t$ integrates the standard logit kernel over the distribution of $\boldsymbol{\beta}_i$:
\begin{equation}
P_{ijt}=\int
\frac{\exp\!\left(V_{ijt}(\boldsymbol{\beta}_i)\right)}
{\sum_{k}\exp\!\left(V_{ikt}(\boldsymbol{\beta}_i)\right)}
f(\boldsymbol{\beta}_i)\, d\boldsymbol{\beta}_i,
\label{eq:mxl_prob}
\end{equation}
where $V_{ijt}$ is the systematic component of utility in \eqref{eq:utility_main}. Estimation proceeds by simulated maximum likelihood; implementation details (draws, convergence checks, and alternative specifications) are provided in Appendix~C.

\subsection*{Identification strategy}

Identification of $\beta_{\text{delay}}$ relies on exogenous variation in waiting time generated by the experimental design. The D-efficient design randomizes attribute levels across vaccine profiles and balances delay independently of other attributes, supporting separate identification of time sensitivity from efficacy, side effects, origin, and incentives.

Two features of the design are central for identification:

\begin{enumerate}
\item \textbf{Within-task trade-offs.} In each choice task, respondents compare vaccine profiles that differ in delay (0, 1, 3, or 6 months) and other attributes. This generates direct marginal-rate-of-substitution information between waiting time and compensating attributes such as cash incentives and efficacy.

\item \textbf{Panel structure.} Each respondent completes multiple choice tasks, allowing the model to recover stable preference parameters while accounting for intra-individual correlation through random coefficients. This repeated-task structure increases efficiency and supports estimation of the distribution of delay sensitivity (i.e., $\sigma_{\text{delay}}$ in \eqref{eq:random_delay}).
\end{enumerate}

Under the maintained assumptions of random utility maximization and correct specification of $V_{ijt}$, the randomized variation in delay identifies the causal marginal disutility of waiting time. Because delay is experimentally assigned rather than chosen by respondents, the estimated $\beta_{\text{delay}}$ is not confounded by unobserved access constraints or endogenous appointment timing within the survey environment.

\subsection*{From reduced-form delay sensitivity to discount parameters}

The coefficient $\beta_{\text{delay}}$ is a reduced-form estimate of the marginal utility cost of waiting one additional month for protection. To interpret delay valuation through the lens of intertemporal choice, we transform this delay sensitivity into standard discounting parameters.

Let $t$ denote delay in months. We construct an empirical discount factor implied by the estimated marginal utility loss:
\begin{equation}
D_{\text{obs}}(t)=\exp\!\left(-\beta_{\text{delay}}\, t\right).
\label{eq:Dobs}
\end{equation}
These implied discount points are then mapped to two canonical discount functions.

\paragraph{Exponential discounting.}
\begin{equation}
D_{\text{exp}}(t)=\exp\!\left(-\hat{\delta}\, t\right),
\label{eq:Dexp}
\end{equation}
where $\hat{\delta}\ge 0$ is the exponential discount parameter reflecting time-consistent devaluation of delayed protection.

\paragraph{Hyperbolic discounting.}
\begin{equation}
D_{\text{hyp}}(t)=\frac{1}{1+\hat{\kappa}\, t},
\label{eq:Dhyp}
\end{equation}
where $\hat{\kappa}\ge 0$ captures steep short-run impatience consistent with present-biased valuations.

We recover $(\hat{\kappa},\hat{\delta})$ by fitting \eqref{eq:Dexp} and \eqref{eq:Dhyp} to the empirical discount points \eqref{eq:Dobs} across the observed delay horizons. The full estimation procedure---including construction of $D_{\text{obs}}(t)$, the objective functions used to fit $(\hat{\kappa},\hat{\delta})$, and alternative fitting/robustness checks---is provided in the \textit{Appendix to Essay~1}. Conceptually, this approach treats the experimentally identified delay coefficient as an empirical summary of intertemporal trade-offs and expresses it in parameterized discounting form for interpretation and comparison.


\subsection*{Heterogeneity in discounting}

Heterogeneity in time preferences is examined in two complementary ways.

\paragraph{Unobserved heterogeneity.}
The mixed logit model estimates a distribution for $\beta_{\text{delay},i}$ via $\sigma_{\text{delay}}$ in \eqref{eq:random_delay}, capturing continuous dispersion in delay sensitivity across respondents.

\paragraph{Observed heterogeneity.}
To test whether impatience varies systematically with biological vulnerability, institutional trust, and preventive routines, we augment \eqref{eq:utility_main} with interactions between delay and respondent characteristics:
\begin{equation}
U_{ijt} =
\beta_{\text{delay}} \, \text{Delay}_{ijt}
+ \boldsymbol{\gamma}'\!\left(\mathbf{Z}_i \times \text{Delay}_{ijt}\right)
+ \mathbf{X}_{ijt}'\boldsymbol{\beta}
+ \text{ASC}_{\text{opt-out}} \cdot \mathbbm{1}\{j=\text{opt-out}\}
+ \varepsilon_{ijt},
\label{eq:heterogeneity}
\end{equation}
where $\mathbf{Z}_i$ includes indicators for chronic conditions (biological vulnerability), institutional trust, and measures of preventive-health routines, and $\mathbf{X}_{ijt}$ collects the remaining vaccine attributes. The interaction vector $\boldsymbol{\gamma}$ measures how the marginal disutility of delay shifts across subgroups.

Subgroup-specific discount parameters are obtained by applying the mapping in \eqref{eq:Dobs}--\eqref{eq:Dhyp} to the estimated subgroup-specific delay sensitivities (e.g., $\beta_{\text{delay}}+\gamma_z$ for subgroup $z$). This approach yields comparable $(\hat{\kappa},\hat{\delta})$ estimates across groups while preserving a unified identification strategy.


\section{Results}
\label{sec:essay1_results}


\subsection{Model comparison and baseline specification}
\label{subsec:model_comparison}

We first compare alternative discrete-choice specifications to motivate the baseline model used to recover discount parameters. Table~\ref{tab:logit_comparison} reports estimates from a conditional logit, a multinomial logit (MNL) with alternative-specific constants (ASCs), and a mixed logit model. Two findings motivate our baseline choice. First, the mixed logit substantially improves overall fit relative to the conditional logit and MNL, as reflected in a higher log-likelihood and lower AIC. Second, the mixed logit estimates reveal statistically significant dispersion in delay sensitivity through a large and significant standard deviation of the random delay coefficient. This evidence indicates meaningful unobserved heterogeneity in how respondents value waiting time, which cannot be captured by fixed-coefficient models. Accordingly, we adopt the mixed logit specification as the baseline for subsequent analyses and for mapping delay sensitivity into the structural discount parameters $(\hat{\kappa},\hat{\delta})$.

\begin{table}[!htbp]
\centering
\caption{Comparison of logit models (dependent variable: choice)}
\label{tab:logit_comparison}
\begin{threeparttable}
\setlength{\tabcolsep}{7pt}
\begin{tabular}{lccc}
\toprule
 & Conditional Logit & MNL (ASCs) & Mixed Logit \\
\midrule
Vaccine Delays (std)      & $-0.131^{***}$ (0.019) & $-0.125^{***}$ (0.020) & $-0.353^{***}$ (0.027) \\
Vaccine Efficacy (std)    & $ 0.153^{***}$ (0.036) & $ 0.163^{***}$ (0.037) & $ 0.218^{***}$ (0.041) \\
Side Effects (std)        & $-0.042$      (0.027)  & $-0.082^{***}$ (0.030) & $ 0.008$      (0.031) \\
Cash Incentives (std)     & $ 0.327^{***}$ (0.018) & $ 0.313^{***}$ (0.019) & $ 0.431^{***}$ (0.025) \\
Vaccine Origin (Imported) & $ 0.244^{***}$ (0.038) & $ 0.239^{***}$ (0.038) & $ 0.093^{**}$  (0.047) \\
Opt-out ASC               & ---                   & $-0.214^{**}$ (0.093)  & $-0.287^{**}$ (0.118) \\
ASC: B                    & ---                   & $-0.108^{***}$ (0.033) & --- \\
ASC: C                    & ---                   & $-0.326^{***}$ (0.100) & --- \\
SD. Vaccine Delays (std)  & ---                   & ---                    & $ 1.303^{***}$ (0.073) \\
\midrule
Num. obs.       & 6162 & 6162 & 6162 \\
Log Likelihood  & -6101.771 & -6096.612 & -5615.522 \\
AIC             & 12215.543 & 12207.225 & 11245.044 \\
\bottomrule
\end{tabular}
\begin{tablenotes}[flushleft]
\footnotesize
\item Notes: Standard errors in parentheses. $^{*}p<0.1$, $^{**}p<0.05$, $^{***}p<0.01$.
\end{tablenotes}
\end{threeparttable}
\end{table}



\subsection{Baseline discounting estimates}
\label{subsec:baseline_discounting}

Using the baseline mixed logit model, we estimate a negative and statistically significant mean coefficient on vaccination delay, implying that longer waiting times reduce the probability of choosing a vaccine alternative. Interpreted as a marginal utility effect, the estimate indicates a sizable disutility associated with postponing protection by one month, holding constant efficacy, side effects, origin, and monetary incentives. The mixed logit results further confirm that this time sensitivity varies substantially across respondents, consistent with heterogeneous intertemporal preferences.

We then map the estimated delay sensitivity into standard discounting parameters as described in Section~\ref{ch:essay1_discounting} (Econometric specification and identification). The exponential specification yields an estimated discount parameter of
\[
\hat{\delta}=0.160,
\]
while the hyperbolic specification yields
\[
\hat{\kappa}=0.225.
\]
Figure~\ref{fig:discount_curves} compares the observed discount factors implied by the experimental delay trade-offs with the fitted exponential and hyperbolic curves over the observed delay horizon (0--6 months). Both specifications imply monotonic devaluation of delayed protection; however, the hyperbolic curve exhibits a sharper initial decline at short horizons. This pattern indicates pronounced short-run impatience: respondents place disproportionately high value on immediate access relative to protection that is postponed even briefly, consistent with present-biased valuation of vaccination timing.


\begin{figure}[!htbp]
\centering
\includegraphics[width=0.75\textwidth]{Overall Discounting_Observed vs Hyperbolic & Exponential fits (1).png}
\caption{Comparison of observed discounting with hyperbolic and exponential fits.}
\label{fig:discount_curves}
\end{figure}



\subsection{Behavioral Delay Multiplier (BDM)}
\label{subsec:bdm}

While $(\hat{\kappa},\hat{\delta})$ provide structurally interpretable discount parameters, they can be challenging to interpret in terms of the subjective experience of waiting. To provide a more intuitive behavioral translation, we utilize the \textbf{Behavioral Delay Multiplier (BDM)} as defined in Chapter~\ref{ch:common_methods} (Section~\ref{sec:backbone_econometrics}). Recall that the BDM is defined as:
\begin{equation}
\text{BDM}=\frac{1}{1-\hat{\kappa}}.
\label{eq:bdm_def_essay1}
\end{equation}
The BDM serves as an index of subjective time-inflation: it represents the ratio of \emph{perceived} to \emph{objective} waiting time implied by the hyperbolic discounting model. Larger BDM values indicate a greater "psychological time tax" associated with delay. 

Using the baseline estimate $\hat{\kappa}=0.225$, the implied BDM for the overall sample is \textbf{1.29}. This suggests that, on average, a given waiting period is experienced as being approximately 29\% longer than its objective duration (e.g., a 30-day administrative delay is perceived as a 38.7-day experiential burden).

To illustrate the distributional implications of this metric, we compute BDM values for selected policy-relevant subgroups (Table~\ref{tab:bdm_selected}). The results demonstrate that subjective time inflation is significantly higher among respondents with low institutional trust (BDM$=1.55$) compared to those with high institutional trust (BDM$=1.19$). This gap highlights a critical policy finding: institutional trust acts as a form of intertemporal "buffer," reducing the perceived burden of logistical friction. Similar elevations in BDM are observed among groups with irregular preventive routines and urban residents. These patterns suggest that identical administrative delays do not impose uniform burdens; rather, they generate unequal experiential costs across the population based on their underlying intertemporal preferences.

\begin{table}[!htbp]
\centering
\caption{Behavioral Delay Multiplier (BDM) for selected policy-relevant subgroups}
\label{tab:bdm_selected}
\begin{threeparttable}
\setlength{\tabcolsep}{7pt}
\begin{tabular}{lcc}
\toprule
Subgroup & $\hat{\kappa}$ & BDM $= \dfrac{1}{1-\hat{\kappa}}$ \\
\midrule
Overall sample & 0.225 & 1.29 \\
\addlinespace
Institutional trust: Low & 0.356 & 1.55 \\
Institutional trust: High & 0.160 & 1.19 \\
\addlinespace
Medical condition: Yes & 0.227 & 1.29 \\
Preventive routine proxy: Physical activity (Rare) & 0.308 & 1.45 \\
Residence: Urban & 0.290 & 1.41 \\
\bottomrule
\end{tabular}
\begin{tablenotes}[flushleft]
\footnotesize
\item Notes: BDM is defined as $\text{BDM}=1/(1-\hat{\kappa})$ (Eq.~\ref{eq:bdm_def}). 
$\hat{\kappa}$ values are taken from the subgroup hyperbolic discount estimates (Table~\ref{tab:kappa_delta_subgroups}). 
BDM can be interpreted as the ratio of perceived to objective waiting time implied by present-biased valuation; for example, BDM$=1.55$ implies that a 30-day delay is experienced as approximately 46.5 days.
\end{tablenotes}
\end{threeparttable}
\end{table}

\subsection{Heterogeneity in discounting}
\label{subsec:heterogeneity}



We next examine whether intertemporal discounting varies systematically across demographic characteristics, health behaviors, biological vulnerability, and institutional trust. This subsection reports subgroup-specific delay sensitivity and the implied discount parameters under both the hyperbolic and exponential mappings. Table~\ref{tab:kappa_delta_subgroups} presents the subgroup estimates, and Figure~\ref{fig:kappa_map} provides a visual summary of variation in $\hat{\kappa}$ across groups.


\paragraph{Demographic heterogeneity.}
Discounting differs meaningfully across demographic subgroups. Urban respondents exhibit substantially steeper discounting than rural respondents (higher $\hat{\kappa}$ and $\hat{\delta}$), indicating a stronger preference for immediate protection in urban settings. Age patterns are non-monotonic: the oldest group displays particularly steep discounting, while several middle-aged groups exhibit comparatively lower discount rates. Education-related differences are also evident, suggesting that time valuation of delayed protection varies across socioeconomic strata.

\paragraph{Behavioral routines and health-related heterogeneity.}
Intertemporal preferences also vary with health behaviors and preventive routines. Respondents with irregular preventive activity (e.g., rare physical activity) display markedly steeper discounting, consistent with a larger behavioral burden associated with waiting. In contrast, groups exhibiting more regular preventive routines show comparatively lower discount parameters, indicating greater tolerance for delay. In addition, respondents reporting chronic medical conditions exhibit higher discounting than those without such conditions, consistent with greater perceived biological cost of remaining unprotected during the waiting period.

\paragraph{Institutional trust.}
Institutional trust is among the strongest correlates of delay valuation. Subgroup estimates indicate that respondents with high trust have relatively lower discount parameters, implying greater willingness to wait for protection. By contrast, respondents expressing neutral or uncertain trust display substantially steeper discounting, consistent with heightened present-oriented valuation when confidence in future protection is ambiguous. This pattern aligns with the interpretation that institutional context shapes how delayed health benefits are perceived and discounted.

\paragraph{Summary of heterogeneity patterns.}
Taken together, Table~\ref{tab:kappa_delta_subgroups} and Figure~\ref{tab:kappa_delta_subgroups} show that discounting is not uniform across the population. The steepest short-run impatience (higher $\hat{\kappa}$) tends to cluster among groups with weaker institutional anchoring and less stable preventive routines, as well as among biologically vulnerable respondents. These systematic differences motivate the robustness checks that follow and provide a behavioral foundation for the subsequent essays that (i) monetize delay burdens in terms of compensation (Essay~2) and (ii) integrate time and money into equity-sensitive policy rules (Essay~3).



\begin{table}[!htbp]
\centering
\caption{Hyperbolic ($\kappa$) vs.\ Exponential ($\delta$) parameters across demographic, behavioral, and institutional-trust subgroups}
\label{tab:kappa_delta_subgroups}
\begin{threeparttable}
\setlength{\tabcolsep}{6pt}
\begin{tabular}{l r c c c}
\toprule
Group & $n$ & $\beta_{\text{delay}}$ (std) & $\beta_{\text{delay}}$ (per mo) & $\kappa$ / $\delta$ \\
\midrule
Overall & 1{,}027 & $-0.353^{***}$ (0.027) & $-0.160^{***}$ (0.012) & 0.225 / 0.160 \\
\addlinespace
\textit{Gender} \\
\quad Female & 545 & $-0.359^{***}$ (0.034) & $-0.163^{***}$ (0.015) & 0.230 / 0.163 \\
\quad Male   & 482 & $-0.346^{***}$ (0.036) & $-0.157^{***}$ (0.016) & 0.219 / 0.157 \\
\addlinespace
\textit{Residence} \\
\quad Rural & 555 & $-0.287^{***}$ (0.033) & $-0.130^{***}$ (0.015) & 0.174 / 0.130 \\
\quad Urban & 472 & $-0.433^{***}$ (0.038) & $-0.196^{***}$ (0.017) & 0.290 / 0.196 \\
\addlinespace
\textit{Age group} \\
\quad 18--34 & 217 & $-0.327^{***}$ (0.054) & $-0.148^{***}$ (0.024) & 0.204 / 0.148 \\
\quad 35--44 & 320 & $-0.420^{***}$ (0.042) & $-0.190^{***}$ (0.019) & 0.279 / 0.190 \\
\quad 45--54 & 231 & $-0.292^{***}$ (0.051) & $-0.132^{***}$ (0.023) & 0.177 / 0.132 \\
\quad 55--65 & 136 & $-0.237^{***}$ (0.059) & $-0.107^{***}$ (0.027) & 0.137 / 0.107 \\
\quad 65+    & 123 & $-0.477^{***}$ (0.070) & $-0.216^{***}$ (0.032) & 0.327 / 0.216 \\
\addlinespace
\textit{Education} \\
\quad No formal education & 65  & $-0.442^{***}$ (0.093) & $-0.201^{***}$ (0.042) & 0.298 / 0.201 \\
\quad Primary school      & 226 & $-0.500^{***}$ (0.053) & $-0.227^{***}$ (0.024) & 0.348 / 0.227 \\
\quad Middle school       & 264 & $-0.265^{***}$ (0.049) & $-0.120^{***}$ (0.022) & 0.158 / 0.120 \\
\quad High school         & 210 & $-0.277^{***}$ (0.052) & $-0.125^{***}$ (0.024) & 0.166 / 0.125 \\
\quad College/University  & 160 & $-0.282^{***}$ (0.058) & $-0.128^{***}$ (0.026) & 0.170 / 0.128 \\
\quad Graduate            & 102 & $-0.458^{***}$ (0.063) & $-0.208^{***}$ (0.029) & 0.311 / 0.208 \\
\bottomrule
\end{tabular}
\begin{tablenotes}[flushleft]
\footnotesize
\item Notes: Standard errors in parentheses. $n$ indicates the number of respondents per subgroup.
$p$-values are adjusted using Benjamini--Hochberg false-discovery rate ($q=0.10$); Holm--Bonferroni corrections yield similar results.
Significance levels: $^{*}p<0.1$, $^{**}p<0.05$, $^{***}p<0.01$.
\end{tablenotes}
\end{threeparttable}
\end{table}

\begin{table}[!htbp]
\ContinuedFloat
\centering
\caption[]{Hyperbolic ($\kappa$) vs.\ Exponential ($\delta$) parameters across subgroups (continued)}
\begin{threeparttable}
\setlength{\tabcolsep}{6pt}
\begin{tabular}{l r c c c}
\toprule
Group & $n$ & $\beta_{\text{delay}}$ (std) & $\beta_{\text{delay}}$ (per mo) & $\kappa$ / $\delta$ \\
\midrule
\textit{Smoking status} \\
\quad Never     & 311 & $-0.367^{***}$ (0.043) & $-0.167^{***}$ (0.020) & 0.236 / 0.167 \\
\quad Sometimes & 402 & $-0.392^{***}$ (0.039) & $-0.178^{***}$ (0.018) & 0.256 / 0.178 \\
\quad Often     & 314 & $-0.290^{***}$ (0.044) & $-0.131^{***}$ (0.020) & 0.176 / 0.131 \\
\addlinespace
\textit{Physical activity} \\
\quad Rare        & 297 & $-0.455^{***}$ (0.046) & $-0.206^{***}$ (0.021) & 0.308 / 0.206 \\
\quad Sometimes   & 439 & $-0.304^{***}$ (0.037) & $-0.138^{***}$ (0.017) & 0.186 / 0.138 \\
\quad Quite often & 291 & $-0.325^{***}$ (0.043) & $-0.147^{***}$ (0.020) & 0.202 / 0.147 \\
\addlinespace
\textit{Medical condition} \\
\quad No  & 781 & $-0.300^{***}$ (0.029) & $-0.136^{***}$ (0.013) & 0.183 / 0.136 \\
\quad Yes & 119 & $-0.355^{***}$ (0.070) & $-0.161^{**}$  (0.032) & 0.227 / 0.161 \\
\addlinespace
\textit{Healthcare access} \\
\quad No access & 460 & $-0.335^{***}$ (0.037) & $-0.152^{***}$ (0.017) & 0.211 / 0.152 \\
\quad Access    & 567 & $-0.368^{***}$ (0.034) & $-0.167^{***}$ (0.015) & 0.237 / 0.167 \\
\addlinespace
\textit{Prior vaccine reaction} \\
\quad No  & 950 & $-0.352^{***}$ (0.027) & $-0.159^{***}$ (0.012) & 0.224 / 0.159 \\
\quad Yes & 77  & $-0.372^{***}$ (0.092) & $-0.168^{**}$  (0.042) & 0.240 / 0.168 \\
\addlinespace
\textit{Institutional trust} \\
\quad Low Trust  & 487 & $-0.510^{***}$ (0.052) & $-0.231^{***}$ (0.023) & 0.356 / 0.231 \\
\quad High Trust & 446 & $-0.268^{***}$ (0.054) & $-0.122^{***}$ (0.025) & 0.160 / 0.122 \\
\bottomrule
\end{tabular}
\end{threeparttable}
\end{table}


\clearpage



\begin{figure}[!htbp]
\centering
\includegraphics[width=0.98\textwidth]{Subgroup Variation in Hyperbolic Discount Parameters.png}
\caption{Subgroup variation in the hyperbolic discount parameter $\hat{\kappa}$. Larger values indicate steeper short-run discounting (greater present-oriented valuation of immediate protection).}
\label{fig:kappa_map}
\end{figure}


\clearpage

\section{Discussion and policy implications}
\label{sec:essay1_discussion}

\paragraph{Summary of main findings.}
This essay provides three main findings. First, delayed protection is systematically discounted in vaccination timing preferences, and the observed discounting profile is better characterized by a present-biased (hyperbolic) mapping than by a time-consistent exponential benchmark \cite{Attema2018,vanDerPolCairns2001}. Second, the Behavioral Delay Multiplier (BDM) offers an intuitive translation of the hyperbolic parameter into subjective waiting-time inflation, clarifying the behavioral meaning of estimated impatience. Third, discounting is systematically heterogeneous across demographic characteristics, health behaviors, biological vulnerability, and institutional trust, implying that identical administrative delays can impose unequal experiential burdens across groups.

\subsection*{Interpretation in a vaccination-timing context}
Vaccination timing is not merely a binary choice between uptake and refusal; it is an intertemporal decision in which individuals trade off immediate frictions against delayed and uncertain protection \cite{MacDonald2015,Larson2014}. In applied health economics, discounting parameters are typically treated as summary representations of how future health benefits are devalued over time, reflecting standard practices in economic evaluation where delayed health gains are adjusted to present value, and in discrete choice studies that accommodate nonlinear time preferences in health valuation \cite{Hultkrantz2021,Jonker2018}.In the present setting, the key object is the valuation of \emph{delayed onset of protection} (i.e., waiting time), recovered from experimentally randomized trade-offs within the DCE. The sharper short-horizon decline captured by the hyperbolic mapping is consistent with a behavioral discounting profile often documented in delay-discounting research, where short delays carry disproportionate weight relative to longer horizons \cite{Bickel1999}.

\subsection*{Why heterogeneity matters}
A central contribution of this essay is documenting structured heterogeneity in discounting. Subgroup patterns are not simply descriptive: they indicate that the behavioral cost of waiting is unevenly distributed, which has implications for both uptake forecasting and equity in access. Methodologically, allowing delay sensitivity to vary (via random coefficients) and examining observed heterogeneity via interactions provides a unified framework for comparing groups under the same identification strategy \cite{McFadden1974,Train2009}. Conceptually, these results align with evidence that time-preference parameters measured in survey and choice settings vary systematically across populations and contexts, rather than reflecting a single stable trait \cite{Bechtel2023}.

\subsection*{Implications for vaccination program design}
The findings suggest three practical implications for program implementation. First, because short delays are discounted steeply, operational improvements that reduce \emph{near-term} waiting (e.g., appointment availability and queue management) may yield disproportionately large behavioral returns compared with reducing already-long delays, consistent with convex perceived costs of waiting \cite{Attema2018}. Second, institutional context matters: the strong association between trust and delay valuation implies that timing policies cannot be evaluated independently of perceived credibility and expectations about future access \cite{Larson2014}. Third, because delay burdens are heterogeneous, uniform waiting times may be inequitable in experiential terms even when they are administratively equal, motivating a need to translate time-based burdens into comparable monetary metrics and evaluate alternative instruments under explicit equity objectives.

\subsection*{Limitations and scope}
Two limitations are worth noting. First, the discount parameters are recovered from stated-preference trade-offs and therefore reflect valuation of delayed protection within the experimental framing; external validity depends on how closely this environment matches real-world constraints and risk perceptions \cite{Bechtel2023}. Second, mapping delay sensitivity into exponential or hyperbolic forms imposes parsimonious functional-form assumptions that are analytically convenient but may not fully capture the richness of observed intertemporal preferences \cite{Frederick2002,Halevy2015}.
 Robustness analyses assess sensitivity to alternative codings and specifications.

\subsection*{Link to subsequent essays}
Essay~1 establishes that delays impose substantial behavioral costs and that these costs are systematically heterogeneous. However, the results remain expressed in time-preference units. Essay~2 translates these preferences into monetary terms by estimating compensation required to offset delayed protection (MWTA), and Essay~3 integrates time and money to derive equity-sensitive policy rules for choosing between speeding up access and compensating delays.




\section{Conclusion}
\label{sec:essay1_conclusion}

This essay demonstrates that vaccination timing is systematically shaped by intertemporal discounting. Using experimentally identified delay sensitivity from a discrete choice framework, we recover both exponential and hyperbolic discount parameters. The evidence indicates economically meaningful short-run impatience, with the hyperbolic specification better capturing the sharp initial devaluation of delayed protection.

Importantly, the Behavioral Delay Multiplier (BDM) provides an intuitive behavioral interpretation of this effect: waiting time is experienced as subjectively longer than its objective duration. Moreover, discounting is not uniform across the population. Respondents with weaker institutional trust, less stable preventive routines, and greater biological vulnerability exhibit steeper short-run impatience, implying that identical administrative delays may generate unequal experiential burdens.

These findings establish vaccination timing as an intertemporal health investment problem with structured heterogeneity. However, the results thus far remain expressed in time-based preference parameters. They do not yet quantify how much compensation would offset a given delay, nor how time and monetary instruments should be combined in policy design.

The next essay addresses this gap by translating delay sensitivity into monetary terms, estimating the willingness to accept compensation for postponed protection, and constructing a policy-relevant trade-off menu between speeding up access and offering incentives.



\clearpage

\chapter{Essay 2: The Monetary Valuation of Vaccination Delay}
\label{sec:essay2}

\section{Introduction}
\label{subsec:essay2_overview}

Vaccination is widely recognized as a cornerstone of infectious-disease prevention, yet real-world vaccination behavior is not simply a binary choice between uptake and refusal. A central policy challenge is heterogeneity in \emph{timing}: many individuals delay vaccination even when vaccines are available, thereby remaining exposed to infection risk for longer periods and weakening the timeliness of immunization coverage \cite{Hu2014,Smith2010}.
 Such delays are often shaped by operational frictions---phased eligibility, appointment scarcity, and queue congestion---and they interact with policy instruments such as financial incentives that aim to accelerate uptake \cite{Brewer2017}. For program design, the key question is not only whether incentives raise vaccination, but how much compensation is required to offset the behavioral cost of waiting.

A growing behavioral-economics literature documents steep near-term discounting and present-biased decision-making, implying that short delays can be disproportionately deterrent in preventive-health choices \cite{Frederick2002,Laibson1997,OR1999}. Recent evidence further links impatience and uncertainty discounting to vaccination intentions and uptake patterns \cite{Guillon2024}. However, existing work rarely converts delay aversion into a policy-ready monetary metric that can be compared directly with incentive budgets or operational investments. As a result, policymakers lack quantitative guidance on the \emph{shadow price} of vaccination delay---that is, the compensation required to maintain uptake when waiting times cannot be eliminated.

This chapter fills that gap by monetizing vaccination delay using a discrete choice experiment (DCE) administered to 1{,}027 adults in Wuhan, China. The core estimand is the \emph{marginal willingness-to-accept} (MWTA): the RMB compensation required for respondents to tolerate an additional month of waiting for vaccine-induced protection. Beyond average MWTA, the chapter documents heterogeneity in delay valuation across demographic characteristics, health status and biological vulnerability, and institutional trust, highlighting groups for whom delay is especially costly and for whom targeted policies may be most effective. Finally, the chapter translates estimated preference parameters into program-planning tools by simulating uptake under alternative delay--incentive bundles and summarizing these trade-offs with iso-uptake (iso-acceptance) curves.

The remainder of the chapter proceeds as follows. Section~\ref{subsec:essay2_data} summarizes the data and DCE design, with full experimental protocols and instrument materials provided in Chapter~2. Section~\ref{sec:essay2_model} presents the econometric framework and the mapping from estimated coefficients to MWTA. Section~\ref{sec:essay2_results} reports the main findings, including monetized delay valuation, subgroup heterogeneity, and policy-bundle simulations with iso-acceptance visualizations. Section~\ref{sec:essay2_validity} offers a concise summary of validity and robustness checks, with detailed diagnostics and supplementary tables reported in Chapter~\ref{sec:essay2_validity}. Section~\ref{sec:essay2_discussion} discusses implications for programme design when balancing operational speed-ups against monetary incentives, and Section~\ref{sec:essay2_limitations} concludes with limitations and directions for future research.



\section{Motivation and contribution}
\label{subsec:essay2_motivation}

Timely vaccination is a central operational objective in epidemic preparedness because delays extend the period of exposure to infection risk, slow population-level protection, and can erode the effectiveness of staged rollouts. Yet real-world programmes routinely confront timing frictions---phased eligibility, appointment congestion, and supply fluctuations---that push protection into the future even when vaccines are free and nominally accessible. At the same time, policymakers increasingly consider monetary incentives to accelerate uptake when timeliness is threatened. The resulting design problem is inherently economic: when access is delayed, should programmes \emph{invest in speeding up delivery} (e.g., staffing, outreach, clinic hours) or \emph{compensate individuals for waiting} (e.g., cash payments, vouchers), and how should this trade-off vary across populations?

Despite extensive evidence on vaccine hesitancy and the behavioral foundations of preventive-health decisions, existing work provides limited guidance on the \emph{price of waiting} in vaccination. Research on intertemporal choice in health consistently documents present bias, near-term impatience, and uncertainty discounting, which can make even modest delays behaviorally costly \cite{Laibson1997,ODonoghueRabin2001,AttemaBleichrodt2013,Attema2018}. Separately, field evaluations and randomized trials show that monetary incentives can increase vaccination uptake, though effect sizes vary by context, baseline acceptance, and incentive design \cite{CamposMercade2021,Yip2022,Kim2022,Li2023,Shen2024}. However, these two strands are rarely integrated into a single framework that quantifies how much compensation is required to offset a given delay, especially over the operationally relevant horizon of months rather than days. As a result, policymakers lack a transparent monetary metric for comparing the behavioral cost of delayed access to the fiscal cost of incentives or delivery-system improvements.

This chapter addresses that gap by monetizing vaccination delay using marginal willingness-to-accept (MWTA) estimates derived from a discrete choice experiment (DCE) fielded among 1,027 adults in Wuhan, China. The Wuhan context provides an analytically useful setting for studying delay--cash trade-offs under high-salience epidemic conditions, in which waiting times are both credible and behaviorally meaningful. At the same time, the stated-preference design allows the analysis to isolate \emph{preference-based} delay aversion from contemporaneous system constraints---such as appointment bottlenecks, supply shocks, and information frictions---that can confound revealed-preference measures of timeliness. The chapter complements Essay~1, which maps delay sensitivity into reduced-form discounting measures, by expressing the valuation of timeliness in directly policy-actionable monetary units. It also sets up Essay~3 by providing monetized inputs that can be combined with discounting heterogeneity to derive equity-sensitive rules for deciding when to accelerate access versus when to offer compensation for unavoidable waiting.


The chapter makes three contributions. First, it provides the first robust estimates of the \emph{monetary valuation of vaccination delay} in this context by converting estimated delay sensitivity into MWTA (RMB per month). This ``shadow price of waiting'' translates abstract time preference into a metric that can be directly used for incentive calibration and programme budgeting. Second, it documents systematic \emph{heterogeneity} in delay valuation across demographic characteristics, health behaviors, medical vulnerability, and institutional trust, identifying populations for whom delays are especially costly and for whom targeted prioritization or incentives may yield higher returns \cite{Liu2021}. Third, it operationalizes the estimated preference parameters into a decision-support tool by simulating predicted uptake under delay--cash bundles and constructing \emph{iso-acceptance curves}---frontiers of delay and compensation that sustain target uptake thresholds relevant for programme planning \cite{Abbas2021}. Together, these outputs provide a practical behavioral-economic framework for evaluating timeliness policies during future epidemics, bridging micro-level time preferences and macro-level delivery design.

\section{Theoretical Framework}
\label{subsec:essay2_framework}

Vaccination timing can be conceptualized as an intertemporal health investment decision in which individuals trade off immediate, salient costs—time, inconvenience, effort, and anticipated side effects—against protection that is realized in the future. A large literature in behavioral and health economics shows that preferences often exhibit present bias, so near-term costs are overweighted relative to delayed benefits \cite{Laibson1997,ODonoghueRabin2001}. In vaccination contexts, this implies that even modest postponements can materially reduce perceived utility and contribute to delay despite substantial eventual protection \cite{Guillon2024}. In economic terms, these forces can be summarized as a positive shadow price of timeliness: individuals require compensation to accept delayed access, and the implied compensation rate provides a policy-relevant monetary valuation of waiting.

Uncertainty can further magnify the disutility of delay. When epidemic trajectories, exposure risk, or realized vaccine performance are ambiguous, individuals tend to discount delayed outcomes more steeply \cite{AttemaBleichrodt2013,Andersen2018}. Institutional trust is central in this channel: lower trust can heighten perceived uncertainty about whether protection will be delivered as promised or whether its effectiveness will persist, thereby strengthening impatience toward waiting \cite{Liu2021}. Crucially, these preference-based mechanisms are distinct from structural frictions—such as eligibility rules, appointment congestion, and supply constraints—that also generate delays in practice. This distinction motivates the stated-preference approach adopted in this chapter. By experimentally varying delay and monetary incentives in a controlled choice environment, the design identifies the \emph{behavioral valuation of waiting} net of contemporaneous system constraints that complicate inference from revealed-preference data. The resulting marginal willingness-to-accept (MWTA)—the marginal rate of substitution between waiting time and monetary compensation—thus provides a reduced-form monetary measure of the behavioral cost of delay over operationally relevant monthly horizons, without requiring identification of structural $\beta$--$\delta$ parameters.


\section{Data and experimental design}
\label{subsec:essay2_data}

Essay~2 relies on the same discrete choice experiment (DCE) described in Essay~1 and in the dissertation’s common methods chapter. Here I briefly summarize the data and experimental features most relevant for monetizing vaccination delay and for the policy simulations developed in this chapter; full details on instrument development, sampling, and implementation are reported once in the shared DCE backbone to avoid repetition.

\subsubsection{Setting and fielding context}
The DCE was fielded among adults in Wuhan, China between March and May 2025. Respondents were asked to consider vaccination against a hypothetical ``COVID-like disease'' (CLD), defined as a future respiratory infection with transmissibility and severity similar to COVID-19. This neutral framing reduces anchoring to a specific COVID-19 wave, variant, or vaccine brand and helps interpret the estimated parameters as valuation of timeliness and other vaccine attributes under epidemic-like conditions. Wuhan provides a high-salience context in which waiting time for vaccination is credible and behaviorally meaningful, while the stated-preference design allows preferences to be elicited under controlled access conditions.

\subsubsection{Choice task structure}
Each respondent completed six choice tasks. In each task, individuals chose between two hypothetical vaccine profiles (A and B) and an opt-out alternative. Including an opt-out option enhances realism by permitting non-uptake when neither vaccine is acceptable. Vaccine profiles varied along multiple attributes, enabling identification of trade-offs between waiting time and monetary incentives that form the basis of the MWTA estimates in this chapter.

\subsubsection{Development of attributes and levels}
Attribute development followed established stated-preference good practices and was harmonized across Essays~1--3. Candidate attributes were identified from the vaccination behavior and incentive-policy literature, refined through consultation and pilot testing, and finalized to reflect policy-relevant ranges. The final experiment included five attributes:
\begin{itemize}
  \item \textbf{Vaccination delay}: 0, 1, 3, and 6 months (waiting time until the vaccine becomes available);
  \item \textbf{Monetary incentives}: 0, 50, 200, and 800 RMB (received upon vaccination);
  \item \textbf{Vaccine efficacy}: 50\%, 70\%, and 95\%;
  \item \textbf{Side-effect severity}: mild, moderate, and severe (modeled as an ordered attribute);
  \item \textbf{Vaccine origin}: domestic versus imported.
\end{itemize}
Delay levels were specified in monthly increments to match operationally relevant postponements arising from phased eligibility, appointment congestion, and supply fluctuations. Incentive levels were selected to span small, moderate, and large payments that are straightforward to interpret for program design and budget planning. These design choices ensure that the MWTA estimates correspond to the practical planning horizon over which vaccination programs adjust scheduling capacity and incentive schemes.

\subsubsection{Sample and representativeness}
The analytic sample includes $N=1{,}027$ respondents. Quota sampling on key demographics (e.g., age, gender, and district of residence) was used to approximate the composition of Wuhan’s adult population, supporting subgroup analyses central to this chapter. Summary comparisons to census benchmarks and additional sampling details are provided in the shared DCE appendix to maintain consistency across essays.

\paragraph{Data structure for estimation.}
Because each respondent completes multiple choice tasks, the dataset has a short panel structure with repeated choices at the individual level. All econometric models in this chapter account for within-respondent correlation (e.g., clustered standard errors at the individual level), and the same coding conventions for attributes used in Essay~1 are maintained here to ensure comparability across essays.

\paragraph{Validity checks.}
Internal validity diagnostics (e.g., dominance tests, repeated-choice consistency, and sensitivity to attribute non-attendance) are implemented using the same procedures as in Essay~1, Section~\ref{sec:validity_robustness} Chapter 2 and the dissertation appendices. External-comparison checks and robustness analyses are likewise harmonized across essays to support a coherent empirical narrative.


\section{Econometric specification and identification}
\label{sec:essay2_model}

This chapter uses the same random-utility framework and choice-data structure as Essay~1, but shifts the estimand from discount-parameter mappings to a monetized valuation of delay. The key object of interest is the marginal willingness-to-accept (MWTA) for vaccination delay, defined as the monetary compensation required to offset the disutility of an additional month of waiting. To ensure comparability across essays, the attribute coding conventions and baseline model structure follow the shared DCE backbone.

\subsection{Random utility framework}
Let individual $i \in \{1,\dots,N\}$ face choice tasks $t \in \{1,\dots,T_i\}$ and choose among alternatives $j \in \{A,B,O\}$, where $A$ and $B$ denote the two vaccine profiles and $O$ denotes the opt-out alternative. Utility is assumed to take the standard additive random-utility form:
\begin{equation}
U_{ijt} = V_{ijt} + \varepsilon_{ijt},
\end{equation}
where $V_{ijt}$ is the systematic component and $\varepsilon_{ijt}$ is an unobserved error term.

The deterministic utility is specified as a linear index in vaccine attributes:
\begin{equation}
V_{ijt} = \alpha_j + \beta_{d}\,\textit{Delay}_{ijt}
+ \beta_{c}\,\textit{Cash}_{ijt}
+ \beta_{e}\,\textit{Efficacy}_{ijt}
+ \beta_{s}\,\textit{SideEffects}_{ijt}
+ \beta_{o}\,\textit{Imported}_{ijt},
\label{eq:essay2_Vijt}
\end{equation}
where $\alpha_j$ are alternative-specific constants (ASCs) capturing average preference for selecting a vaccine relative to opting out (with $\alpha_O \equiv 0$ for normalization). $\textit{Delay}$ is measured in months, $\textit{Cash}$ is the monetary incentive (typically scaled per 100 RMB for interpretability), $\textit{Efficacy}$ is coded on a 0--1 scale (and rescaled for interpretation where useful), $\textit{SideEffects}$ captures side-effect severity (entered either ordinally or as dummies), and $\textit{Imported}$ is an indicator for imported origin.

\subsection{Benchmark models (MNL and conditional logit)}
Assuming $\varepsilon_{ijt}$ are i.i.d.\ Type-I extreme value, the probability that respondent $i$ chooses alternative $j$ in task $t$ follows the multinomial logit (MNL) form:
\begin{equation}
P_{ijt} = \Pr(y_{it}=j) =
\frac{\exp(V_{ijt})}{\sum_{k \in \{A,B,O\}} \exp(V_{ikt})}.
\label{eq:essay2_mnl}
\end{equation}
Parameters are estimated by maximum likelihood using the stacked choice-task data, with standard errors clustered at the respondent level to account for repeated choices.

This benchmark specification provides a transparent mapping from attribute coefficients to marginal trade-offs (e.g., between delay and cash), which is essential for MWTA estimation and for simulating uptake under counterfactual delay--incentive bundles. When the data are arranged in long format with alternative-varying attributes, the MNL is numerically equivalent to the conditional logit formulation commonly used in stated-preference applications.

\subsection{Mixed logit specification}
To relax the IIA property and allow for unobserved preference heterogeneity, I also estimate mixed logit (MXL) specifications. In the MXL model, one or more coefficients are allowed to vary across individuals:
\begin{equation}
\beta_i \sim f(\beta \mid \theta),
\end{equation}
where $f(\cdot)$ is a parametric distribution (e.g., normal) characterized by parameters $\theta$. Utility becomes
\begin{equation}
U_{ijt} = \alpha_j + X_{ijt}'\beta_i + \varepsilon_{ijt}.
\end{equation}
Choice probabilities integrate over the distribution of random coefficients:
\begin{equation}
P_{ijt} = \int
\frac{\exp(\alpha_j + X_{ijt}'\beta)}
{\sum_{k}\exp(\alpha_k + X_{ikt}'\beta)}
\, f(\beta \mid \theta)\, d\beta,
\label{eq:essay2_mxl}
\end{equation}
and are approximated via simulation.

Given the short panel structure (six tasks per respondent), the MXL specification is parameterized parsimoniously to ensure empirical stability (e.g., allowing random coefficients on delay and/or cash while keeping other attributes fixed). The MXL is used primarily as a robustness check for the benchmark MNL results and to assess sensitivity to unobserved heterogeneity.

\subsection{Derivation of MWTA}
The marginal willingness-to-accept (MWTA) for a one-month increase in vaccination delay is defined as the monetary compensation that keeps indirect utility constant for a marginal change in delay:
\begin{equation}
MWTA_{\textit{delay}} = -\frac{\beta_{d}}{\beta_{c}}.
\label{eq:essay2_mwta}
\end{equation}
MWTA is reported in RMB per month (accounting for any scaling of $\textit{Cash}$ in estimation). Uncertainty for MWTA is summarized using delta-method standard errors and confidence intervals, and (in robustness checks) by nonparametric bootstrapping clustered at the respondent level.

MWTA is a reduced-form monetary valuation of waiting over the experiment’s monthly horizon. It does not identify structural discounting parameters (e.g., $\beta$--$\delta$) because the design elicits static trade-offs over discrete, operationally relevant waiting times rather than dynamic intertemporal allocations. Its value for policy is interpretability: MWTA places delay on the same monetary scale as incentive payments and delivery costs, enabling direct comparisons between ``speed-up'' and ``compensation'' strategies.

\subsection{Policy-bundle simulations and iso-acceptance curves}
To translate the estimated preference parameters into quantities relevant for programme design, I simulate predicted uptake under counterfactual policy bundles that vary vaccination delay and incentive amounts. For each bundle, I compute the predicted probability of choosing vaccination (relative to opting out) by evaluating \eqref{eq:essay2_mnl} (or \eqref{eq:essay2_mxl}) at the bundle’s attribute levels while holding other attributes at a reference profile. Repeating this calculation over a grid of $(\textit{Delay}, \textit{Cash})$ combinations yields an uptake surface that can be visualized as iso-acceptance curves (or heatmaps). These curves represent the combinations of delay and compensation that sustain a target uptake threshold, providing an interpretable decision-support frontier for comparing ``speed-up'' policies (reducing delay) to ``compensation'' policies (increasing incentives).

\subsection{Identification}
Identification of $\beta_d$ and $\beta_c$ arises from the experimental variation in delay and incentive levels across alternatives and tasks, orthogonal to respondent characteristics by design. Under standard DCE assumptions (stable preferences over tasks, respondents attend to attributes, and the random-utility structure), the delay and cash coefficients capture marginal utilities that can be interpreted as compensatory trade-offs. The MWTA estimand in \eqref{eq:essay2_mwta} is identified as a ratio of these coefficients. As in Essay~1, I assess the credibility of the preference data using internal validity diagnostics (dominance tests, repeated-task consistency, and sensitivity to attribute non-attendance) and report robustness to alternative specifications.

\section{Main results}
\label{sec:essay2_results}


\subsection{Attribute effects}
\label{subsec:essay2_attribute_effects}

Table~\ref{tab:main_models_std} reports the main preference estimates from the multinomial logit (MNL) and mixed logit (MXL) models using standardized attributes. Across both specifications, the estimated coefficients exhibit the expected signs and are statistically significant at conventional levels. Vaccination delay enters negatively, indicating that longer waiting times reduce utility and lower the probability of vaccine uptake. Monetary incentives enter positively, implying that financial compensation increases the attractiveness of vaccination. Higher vaccine efficacy increases utility, whereas more severe side effects substantially decrease it. 

The MXL specification relaxes the independence-of-irrelevant-alternatives assumption and allows for unobserved heterogeneity in key attributes. The mean effects are consistent with the MNL estimates, supporting the stability of the benchmark results. Estimated standard deviations for delay and cash suggest meaningful heterogeneity in time and incentive preferences, motivating the subgroup analysis that follows.

\begin{table}[htbp]
\centering
\caption{Main model estimates (standardized attributes): Multinomial Logit (MNL) and Mixed Logit (MXL)}
\label{tab:main_models_std}
\begin{threeparttable}
\begin{tabular}{lcc}
\toprule
 & MNL & MXL (Mean) \\
\midrule
Vaccination delay (months) & -0.412*** (0.031) & -0.398*** (0.045) \\
Monetary incentive (RMB, standardized) & 0.523*** (0.037) & 0.508*** (0.054) \\
Vaccine efficacy & 0.684*** (0.049) & 0.671*** (0.062) \\
Side-effect severity & -0.731*** (0.052) & -0.715*** (0.066) \\
Imported origin & 0.128** (0.058) & 0.120** (0.061) \\
Alternative-specific constant (vaccine) & 0.842*** (0.071) & 0.815*** (0.083) \\
\midrule
Std. dev. (Delay) & -- & 0.192** (0.081) \\
Std. dev. (Cash) & -- & 0.214** (0.093) \\
\midrule
Observations & 18,486 & 18,486 \\
Individuals & 1,027 & 1,027 \\
Log-likelihood & -10,842 & -10,615 \\
\bottomrule
\end{tabular}
\begin{tablenotes}
\footnotesize
\item Notes: Standard errors in parentheses. MNL standard errors are clustered at the individual level. MXL estimates are based on 500 simulated draws. *** $p<0.01$, ** $p<0.05$, * $p<0.1$.
\end{tablenotes}
\end{threeparttable}
\end{table}
\clearpage

\subsection{Monetized delay aversion}
\label{subsec:essay2_monetized}

To translate the estimated preferences into policy-relevant monetary units, I convert the unstandardized MNL coefficients into cash--attribute trade-offs. The central quantity of interest is the marginal willingness-to-accept (MWTA) for vaccination delay, defined as the monetary compensation required to offset the utility loss from an additional month of waiting:
\begin{equation}
MWTA_{\textit{delay}} \;=\; -\frac{\beta_{\textit{delay}}}{\beta_{\textit{cash}}}.
\label{eq:mwta}
\end{equation}
I also report the corresponding willingness-to-wait (WTW), which expresses the same trade-off in time units (months) per unit of cash:
\begin{equation}
WTW_{x} \;=\; -\frac{\beta_{x}}{\beta_{\textit{delay}}}.
\label{eq:wtw}
\end{equation}
Because MWTA and WTW are nonlinear ratios of estimated coefficients, statistical uncertainty is summarized using cluster-robust standard errors and 95\% confidence intervals constructed via the delta method at the respondent level.

Table~\ref{tab:tradeoffs_mwta_wtw} reports the resulting trade-off estimates. Because MWTA is a monetary trade-off defined in the original units of time and currency, it must be computed from \emph{unstandardized} coefficients; standardizing attributes rescales the parameters by their sample variability and would therefore distort the RMB-per-month interpretation of delay valuation.
The MWTA for a one-month vaccination delay is \textbf{56.35 RMB} (95\% CI: \textbf{35.10--77.60}), indicating a sizable shadow price of timeliness. Interpreted over operationally relevant horizons, this estimate implies that a three-month postponement corresponds to a behavioral cost of approximately 169 RMB (and about 338 RMB for a six-month postponement), holding other vaccine attributes constant. The inverse trade-off suggests that monetary incentives can partially compensate for delayed access: expressed as WTW, cash incentives translate into additional tolerated waiting time, providing an interpretable exchange rate between compensation and delay.

Table~\ref{tab:tradeoffs_mwta_wtw} also monetizes other vaccine attributes on the same scale. A 10-percentage-point increase in vaccine efficacy is valued at \textbf{43.97 RMB}, while an increase in side-effect severity by one level (mild $\rightarrow$ moderate $\rightarrow$ severe) requires approximately \textbf{70.64 RMB} in compensation. Together, these results show that waiting time is a quantitatively meaningful component of vaccine utility and provide the monetary inputs needed for the subgroup analyses and policy-bundle simulations that follow.

\begin{table}[htbp]
\centering
\setlength{\tabcolsep}{4pt}
\small
\caption{Trade-off estimates (MNL, unstandardized): MWTA and WTW with cluster-robust 95\% confidence intervals}
\label{tab:tradeoffs_mwta_wtw}
\begin{threeparttable}
\begin{tabular}{lccc|ccc}
\toprule
& \multicolumn{3}{c}{MWTA} & \multicolumn{3}{c}{WTW} \\
\cmidrule(lr){2-4}\cmidrule(lr){5-7}
Attribute & RMB & SE & 95\% CI & months & SE & 95\% CI \\
\midrule
Vaccine delays (per month) & 56.35 & 10.84 & [35.10,\; 77.60] & -1.00 & 0.17 & [-1.33,\; -0.67] \\
Monetary incentives (per 100 RMB) & -100.00 & 0.06 & [-1.11,\; -0.89] & 2.00 & 0.00 & [0.01,\; 0.02] \\
Vaccine efficacy (per 10 pp) & -43.97 & 11.40 & [-66.31,\; -21.63] & 0.78 & 0.18 & [0.43,\; 1.13] \\
Side effects (per severity level) & 70.64 & 27.13 & [17.47,\; 123.81] & -1.25 & 0.49 & [-2.21,\; -0.30] \\
Origin (Imported) & -237.93 & 43.45 & [-323.08,\; -152.78] & 4.22 & 0.84 & [2.58,\; 5.86] \\
\bottomrule
\end{tabular}
\begin{tablenotes}
\footnotesize
\item Notes: MWTA (marginal willingness-to-accept) and WTW (willingness-to-wait) are computed as $-\beta_x/\beta_{\textit{cash}}$ and $-\beta_x/\beta_{\textit{delay}}$, respectively, using unstandardized MNL coefficients. Cluster-robust standard errors and 95\% confidence intervals are calculated at the respondent level using the delta method. Vaccine efficacy was coded 0--1 and rescaled for interpretability to 10-percentage-point increments. Side effects represent the change per increase in severity level (mild $\rightarrow$ moderate $\rightarrow$ severe). Confidence intervals convey statistical significance; $p$-values are omitted because MWTA and WTW are nonlinear ratios of estimated coefficients.
\end{tablenotes}
\end{threeparttable}
\end{table}


\clearpage

\subsection{Heterogeneity in delay valuation}
\label{subsec:essay2_heterogeneity}

While the overall MWTA provides a useful benchmark, substantial heterogeneity emerges across demographic, behavioral, and institutional dimensions. Table~\ref{tab:mwta_subgroups} reports subgroup-specific marginal willingness-to-accept (MWTA) for a one-month increase in vaccination delay, derived from multinomial logit models with vaccination delay and cash incentives entered linearly. Confidence intervals are cluster-robust at the respondent level and constructed using the delta method.

For the overall sample, the MWTA is 44.8 RMB per month (95\% CI: 28.9--60.7), indicating that respondents require approximately this amount of compensation to tolerate an additional month of waiting. However, this average masks meaningful gradients across groups.

\textit{Demographics.} Women exhibit a higher MWTA (55.7 RMB) than men (33.1 RMB), and urban residents display higher delay valuation (57.7 RMB) than rural residents (34.7 RMB). Education and age differences also emerge: lower-education respondents (51.2 RMB) and younger respondents (<55 years; 49.8 RMB) value timeliness more than higher-education respondents (28.1 RMB) and older respondents ($\geq$55 years; 30.8 RMB), although confidence intervals partially overlap.

\textit{Health behavior.} Behavioral heterogeneity is particularly pronounced. Smokers require 55.0 RMB per month to accept delay, compared with 38.4 RMB among non-smokers. Differences across physical activity groups are substantial: individuals with low physical activity exhibit the highest MWTA (80.1 RMB; 95\% CI: 44.6--115.6), nearly three times that of the mid-activity group (27.6 RMB; 95\% CI: 5.1--50.1).

\textit{Institutional trust.} The most striking contrast appears along institutional trust. Respondents with low institutional trust display an MWTA of 62.3 RMB (95\% CI: 37.1--87.5), more than double that of high-trust respondents (28.7 RMB; 95\% CI: 6.5--50.8). This pattern is consistent with the conceptual framework in which lower trust heightens perceived uncertainty regarding delayed protection, thereby increasing impatience toward waiting.

\textit{Medical condition and self-rated health.} Individuals reporting a medical condition exhibit higher MWTA (57.5 RMB) relative to those without (43.4 RMB), though the confidence interval is wider due to the smaller subgroup size. Differences by self-rated health are comparatively modest (44.8 RMB vs 44.9 RMB).

Figure~\ref{fig:mwta_forest} provides a compact visualization of these subgroup differences. Point estimates and confidence intervals are plotted against a vertical reference line marking the overall MWTA, making clear which groups place above-average monetary value on timeliness. The figure highlights especially elevated delay costs among low-trust respondents and those with low physical activity, reinforcing the conclusion that the welfare burden of delay is unevenly distributed.

Taken together, these results indicate that the behavioral cost of delay is far from uniform. Delay schedules that appear moderate at the aggregate level may impose disproportionate welfare burdens on specific subpopulations, particularly those with lower institutional trust or lower preventive engagement. These heterogeneity patterns motivate targeted prioritization or compensation strategies, explored further in the policy simulations below.

\begin{table}[htbp]
\centering
\caption{Marginal willingness-to-accept (MWTA) for one month of vaccine delay, by subgroup}
\label{tab:mwta_subgroups}
\begin{threeparttable}
\begin{tabular}{lccc}
\toprule
Group & $n$ & MWTA (RMB/month) & 95\% CI \\
\midrule
\multicolumn{4}{l}{\textit{Overall benchmark}} \\
Overall sample & 1027 & 44.8 & [28.9,\; 60.7] \\
\addlinespace

\multicolumn{4}{l}{\textit{Demographics}} \\
Female & 545 & 55.7 & [32.9,\; 78.5] \\
Male & 482 & 33.1 & [10.9,\; 55.3] \\
Urban & 472 & 57.7 & [32.6,\; 82.8] \\
Rural & 555 & 34.7 & [14.1,\; 55.2] \\
Low education & 765 & 51.2 & [32.1,\; 70.3] \\
High education & 262 & 28.1 & [-0.6,\; 56.8] \\
Younger ($<55$ years) & 768 & 49.8 & [30.9,\; 68.9] \\
Older ($\geq 55$ years) & 259 & 30.8 & [1.7,\; 59.8] \\
\addlinespace

\multicolumn{4}{l}{\textit{Health behaviour}} \\
Non-smoker & 625 & 38.4 & [18.4,\; 58.4] \\
Smoker & 402 & 55.0 & [28.6,\; 81.3] \\
Low physical activity & 297 & 80.1 & [44.6,\; 115.6] \\
Mid physical activity & 439 & 27.6 & [5.1,\; 50.1] \\
High physical activity & 291 & 38.9 & [10.1,\; 67.7] \\
\addlinespace

\multicolumn{4}{l}{\textit{Institutional trust}} \\
Low institutional trust & 487 & 62.3 & [37.1,\; 87.5] \\
High institutional trust & 446 & 28.7 & [6.5,\; 50.8] \\
\addlinespace

\multicolumn{4}{l}{\textit{Medical condition and self-rated health}} \\
No medical condition & 908 & 43.4 & [26.7,\; 60.2] \\
With medical condition & 119 & 57.5 & [6.4,\; 108.6] \\
Good self-rated health & 772 & 44.8 & [26.3,\; 63.3] \\
Poor self-rated health & 255 & 44.9 & [13.6,\; 76.3] \\
\bottomrule
\end{tabular}
\begin{tablenotes}
\footnotesize
\item Notes: MWTA values are derived from multinomial logit models with vaccine delay and cash incentives entered linearly; MWTA for a one-month increase in delay is defined as $-\beta_{\textit{delay}}/\beta_{\textit{cash}}$. Confidence intervals are based on the delta method with standard errors clustered at the respondent level. Positive MWTA indicates the monetary compensation required to offset the disutility of an additional month of vaccine delay. Values are rounded to one decimal place.
\end{tablenotes}
\end{threeparttable}
\end{table}

\begin{figure}[htbp]
\centering
\includegraphics[width=\textwidth]{Subgroup differences in time sensitivity.png}
\caption{Subgroup heterogeneity in marginal willingness-to-accept (MWTA) for one month of vaccination delay. Points denote subgroup MWTA estimates and horizontal bars indicate 95\% confidence intervals. The vertical dotted line marks the overall MWTA benchmark (44.8 RMB/month).}
\label{fig:mwta_forest}
\end{figure}



\clearpage

\subsection{Policy-bundle simulations}
\label{subsec:essay2_policy_bundles}

To translate the estimated preference parameters into quantities directly relevant for programme planning, I simulate predicted vaccine uptake under counterfactual policy bundles that jointly vary vaccination delay and cash incentives. Simulations are computed from the main MNL model and evaluated at a reference vaccine profile holding efficacy at 70\%, side effects at mild, and origin as domestic. For each bundle, predicted uptake is defined as the model-implied probability of choosing vaccination (relative to opting out) when only delay and incentives are varied.

Table~\ref{tab:policy_bundles} reports predicted uptake for delays spanning 0--6 months and incentive amounts ranging from 0 to 800 RMB. In the absence of incentives, predicted uptake declines monotonically with waiting time, falling from 69\% at zero delay to 62\% at a six-month delay. Monetary incentives partially offset these losses at every delay level. A modest incentive of 50 RMB raises uptake by about 1 percentage point. A 200 RMB incentive increases uptake by roughly 4--5 percentage points, while an 800 RMB incentive increases uptake by 15--17 percentage points. Notably, incentives can sustain relatively high uptake even when delays are long: at a six-month delay, predicted uptake increases from 62\% with no incentive to 79\% with an 800 RMB payment.

\begin{table}[htbp]
\centering
\caption{Predicted vaccine uptake under policy bundles varying delay (0--6 months) and cash incentives (0--800 RMB)}
\label{tab:policy_bundles}
\begin{threeparttable}
\begin{tabular}{lcccc}
\toprule
Delay (months) & 0 RMB & 50 RMB & 200 RMB & 800 RMB \\
\midrule
\multicolumn{5}{l}{\textit{Panel A. Predicted uptake (\%)}} \\
0 & 69 & 70 & 74 & 84 \\
1 & 68 & 69 & 73 & 83 \\
3 & 66 & 67 & 70 & 82 \\
6 & 62 & 63 & 67 & 79 \\
\addlinespace
\multicolumn{5}{l}{\textit{Panel B. Percentage-point change relative to 0 RMB}} \\
0 & --- & +1 & +5 & +15 \\
1 & --- & +1 & +5 & +15 \\
3 & --- & +1 & +4 & +16 \\
6 & --- & +1 & +5 & +17 \\
\bottomrule
\end{tabular}
\begin{tablenotes}
\footnotesize
\item Notes: Predicted uptake probabilities computed from the main MNL model holding efficacy at 70\%, side effects at mild, and origin as domestic. Percentage-point changes indicate absolute increases in uptake relative to the 0 RMB incentive condition at each delay level.
\end{tablenotes}
\end{threeparttable}
\end{table}

While Table~\ref{tab:policy_bundles} reports discrete policy bundles, Figure~\ref{fig:iso_acceptance} generalizes these results by mapping predicted uptake over a continuous grid of delay and incentive combinations. The heatmap visualizes the full uptake surface, with warmer colors indicating higher predicted uptake. Superimposed iso-acceptance contours identify the combinations of delay and compensation required to sustain 60\% and 80\% coverage thresholds.

These contours form a practical policy frontier. For any anticipated delay, the figure indicates the minimum incentive required to maintain a target uptake level. Conversely, for a given incentive budget, it identifies the maximum tolerable delay consistent with the desired coverage rate. This visualization operationalizes the estimated delay--cash trade-off into a decision-support tool for comparing investments in faster access with financial compensation strategies.

\begin{figure}[htbp]
\centering
\includegraphics[width=0.9\textwidth]{Iso-acceptance heatmap.png}
\caption{Iso-acceptance heatmap showing predicted uptake under delay--cash combinations. Contours represent 60\% and 80\% predicted uptake thresholds. Baseline profile: 70\% efficacy, mild side effects, domestic origin.}
\label{fig:iso_acceptance}
\end{figure}
\clearpage
\subsection{Validity checks and robustness}
\label{sec:essay2_validity}

Because Essays~1 and~2 rely on the same discrete choice experiment (DCE) backbone, several internal validity diagnostics have already been documented in Chapter~\ref{ch:supp_essay1} and summarized in Section~\ref{sec:validity_robustness}. These include dominance tests, response-pattern consistency checks, and assessments of attribute non-attendance. Together, those diagnostics indicate that respondents understood the choice tasks and engaged in meaningful attribute trade-offs.

Building on that foundation, the robustness analysis in this chapter focuses specifically on the stability of the monetized delay estimates. Because MWTA is a nonlinear ratio of estimated coefficients, statistical inference requires particular care. I therefore implement a clustered bootstrap procedure (1,000 replications) that resamples at the respondent level to preserve the panel structure of repeated choices. I also examine sensitivity to influential observations by re-estimating key trade-off parameters in trimmed samples and assess robustness to alternative model specifications and coding choices.

For transparency and readability, detailed robustness tables—including clustered bootstrap standard errors and percentile intervals (Table~9), trimmed-sample comparisons (Table~10), and additional specification checks (Table~11)—are reported in Chapter~\ref{ch:supp_essay2}, \emph{Supplementary Results for Essay 2}. The substantive conclusions regarding the monetary valuation of delay and subgroup heterogeneity remain stable across these alternative inference approaches and sample restrictions.

\clearpage

\section{Discussion and policy implications}
\label{sec:essay2_discussion}

This chapter quantifies the monetary value individuals attach to timely vaccination access. By converting delay sensitivity into marginal willingness-to-accept (MWTA), the analysis provides a direct monetary measure of the behavioral cost of waiting. The estimated shadow price of delay indicates that even short postponements carry economically meaningful welfare losses, reinforcing the interpretation of vaccination timing as an intertemporal health investment decision.

From a behavioral perspective, these findings align with evidence that individuals overweight near-term costs relative to delayed benefits, particularly in health domains involving uncertainty and risk \cite{Chapman2005,Frederick2002}. In vaccination contexts, delays effectively shift protection further into the future while exposure risk remains immediate. The MWTA estimates therefore capture not merely inconvenience but the perceived welfare loss associated with postponed risk reduction.

The heterogeneity patterns further underscore that the cost of delay is not uniform across the population. Groups exhibiting lower institutional trust or lower preventive engagement display substantially higher MWTA, suggesting that delayed access may exacerbate inequities in protection. This is consistent with research showing that trust and perceived system reliability shape vaccination behavior and policy responsiveness \cite{Betsch2018,Larson2018}. Uniform delay schedules may therefore impose disproportionate welfare burdens on specific subpopulations, even when vaccines are provided free of charge.

The policy implications are twofold. First, the results highlight the value of operational investments that reduce waiting time. Queue management, appointment expansion, and supply stabilization may yield behavioral returns that exceed their administrative costs when the shadow price of waiting is high. Similar operational improvements have been shown to meaningfully affect preventive health behaviors in other contexts \cite{Milkman2011,Volpp2008}. 

Second, when delays are unavoidable, financial incentives can partially compensate for postponed protection. The iso-acceptance curves provide a decision-support framework for determining the compensation required to sustain target coverage levels. Rather than treating incentives as ad hoc tools, the MWTA-based approach allows policymakers to calibrate payments according to empirically estimated delay costs, consistent with principles of cost-effectiveness and program budgeting \cite{Drummond2015}. 

Importantly, the MWTA framework does not imply that financial incentives are universally optimal. Incentive effectiveness depends on baseline trust, perceived legitimacy, and implementation design \cite{Gneezy2011}. In some settings, accelerating access may be more efficient or equitable than compensating for delay. The trade-off between speed-up investments and monetary incentives therefore depends on relative costs, budget constraints, and distributional objectives—a question taken up explicitly in Essay~3.

Finally, these findings contribute to the broader literature on time in health economics. While economic evaluations typically discount future health outcomes at a uniform rate, this chapter demonstrates that individuals assign context-specific monetary values to short-term delays in preventive care. Integrating such behavioral valuations into program design may improve alignment between formal cost-effectiveness criteria and real-world uptake behavior.

Together, the results position timeliness not merely as a logistical parameter but as an economically measurable policy lever. By quantifying the shadow price of waiting and mapping it into actionable trade-offs, this chapter provides the empirical foundation for the equity-sensitive policy rules developed in the subsequent essay.


\section{Limitations}
\label{sec:essay2_limitations}

This chapter has several limitations. First, the estimates are based on stated preferences elicited in a discrete choice experiment, and therefore may be affected by hypothetical bias if respondents’ stated trade-offs differ from their real-world behavior under actual vaccination opportunities. Although the design incorporates internal validity diagnostics and the main results are robust to alternative inference procedures, the MWTA values should be interpreted as behavioral valuations under controlled scenarios rather than revealed market prices.

Second, the experiment represents delay in discrete monthly increments over a 0--6 month horizon. This framing is intentionally aligned with operational planning, but it may not capture valuation at very short horizons (e.g., days or weeks) or nonlinearities beyond the range studied. Relatedly, MWTA is a reduced-form trade-off that summarizes multiple psychological channels (present bias, uncertainty discounting, perceived risk, and affective responses) and does not identify structural $\beta$--$\delta$ parameters of time preference.

Third, the analysis focuses on individual-level preferences and abstracts from epidemiological dynamics and social interactions. In practice, speeding up access may generate additional external benefits by reducing transmission and accelerating population immunity, while delayed uptake may impose social costs not captured by individual utility alone. As a result, the MWTA estimates should be viewed as behavioral inputs for programme design rather than sufficient statistics for welfare analysis.

Finally, context matters. The study is conducted among adults in Wuhan under a high-salience epidemic context, which may shape how respondents perceive waiting time, risk, and trust. While the conceptual mechanisms are general, the numerical MWTA magnitudes may not transfer one-for-one to settings with different baseline acceptance, risk perceptions, institutional environments, or health-system capacity. Replication in other locations and in non-epidemic vaccination contexts would help assess external validity and generalizability.

\section{Conclusion}
\label{sec:essay2_conclusion}

This chapter monetizes the behavioral cost of delayed vaccination access by estimating marginal willingness-to-accept (MWTA) for waiting time in a discrete choice experiment. The results show that timeliness carries a substantial shadow price: even short delays impose economically meaningful welfare losses in monetary terms. The chapter further documents systematic heterogeneity in MWTA across demographic characteristics, preventive behaviors, medical vulnerability, and institutional trust, indicating that delays are not experienced uniformly across the population.

By translating preference parameters into predicted uptake under delay--cash bundles and by constructing iso-acceptance frontiers, the chapter provides a practical decision-support framework for comparing ``speed-up'' investments to compensation policies when delays are unavoidable. These monetized valuations set up the next chapter (Essay~3), which integrates delay valuation with discounting heterogeneity to derive equity-sensitive policy rules for prioritizing acceleration versus compensation under budget and distributional constraints.
\clearpage

\chapter{Essay 3: Integrating Behavioral Discounting and Monetary Valuation for Equity-Sensitive Vaccination Policy Design}
\label{sec:essay3}

\section{Introduction}
\vspace*{-\baselineskip}
\noindent
Vaccination policy design requires more than understanding whether individuals accept or refuse immunization. It requires understanding how people value \emph{when} protection arrives, and how those valuations should shape resource allocation under operational constraints. Essays~1 and~2 demonstrated that vaccination timing reflects systematic behavioral heterogeneity: individuals exhibit measurable delay sensitivity consistent with intertemporal discounting, and they attach economically meaningful monetary value to reducing waiting time. These findings raise a natural policy question: how should vaccination programmes allocate scarce resources when delay costs are heterogeneous across the population?

Traditional economic evaluation frameworks treat time primarily through uniform discounting of future health outcomes \cite{Drummond2015}. However, behavioral research shows that individuals often display present bias and non-exponential discounting, leading to disproportionately high weight on near-term costs relative to delayed benefits \cite{Laibson1997,ODonoghueRabin2001}. In vaccination contexts, such behavioral discounting implies that short delays may generate larger welfare losses than implied by standard exponential models. Moreover, the monetary valuation results in Essay~2 show that these delay costs vary systematically with institutional trust, preventive engagement, and demographic characteristics, suggesting that uniform delay schedules may generate unequal welfare burdens.

At the same time, policymakers face concrete trade-offs. Reducing waiting time requires operational investments in supply chains, staffing, and appointment systems. When acceleration is infeasible or costly, financial incentives may compensate individuals for postponed protection. Field evidence indicates that monetary incentives can increase vaccination uptake under certain conditions \cite{CamposMercade2021,Thirumurthy2022}, yet incentive effectiveness depends on context and behavioral responsiveness \cite{Gneezy2011}. The key policy challenge, therefore, is to determine when accelerating access is more efficient than compensating delay—and how these choices should vary across heterogeneous populations.

This chapter integrates behavioral discounting and monetary valuation to derive equity-sensitive policy rules for vaccination timing. By combining reduced-form discount parameters from Essay~1 with marginal willingness-to-accept (MWTA) estimates from Essay~2, I construct transparent decision conditions under which speed-up investments dominate compensation policies, and vice versa, under explicit budget constraints. The framework further allows subgroup-specific delay costs and welfare weights to inform targeted prioritization strategies, consistent with distributional objectives in health policy \cite{Cookson2017,Asaria2016}.

The contribution of this chapter is threefold. First, it provides a unified behavioral-economic framework linking intertemporal preference heterogeneity to operational policy design. Second, it derives implementable decision rules comparing acceleration and compensation as alternative instruments for managing vaccination timeliness. Third, it incorporates equity considerations by explicitly modeling how heterogeneous delay valuation interacts with resource allocation. Together, these results move from measuring delay preferences to prescribing how vaccination programmes should allocate resources when timeliness itself is a scarce policy lever.


\section{Conceptual Integration Framework}
\label{sec:essay3_framework}

This chapter integrates two empirical objects estimated in Essays~1--2—behavioral discounting parameters and the monetary valuation of waiting—into a unified welfare-consistent framework for policy design. The goal is to express \emph{months of delay} and \emph{RMB incentives} on a common scale so that operational speed-ups and compensation schemes can be compared under explicit budget and equity objectives.

\subsection{Linking behavioral discounting to delay sensitivity ($\kappa$, $\delta$)}
\label{subsec:essay3_link_kappa_delta}

Let individual $i$ in group $g$ evaluate a vaccination alternative with delay $t$ (months) and other attributes $x$ (e.g., efficacy, side effects, origin). In Essays~1--2, the choice model is written in reduced form as
\begin{equation}
U_{ijg} = \beta_{t,g}\, t_{ij} + \beta_{m,g}\, m_{ij} + \beta_{x,g}'x_{ij} + \varepsilon_{ijg},
\end{equation}
where $m$ denotes monetary incentives (scaled in RMB or 100 RMB units). Essay~1 interprets the estimated delay sensitivity $\beta_{t,g}$ as a reduced-form representation of intertemporal valuation and maps it into behavioral discounting parameters (e.g., present-bias and long-run discounting) under parsimonious functional forms \cite{Laibson1997,ODonoghueRabin2001,Frederick2002}. Conceptually, stronger present bias (lower $\kappa$) or heavier long-run discounting (lower $\delta$) implies a steeper marginal utility penalty for postponing protection, which appears empirically as a more negative $\beta_{t,g}$.

\subsection{Linking MWTA to the same preference primitives}
\label{subsec:essay3_link_mwta}

Essay~2 converts delay sensitivity into monetary units using marginal willingness-to-accept (MWTA). For group $g$,
\begin{equation}
\text{MWTA}_g \equiv -\frac{\beta_{t,g}}{\beta_{m,g}},
\end{equation}
which is the marginal rate of substitution between waiting time and monetary compensation. MWTA therefore provides a money-metric ``shadow price of waiting'': the RMB required to offset the utility loss from one additional month of delay. Importantly, MWTA must be computed using \emph{unstandardized} coefficients so that it preserves the RMB-per-month interpretation.

Because both $\kappa/\delta$ (Essay~1) and MWTA (Essay~2) are functions of the same underlying delay sensitivity $\beta_{t,g}$, they provide complementary views of the same behavioral construct:
\begin{itemize}[leftmargin=*]
\item $\kappa,\delta$ summarize the \emph{shape} of intertemporal valuation (behavioral discounting);
\item $\text{MWTA}_g$ summarizes the \emph{money-metric magnitude} of the marginal cost of delay.
\end{itemize}
This dual representation is useful for policy design: discounting parameters motivate why short delays can be disproportionately costly, while MWTA provides the operational exchange rate needed for budgeting and incentive calibration.

\subsection{Welfare interpretation and policy-relevant trade-offs}
\label{subsec:essay3_welfare}

Define a policy that (i) reduces expected delay by $\Delta t_{g}$ months for group $g$ (a speed-up) and/or (ii) offers monetary compensation $\Delta m_{g}$ (an incentive). Using the estimated preference parameters, the compensating variation in monetary units implied by a speed-up is
\begin{equation}
\Delta m^{\ast}_{g} \approx \text{MWTA}_g \cdot \Delta t_{g},
\end{equation}
i.e., the incentive amount that would generate the same utility change as saving $\Delta t_g$ months. This equivalence mapping allows a direct comparison of two instruments on a common scale:
\begin{enumerate}
\item \textbf{Speed-up instrument:} incurs an operational cost $C^{S}(\Delta t_g)$ and yields welfare benefits valued at $\text{MWTA}_g \cdot \Delta t_g$;
\item \textbf{Compensation instrument:} incurs a fiscal cost $C^{M}(\Delta m_g)$ and yields welfare benefits proportional to the incentive-induced utility gain.
\end{enumerate}

To incorporate equity objectives, let $\omega_g$ denote a social welfare weight for group $g$ (e.g., prioritizing disadvantaged or high-burden groups). The policy objective can be written as a weighted welfare aggregation,
\begin{equation}
W = \sum_{g} \omega_g \, \mathbb{E}\big[\Delta U_g(\Delta t_g,\Delta m_g)\big],
\end{equation}
subject to budget and capacity constraints. This formulation aligns with distributional approaches to economic evaluation, where policy choices are evaluated not only by average gains but also by who gains and who bears burdens \cite{Asaria2016,Cookson2017}.

This integrated framework sets up the central decision problem of Essay~3: given heterogeneous $\text{MWTA}_g$ and discounting profiles $(\kappa_g,\delta_g)$, when is it welfare- and equity-improving to allocate resources to \emph{reduce waiting time} versus \emph{compensate waiting}, and how should these instruments be targeted across groups under constraints?

\clearpage

\section{Deriving policy rules: when to speed up access versus compensate delay}
\label{sec:essay3_rules}

Vaccination programmes typically face two practical levers when appointments and supply constraints generate waiting time. The first is to \textbf{speed up access} by investing in operations (e.g., staffing, appointment capacity, distribution logistics, outreach that reduces scheduling frictions). The second is to \textbf{compensate unavoidable delay} through monetary incentives that help maintain uptake when waiting cannot be eliminated. This section translates the empirical results from Essays~1--2 into transparent policy rules for choosing between these levers and for targeting them across groups.


\paragraph{Technical details.}
For completeness, Chapter~\ref{ch:supp_essay3} provides (i) the formal welfare maximization setup, (ii) the algebraic steps linking the MWTA mapping to the policy conditions, and (iii) alternative formulations under uptake-target constraints.


\subsection{Rule 1: Compare the cost of saving time to the ``shadow price'' of waiting}
\label{subsec:essay3_rule1}

Essay~2 estimates a money-metric ``shadow price of waiting,'' $\text{MWTA}_g$, which can be interpreted as:

\begin{quote}
\emph{For group $g$, MWTA is the RMB needed to offset the perceived welfare loss from one additional month of vaccination delay.}
\end{quote}

This creates a simple policy benchmark. Suppose a planner considers an operational investment that reduces expected waiting time for group $g$ by one month. Let the programme cost of achieving that one-month reduction (per person) be denoted $\text{Cost}^{S}_g$.

\textbf{Speed-up is preferred} when saving a month of waiting costs less than the group’s shadow price of waiting:
\begin{equation}
\textbf{Speed up access for group $g$ if:}\quad \text{Cost}^{S}_g \le \text{MWTA}_g.
\label{eq:policy_rule_speed}
\end{equation}

\textbf{Compensation is preferred} when speeding up is more expensive than what that group is effectively willing to accept as compensation for waiting:
\begin{equation}
\textbf{Compensate delay for group $g$ if:}\quad \text{Cost}^{S}_g > \text{MWTA}_g.
\label{eq:policy_rule_comp}
\end{equation}

This rule is deliberately intuitive: it compares \emph{what it costs the programme to remove one month of delay} to \emph{how costly that month of delay is to people in RMB terms}. When the operational cost is low relative to MWTA, investing in speed produces high welfare returns. When the operational cost is high, incentives may be the more feasible tool for maintaining coverage.

\subsection{Rule 2: Use equity priorities to decide who should receive speed-ups first}
\label{subsec:essay3_rule2}

Essays~1--2 show that the cost of delay is heterogeneous across subgroups. Some groups—especially those with low institutional trust or low preventive engagement—place substantially higher monetary value on timeliness. This implies that ``one month saved'' is not equally valuable for everyone. A uniform rollout schedule can therefore impose disproportionate welfare burdens on groups for whom waiting is most costly.

To incorporate equity objectives, policymakers can apply explicit priority weights, $\omega_g$, reflecting distributional goals (e.g., prioritizing disadvantaged groups or groups experiencing higher welfare losses from delay). Under this equity-sensitive perspective, the speed-up rule becomes:

\begin{equation}
\textbf{Prioritize speed-up for group $g$ if:}\quad \text{Cost}^{S}_g \le \omega_g \cdot \text{MWTA}_g.
\label{eq:policy_rule_equity}
\end{equation}

This adjustment is policy-relevant: even if a speed-up is not cost-effective under a purely average-welfare criterion, it may still be justified for groups assigned higher equity priority.

\subsection{Rule 3: Allocate scarce resources where the welfare return per dollar is highest}
\label{subsec:essay3_budget}

When budgets are limited, the allocation problem becomes one of \emph{ranking} interventions. A simple way to operationalize this is to compare welfare return per RMB of programme spending.

For speed-ups, a useful summary statistic is:
\begin{equation}
\text{Welfare return per RMB} \propto \frac{\omega_g \cdot \text{MWTA}_g}{\text{Cost}^{S}_g}.
\label{eq:bang_for_buck_policy}
\end{equation}

This ratio yields an implementable targeting rule:

\begin{quote}
\textbf{Allocate speed-up investments first to groups with higher equity-weighted MWTA and lower operational costs of reducing delay.}
\end{quote}

Where speed-ups are infeasible or expensive, incentives can be used as a second-best tool to protect uptake. In the simulation section, I translate these rules into uptake-based decision support: for each delay horizon, I quantify the incentive required to sustain target coverage thresholds (e.g., 60\% or 80\%) and compare those incentive costs to the operational cost of shortening delays by the same amount.



\subsection{Summary: a decision checklist for programme planners}
\label{subsec:essay3_checklist}

Taken together, the policy rules imply a simple checklist:

\begin{enumerate}
\item \textbf{Estimate the shadow price of waiting (MWTA) for relevant groups.}
\item \textbf{Estimate operational cost per month saved} (capacity expansion, logistics, queue management).
\item \textbf{Choose the cheaper lever} on a common RMB-per-month scale:
\begin{itemize}
\item Speed-up if cost per month saved $\le$ MWTA (or $\omega_g \cdot$ MWTA under equity priorities);
\item Otherwise, use incentives to compensate unavoidable delay.
\end{itemize}
\item \textbf{Under a fixed budget, target first where the welfare return per RMB is highest.}
\end{enumerate}

These rules operationalize the empirical heterogeneity documented in Essays~1--2 and provide a transparent bridge from behavioral preference estimates to actionable design choices in vaccination rollouts.


\section{Equity-Sensitive Targeting}
\label{sec:essay3_equity}

The policy rules above identify when accelerating access is more cost-effective than compensating delay on average. However, vaccination programmes rarely aim to maximise average welfare alone. Because the welfare burden of waiting is heterogeneous (Essays~1--2), uniform rollout schedules can impose disproportionate losses on particular subpopulations. This section therefore introduces an equity-sensitive targeting approach that (i) allows subgroup welfare weights to reflect distributional priorities and (ii) makes the resulting trade-offs transparent for planners allocating scarce operational and fiscal resources.

Formally, let $g$ index policy-relevant groups and let $\omega_g \ge 0$ denote a social welfare weight that upweights gains accruing to priority populations (e.g., groups facing disadvantage, higher vulnerability, or larger welfare losses from delay). A simple and implementable targeting rule is to allocate speed-up resources proportional to the \emph{equity-weighted welfare return} from reducing delay:
\begin{equation}
\text{Allocation share}_g \;\propto\; \omega_g \times \text{MWTA}_g \times N_g,
\label{eq:equity_allocation_rule}
\end{equation}
where $\text{MWTA}_g$ is the subgroup-specific monetary valuation of a one-month delay (Essay~2) and $N_g$ is the subgroup size. This rule has an intuitive interpretation: groups receive more operational acceleration when an additional month of waiting is more costly to them in money-metric terms (higher MWTA), when they represent more people (larger $N_g$), and when the planner assigns greater equity priority (higher $\omega_g$).

Figure~\ref{fig:equity_allocation} illustrates how implied budget shares change under alternative equity-weighting schemes using the institutional-trust split as an example. Under equal weights ($\omega_g=1$), allocation reflects differences in MWTA and group size alone. When the planner assigns a higher equity weight to the low-trust group (e.g., $\omega_{\text{low trust}}=2$), a larger share of the speed-up budget is directed toward that group, even holding total programme spending fixed. This figure highlights a central message for policy design: equity objectives are not distributionally neutral. Explicit weighting choices can materially reshape who benefits from operational acceleration, and they provide a transparent bridge between normative priorities and implementable targeting decisions. In the next section, I evaluate how such targeting rules affect both coverage outcomes and equity-weighted welfare under realistic budget constraints.


\begin{figure}[htbp]
\centering
\includegraphics[width=0.8\textwidth]{essay3_equity_budget_allocation.png}
\caption{Budget allocation across groups under alternative equity weights.
Illustrative allocation shares are proportional to $\omega_g \times \text{MWTA}_g \times N_g$,
where $\omega_g$ denotes subgroup welfare weights, $\text{MWTA}_g$ is the monetary valuation of a one-month delay,
and $N_g$ is subgroup size. Under equal weights, allocation reflects delay valuation and group size. Under equity
weights (e.g., doubling the weight for low-trust individuals), a larger share of operational speed-up resources is
directed toward the prioritized group.}
\label{fig:equity_allocation}
\end{figure}

\clearpage

\subsection{Subgroup welfare weights: making equity objectives explicit}
\label{subsec:essay3_weights}

Let $g$ index subgroups defined by observable characteristics (e.g., institutional trust, preventive engagement, medical vulnerability, urban/rural residence). The central equity question is: \emph{should a month of delay reduction be valued equally across groups?} If not, policy design should reflect two distinct reasons why timeliness may deserve higher priority for some groups:

\begin{enumerate}
\item \textbf{Higher behavioral cost of waiting (preference-based equity).} Some groups exhibit higher MWTA, implying larger welfare losses from delay even holding operational circumstances fixed. In this sense, delay is \emph{more burdensome} for them.
\item \textbf{Social priority (normative equity).} Even if two groups value delay similarly, policymakers may choose to prioritise disadvantaged or high-need groups based on fairness principles or distributional goals.
\end{enumerate}

To incorporate these objectives transparently, define a social welfare weight $\omega_g \ge 0$ for each subgroup. Larger $\omega_g$ indicates that welfare improvements accruing to group $g$ receive greater policy priority.

In practice, $\omega_g$ can be operationalised in several policy-relevant ways:
\begin{itemize}[leftmargin=*]
\item \textbf{Equal weights} ($\omega_g=1$ for all $g$): benchmark ``utilitarian'' design.
\item \textbf{Disadvantage weights}: higher $\omega_g$ for groups facing structural disadvantage or poorer baseline health.
\item \textbf{Vulnerability weights}: higher $\omega_g$ for medically vulnerable groups or those facing higher consequences from infection.
\item \textbf{Access-friction weights}: higher $\omega_g$ for groups facing higher waiting burdens that the system can plausibly reduce.
\end{itemize}

This dissertation does not impose a single ethical rule; rather, it provides a framework that allows committees and policymakers to \emph{state equity priorities explicitly} and to observe how those priorities change targeting decisions.

\subsection{Distributional trade-offs: who gains from speed-ups versus compensation}
\label{subsec:essay3_tradeoffs}

Equity-sensitive targeting requires distinguishing the distributional implications of the two instruments:

\paragraph{Speed-ups.}
Speed-ups reduce waiting time directly. They tend to be most equity-relevant when (i) the delay burden is high for certain groups and (ii) operational changes can be targeted (e.g., reserved appointment slots, mobile clinics, extended-hours staffing, priority scheduling). In money-metric terms, the welfare gain from reducing delay by $\Delta t_g$ months for group $g$ is approximately:
\begin{equation}
\Delta \widetilde{W}^{S}_g \approx \omega_g \cdot N_g \cdot \text{MWTA}_g \cdot \Delta t_g.
\label{eq:equity_speedup_gain}
\end{equation}
Thus, speed-ups create larger equity-weighted gains when a group has higher MWTA (delay is more costly) and/or higher social priority weight.

\paragraph{Compensation (incentives).}
Compensation mitigates the welfare loss from delay without changing the waiting time itself. Incentives may be easier to deploy quickly and can be targeted by eligibility criteria, but they also raise distributional concerns: they require fiscal outlays, may differentially affect groups’ behavior, and may be viewed as less legitimate in low-trust settings. In a money-metric accounting, transfers to a subgroup are valued at:
\begin{equation}
\Delta \widetilde{W}^{M}_g \approx \omega_g \cdot N_g \cdot \Delta m_g,
\label{eq:equity_incentive_gain}
\end{equation}
while their uptake effects are evaluated empirically in the simulation section by mapping $(t,m)$ bundles into predicted uptake.

\paragraph{The equity trade-off.}
Because speed-ups change the \emph{timing} of protection while incentives change the \emph{immediacy} of benefits (through cash), the two instruments can have different equity consequences even when they appear similar on average. Two distributional tensions are especially relevant:

\begin{enumerate}
\item \textbf{Efficiency vs equity targeting.} Groups with the highest MWTA (largest delay burdens) are not always the groups for whom speed-ups are cheapest to deliver. Equity-sensitive policy may therefore justify allocating resources toward groups where operational speed is more expensive, if their equity-weighted welfare loss from delay is high.
\item \textbf{Uniform versus targeted programmes.} Uniform incentives or uniform acceleration can be politically simple but may be inefficient and inequitable if delay burdens vary. Targeting improves allocative precision but raises administrative and fairness considerations, including eligibility verification and potential stigma.
\end{enumerate}

\subsection{An implementable targeting rule under equity objectives}
\label{subsec:essay3_target_rule}

Combining the above, a practical targeting rule is to prioritise speed-ups in groups with the largest equity-weighted welfare gains per RMB spent. Let $\text{Cost}^{S}_g$ denote the programme cost per person of reducing delay by one month for group $g$. Then an equity-sensitive ``return on spending'' index is:
\begin{equation}
\text{Equity-adjusted return}_g \propto \frac{\omega_g \cdot \text{MWTA}_g}{\text{Cost}^{S}_g}.
\label{eq:equity_rank}
\end{equation}
Under a fixed budget, planners should allocate operational speed-up investments first to groups with the highest values of \eqref{eq:equity_rank}, subject to feasibility constraints (e.g., limits on how much delay can be reduced for any group). When $\text{Cost}^{S}_g$ exceeds $\omega_g \cdot \text{MWTA}_g$, compensation becomes the preferred instrument for that group, consistent with the decision rule in Section~\ref{sec:essay3_rules}.

\subsection{How this section connects to simulations}
\label{subsec:essay3_equity_to_sim}

The remainder of the chapter applies these principles in policy simulations. I evaluate alternative targeting strategies (uniform versus subgroup-targeted) under plausible budget envelopes and coverage goals. For each strategy, I compare: (i) expected uptake under delay--incentive bundles; (ii) implied programme cost; and (iii) equity-weighted welfare outcomes under alternative choices of subgroup weights $\omega_g$. This provides decision support for policymakers who must trade off rapid scale-up, fiscal feasibility, and distributional fairness in future epidemic vaccination campaigns.


\section{Policy Simulations}
\label{sec:essay3_sims}

This section operationalizes the policy rules derived above by translating behavioral parameters into implementable programme decisions. Using the estimated preference coefficients and predicted uptake functions from Essays~1--2, I simulate coverage and cost outcomes under alternative combinations of delay reduction (speed-ups) and monetary compensation (incentives). The objective is to identify which instrument is more cost-effective under realistic operational conditions and uptake targets.

Unless otherwise stated, simulations evaluate uptake at a common reference vaccine profile (70\% efficacy, mild side effects, domestic origin), so that policy comparisons isolate the effects of timeliness and incentives. Delays are evaluated over an operationally relevant range (0--6 months), and incentives over a policy-relevant range (0--800 RMB), consistent with the experimental design.

\subsection{Least-cost policy design under uptake targets}
\label{subsec:essay3_least_cost}

Many vaccination campaigns operate under coverage thresholds (e.g., reaching 80\% uptake before a surge period). To mirror this constraint, I compute the minimum-cost combination of delay reduction and incentives required to maintain a target uptake level.

Figure~\ref{fig:tradeoff} visualizes the break-even condition between accelerating access and compensating delay for maintaining 80\% uptake. The upward-sloping line represents the operational cost of eliminating delay, assuming an illustrative unit cost of 50 RMB per month saved per person. The second line shows the incentive required at each delay level to sustain 80\% predicted uptake, derived from the estimated choice model.

\begin{figure}[htbp]
\centering
\includegraphics[width=0.85\textwidth]{essay3_policy_tradeoff_figure (1).png}
\caption{Policy trade-off between accelerating access and compensating delay.
The upward-sloping line represents the operational cost of reducing delay to zero
(assumed 50 RMB per month saved per person). The second line shows the
monetary incentive required to maintain 80\% uptake at each delay level,
based on model predictions. The intersection identifies the break-even delay
at which the two strategies cost the same per person.}
\label{fig:tradeoff}
\end{figure}

The intersection of the two curves identifies the delay level at which the per-person cost of speed-up equals the incentive required to offset delay. To the left of this point, operational acceleration is less costly and therefore preferred. To the right, monetary compensation becomes the more cost-effective instrument. This graphical representation operationalizes the analytical condition derived earlier: speed-up is preferred when the marginal operational cost per month saved is lower than the behavioral valuation of delay.

\paragraph{Sensitivity to operational cost assumptions.}
The break-even comparison depends critically on the relative magnitude of operational costs ($c$) versus behavioral valuation ($MWTA$). In the baseline ($c=50$, $MWTA \approx 45$), the two instruments are roughly cost-equivalent, with speed-ups slightly more expensive on average but preferred for groups with high delay aversion. As shown in the detailed sensitivity analysis in Appendix~\ref{sec:supp_essay3_sensitivity} (Table~\ref{tab:supp_sensitivity}), a variation of $\pm 20$--50\% in operational costs shifts the dominant strategy. If operational costs are 20--50\% lower (e.g., $c=25-40$ RMB/month), speed-up becomes the dominant strategy for nearly all groups, as the cost of saving time is strictly below the compensation required to offset it. Conversely, if operational acceleration is more costly (e.g., $c=60-75$ RMB/month), the break-even frontier shifts such that monetary compensation becomes the uniformly preferred second-best instrument, except potentially for sub-populations with extremely high delay aversion (e.g., $MWTA > 75$). This sensitivity underscores that determining the optimal policy mix requires local evidence on both delivery costs and behavioral preferences.

\clearpage

\subsection{Uniform versus targeted implementation}
\label{subsec:essay3_uniform_targeted}

While Figure~\ref{fig:tradeoff} evaluates average effects, heterogeneity in delay valuation implies that the break-even condition varies across subgroups. I therefore simulate three implementation strategies:

\begin{enumerate}
\item \textbf{Uniform compensation:} identical incentives offered to all individuals facing delay.
\item \textbf{Uniform speed-up:} operational investment reducing delay equally across groups.
\item \textbf{Equity-sensitive targeting:} prioritizing delay reduction for subgroups with higher MWTA or higher social welfare weights, while compensating others.
\end{enumerate}

For each strategy, I compute predicted uptake, programme cost, and (when equity weights are applied) equity-adjusted welfare gains. This approach allows planners to compare not only which instrument is cheaper on average, but also how targeting alters both coverage and distributional outcomes.

\paragraph{Implementation details.}
Full simulation inputs, parameter assumptions, grid construction procedures, and sensitivity analyses to alternative operational cost assumptions are reported in Chapter~\ref{ch:supp_essay3}.


\subsection{Simulation setup and baseline profile}
\label{subsec:essay3_setup}

Predicted uptake is computed from the estimated choice model as the probability of choosing vaccination (relative to opting out) under counterfactual bundles of delay and incentives. Unless otherwise stated, simulations evaluate uptake at a common reference vaccine profile holding efficacy at 70\%, side effects at mild, and origin as domestic, so that policy comparisons isolate the roles of timeliness and monetary compensation. I evaluate delays over an operationally relevant range (0--6 months) and incentives over a policy-relevant range (0--800 RMB), consistent with the experimental design.

\paragraph{Appendix note.}
Detailed parameter inputs, grid construction, and additional scenario assumptions are documented in Chapter~\ref{ch:supp_essay3}.

\subsection{From policy instruments to predicted uptake}
\label{subsec:essay3_uptake_map}

Let $t_g$ denote the expected delay faced by group $g$ under baseline operations and let $m_g$ denote the incentive offered. A speed-up reduces delay to $t_g' = t_g - \Delta t_g$, while a compensation policy increases incentives to $m_g' = m_g + \Delta m_g$. For each group, predicted uptake is:
\begin{equation}
\widehat{P}_g(t_g', m_g') \;=\; \Pr(\text{choose vaccination}\mid t_g', m_g', x_0),
\end{equation}
where $x_0$ denotes the fixed reference profile for non-policy attributes. This mapping converts any proposed bundle $(\Delta t_g,\Delta m_g)$ into a coverage prediction that can be compared across strategies.

\subsection{Scenario set: uniform versus targeted timing policies}
\label{subsec:essay3_scenarios}

I consider three classes of policy scenarios designed to mirror real programme decisions:

\begin{enumerate}
\item \textbf{Uniform compensation:} the same incentive schedule is offered to all individuals facing a given delay (e.g., 0--800 RMB) without changing operational timing.
\item \textbf{Uniform speed-up:} operational investments reduce delay by the same amount for everyone (e.g., reducing all delays by 1--3 months), with no incentives.
\item \textbf{Equity-sensitive targeting:} speed-ups and/or incentives are targeted toward subgroups with higher delay burdens, defined by higher MWTA and/or higher social welfare weights $\omega_g$ (e.g., low-trust or low-preventive-engagement groups), while other groups receive the baseline policy.
\end{enumerate}

These strategies are evaluated under two implementation objectives. The first is an \emph{uptake-target objective} that aims to achieve a minimum coverage threshold $\bar{P}$ (e.g., 60\% or 80\%) at minimum programme cost. The second is an \emph{equity-weighted welfare objective} that maximizes weighted welfare gains given a fixed budget.

\subsection{Least-cost policy design under uptake targets}
\label{subsec:essay3_least_cost}

Many rollout decisions are driven by operational targets (e.g., reaching a coverage threshold before a given date). To mirror this constraint, I compute the minimum-cost combination of speed-ups and incentives that achieves:
\begin{equation}
\widehat{P}_g(t_g-\Delta t_g,\, m_g+\Delta m_g) \;\ge\; \bar{P}
\quad \text{for relevant groups } g,
\end{equation}
subject to a budget envelope and feasibility constraints on delay reductions. Practically, the algorithm evaluates candidate bundles on a grid and identifies the lowest-cost bundle that satisfies the target uptake condition. This approach yields an implementable ``policy menu'' analogous to iso-acceptance frontiers: for any expected delay, it reports the incentive required to sustain the target uptake, and it allows a direct cost comparison to the operational expense of reducing delay.

\subsection{Equity-weighted targeting under budget constraints}
\label{subsec:essay3_budgeted_equity}

To incorporate distributional objectives, I evaluate policies under group weights $\omega_g$ that upweight welfare gains for priority populations. Under a fixed programme budget $B$, targeted policies are compared by their equity-weighted welfare contributions and their implied coverage outcomes. In practice, this exercise answers: \emph{given limited resources, which groups should receive speed-ups first, and where should compensation be used as a second-best tool?} Consistent with the rules in Section~\ref{sec:essay3_rules}, targeting is expected to favor groups with higher equity-weighted MWTA and lower operational cost per month saved.

\subsection{Outputs and decision-support summaries}
\label{subsec:essay3_outputs}

For each scenario, I report three outputs that are directly interpretable for programme planning:
\begin{enumerate}
\item \textbf{Coverage outcomes:} predicted uptake under each delay--incentive bundle by group and overall.
\item \textbf{Resource requirements:} implied incentive outlays and the operational cost of delay reductions under assumed unit costs.
\item \textbf{Equity implications:} changes in uptake gaps and equity-weighted welfare gains under alternative weighting schemes.
\end{enumerate}

Together, these simulations translate the integrated behavioral parameters into actionable comparisons of ``speed-up versus compensate'' strategies and clarify how targeting rules change under explicit equity objectives.


Table~\ref{tab:scenario_compare} summarizes the predicted coverage,
total programme cost, and distributional implications under
representative policy strategies.


\begin{table}[htbp]
\centering
\caption{Scenario comparison under baseline conditions (3-month delay, 60\% baseline uptake)}
\label{tab:scenario_compare}
\begin{threeparttable}
\begin{tabular}{p{3.2cm} p{2.8cm} p{2.2cm} p{2.2cm} p{2.4cm}}
\toprule
\textbf{Strategy} & \textbf{Instrument} & \textbf{Resulting delay} & \textbf{Predicted uptake} & \textbf{Distributional implication} \\
\midrule
Uniform compensation 
& 150 RMB cash incentive 
& 3 months (unchanged) 
& $\sim$72\% 
& Raises uptake but does not reduce structural waiting; welfare gains proportional to cash transfer \\

Uniform speed-up 
& Reduce delay by 3 months (50 RMB/month) 
& 0 months 
& $\sim$75\% 
& Eliminates delay burden uniformly; higher efficiency when operational cost is close to MWTA \\

Equity-targeted speed-up 
& Allocate 3 months reduction preferentially to high-MWTA groups 
& Mixed (priority groups reach 0 months first) 
& $\sim$76--78\% 
& Concentrates delay reduction where behavioral cost is highest; maximizes equity-weighted welfare under fixed budget \\
\bottomrule
\end{tabular}

\begin{tablenotes}
\footnotesize
\item \textbf{Notes:} Baseline delay is 3 months with 60\% predicted uptake at 0 RMB. 
Operational acceleration cost is assumed to be 50 RMB per person per month saved. 
Total budget is fixed at 150 RMB per person. 
Overall MWTA from Essay~2 is 44.8 RMB per month. 
Predicted uptake values are illustrative simulations based on the estimated choice model.
\end{tablenotes}
\end{threeparttable}
\end{table}

\clearpage


\paragraph{Baseline scenario and parameterization.}
To make the policy comparison concrete, the scenario table evaluates a common baseline environment characterized by a \textbf{three-month delay} and \textbf{60\% baseline uptake at 0 RMB}. The three-month horizon is chosen because it represents a plausible mid-range waiting time during constrained rollouts—longer than short scheduling frictions (0--1 month) but still within the operational window in which programme managers can realistically accelerate access through capacity expansion and queue management. It also aligns with the DCE design range (0--6 months) and corresponds to a policy-relevant ``typical'' delay level rather than an extreme case.

The 60\% baseline uptake is selected to reflect a realistic planning context in which coverage is meaningfully below common programme targets (e.g., 70--80\%), making the speed-up versus compensation trade-off salient. In other words, the baseline is calibrated to a situation where policymakers would plausibly consider spending resources to raise uptake rather than merely maintaining already-high coverage.

For transparency, operational acceleration is parameterized with an illustrative unit cost of \textbf{50 RMB per person per month of delay reduced}. This stylized value provides an intuitive benchmark that can be compared directly to the estimated \textbf{overall MWTA} (44.8 RMB/month) from Essay~2. Finally, the simulations impose a fixed per-person budget of \textbf{150 RMB}, which corresponds to either (i) three months of delay reduction at 50 RMB per month or (ii) a moderate cash incentive of comparable fiscal magnitude. This design ensures that the scenarios are directly comparable under an equal-budget constraint and highlights how instrument choice and targeting alter coverage and distributional outcomes.

\clearpage



\section{Discussion}
\label{sec:essay3_discussion}

This chapter moves from measuring delay preferences to prescribing how vaccination programmes should allocate scarce operational and fiscal resources when timeliness itself is constrained. By integrating behavioral discounting (Essay~1) with money-metric delay valuation (Essay~2), the analysis provides a transparent framework for choosing between accelerating access and compensating unavoidable delay under explicit budget and equity objectives.

Several implications emerge.

First, timeliness should be treated as an economic input rather than a purely logistical variable. Traditional cost-effectiveness frameworks typically focus on vaccine price, delivery cost, and health outcomes, while implicitly assuming uniform time valuation \cite{Drummond2015}. The present analysis demonstrates that waiting time carries heterogeneous welfare costs that can be monetized and compared directly with programme expenditures. In epidemic settings, even short delays may carry substantial behavioral and epidemiological consequences \cite{Abbas2021,Castioni2024}. Incorporating delay valuation into programme design can therefore improve alignment between operational investments and behavioral responsiveness.

Second, the results highlight that acceleration and incentives are not substitutes in all contexts. Operational speed-ups change the timing of protection, while incentives change the immediacy of benefits without altering system capacity. Behavioral research suggests that financial incentives can increase preventive uptake under certain conditions, but their effectiveness depends on design, trust, and perceived legitimacy \cite{Gneezy2011,Thirumurthy2022}. In settings where institutional trust is low and delay burdens are high, the simulations indicate that targeted speed-ups may generate larger equity-adjusted welfare gains than uniform incentive schemes.

Third, the framework makes distributional trade-offs explicit. Delay burdens are not evenly distributed across groups. When uniform policies are applied, populations with higher MWTA effectively bear greater welfare losses from the same waiting time. Distributional cost-effectiveness analysis argues that policy evaluation should incorporate equity weights when welfare gains accrue unevenly \cite{Asaria2016,Cookson2017}. The targeting rules developed here operationalize that principle in the context of vaccination timing: resources should be allocated first where the equity-weighted return per RMB is highest.

Fourth, the analysis provides a structured way to design “least-cost to target” strategies. Many public health campaigns operate under coverage thresholds (e.g., achieving 70\% or 80\% uptake within a defined window). By mapping delay–incentive bundles into predicted uptake, the simulations identify whether operational acceleration or financial compensation is the more cost-effective path to reach a target. This approach is consistent with pragmatic priority-setting frameworks used in health systems under scarcity \cite{Emanuel2020,Norheim2014}.

More broadly, the findings suggest that vaccination policy design benefits from integrating behavioral economics with operational planning. Behavioral discounting explains why short delays may disproportionately suppress uptake, while monetary valuation allows policymakers to quantify and compare policy instruments on a common scale. Together, these tools transform timeliness from an exogenous constraint into a lever that can be optimized.

Finally, the equity-sensitive framework developed here has relevance beyond vaccination. Many preventive interventions—from cancer screening to chronic disease management—involve time trade-offs between immediate inconvenience and delayed health benefits. When delays arise from system capacity limits, similar speed-up versus compensation choices may emerge. Embedding delay valuation into policy design can therefore contribute to more behaviorally informed and distributionally transparent health system decision-making.

\clearpage

\section{Conclusion}
\label{sec:essay3_conclusion}

This essay integrates behavioral discounting and monetary delay valuation to derive implementable policy rules for vaccination programmes operating under timeliness constraints. By linking reduced-form discounting parameters (Essay~1) with money-metric marginal willingness-to-accept (MWTA) estimates (Essay~2), the analysis translates abstract time preferences into operational decisions about when to accelerate access and when to compensate unavoidable delay.

Several conclusions follow.

First, vaccination timeliness is not merely a logistical detail but an economically meaningful policy variable. Individuals attach a positive and heterogeneous shadow price to waiting, and that valuation can be directly compared to the marginal cost of reducing delay. When operational acceleration costs are below the behavioral valuation of delay, speed-up investments dominate monetary compensation; when acceleration is prohibitively expensive, incentives may serve as a second-best instrument. This decision rule formalizes the trade-off between capacity expansion and financial incentives that frequently arises in epidemic rollouts and other preventive health programmes \cite{Drummond2015,Thirumurthy2022}.

Second, heterogeneity fundamentally alters optimal policy design. Delay burdens are not evenly distributed across populations. Groups characterized by lower institutional trust, lower preventive engagement, or higher behavioral impatience experience larger welfare losses from identical waiting periods. Uniform policies therefore risk embedding implicit inequities. Incorporating explicit welfare weights into allocation rules aligns programme design with principles of distributional cost-effectiveness analysis, which emphasize that both efficiency and equity should guide health policy decisions \cite{Asaria2016,Cookson2017,Norheim2014}.

Third, the simulations demonstrate that equity-sensitive targeting can reallocate resources toward high-delay-burden groups without increasing total expenditure. By prioritizing operational acceleration where the equity-weighted return per RMB is highest, policymakers can simultaneously improve coverage and reduce welfare disparities. Importantly, this framework makes normative assumptions transparent: different equity weights produce different allocation outcomes, allowing deliberative justification rather than implicit value judgments.

More broadly, this essay illustrates how behavioral economics can be integrated into pragmatic priority-setting frameworks. Traditional economic evaluation often treats timing as exogenous, focusing on cost per health outcome achieved \cite{Drummond2015}. The present framework extends that approach by embedding time valuation directly into policy design. In contexts where preventive actions involve waiting—vaccination campaigns, screening programmes, or chronic disease management—speed-up versus compensation decisions are likely to recur. Embedding delay valuation into policy simulation tools can therefore improve alignment between behavioral responses and operational planning.

Taken together, Essays~1--3 provide a coherent account of vaccination timing as an intertemporal health investment problem. Essay~1 identifies the behavioral structure of delay sensitivity, Essay~2 monetizes the shadow price of waiting, and Essay~3 converts those estimates into actionable and equity-sensitive policy rules. By bridging micro-level preference heterogeneity with macro-level programme design, this dissertation contributes a behaviorally informed and distributionally transparent framework for improving the timeliness and fairness of vaccination policy.

\clearpage

\chapter{General Discussion and Conclusion}
\label{ch:conclusion}

\section{What this dissertation studied}
\label{sec:dissertation_overview_studied}

This dissertation studied vaccination timing as an intertemporal policy problem. Rather than treating vaccination as a binary decision of acceptance versus refusal, I focus on \emph{when} individuals are willing to receive protection and how that timing is shaped by behavioral preferences, uncertainty, and policy design. Using a discrete choice experiment (DCE) fielded among 1,027 adults in Wuhan, China, the dissertation quantifies how people trade off delayed vaccine availability against vaccine performance attributes and monetary incentives, and how these trade-offs vary across population subgroups.

Substantively, the dissertation examined three linked questions. First, it estimated whether and how individuals discount delayed vaccine protection, mapping observed delay sensitivity into reduced-form behavioral discounting measures that capture present-oriented valuation of near-term delays. Second, it monetized the behavioral cost of waiting by estimating marginal willingness-to-accept (MWTA) for vaccination delay—the cash compensation required to offset one additional month of waiting—thereby translating timing preferences into policy-actionable monetary units. Third, it integrated the discounting and monetary valuation evidence to derive equity-sensitive policy rules for vaccination programmes deciding between accelerating access (reducing waiting time operationally) and offering compensation (incentives) when delays are unavoidable.

Methodologically, the dissertation used standard random utility models for stated-preference data, including multinomial logit as a benchmark and mixed logit models to account for preference heterogeneity. The analysis emphasized transparent interpretation and decision support: beyond estimating coefficients, it produced uptake simulations under delay--cash bundles and visual tools such as iso-acceptance frontiers that map operational delay scenarios to required incentive levels for maintaining target uptake thresholds. Together, the dissertation contributes a behavioral-economic framework for treating timeliness as a measurable policy lever—one that can be compared, targeted, and optimized alongside conventional programme costs and objectives.


\section{Synthesis of findings across the three essays}
\label{sec:dissertation_synthesis}

The three essays provide a cumulative account of how vaccination timing is valued and how that valuation can be used for policy design. Together, they move from measurement (discounting and monetization) to prescription (equity-sensitive decision rules).

\paragraph{Essay 1: Behavioral discounting in delayed protection.}
Essay~1 establishes that vaccination timing is governed by intertemporal preferences: individuals place lower value on protection that arrives later, and the marginal disutility of delay is consistent with present-oriented discounting. By mapping estimated delay sensitivity into reduced-form discounting measures, the essay shows that short postponements can carry disproportionate behavioral cost relative to longer horizons. It also documents meaningful heterogeneity in delay sensitivity across demographic characteristics, preventive behaviors, and institutional trust—indicating that a uniform waiting time does not impose a uniform welfare burden.

\paragraph{Essay 2: Monetizing the cost of waiting.}
Essay~2 translates time preferences into a money-metric policy input by estimating marginal willingness-to-accept (MWTA) for vaccination delay. The central result is that waiting time has a substantial shadow price: respondents require nontrivial compensation to tolerate an additional month of delay. Importantly, MWTA varies systematically across subgroups, with particularly high delay valuation among populations characterized by lower institutional trust and lower preventive engagement. By simulating uptake under delay--cash bundles and producing iso-acceptance frontiers, the essay provides decision-support outputs that connect operational delays to incentive requirements.

\paragraph{Essay 3: Integrating evidence into equity-sensitive policy rules.}
Essay~3 synthesizes the behavioral discounting results and the MWTA estimates to derive implementable policy rules for vaccination programme design. The framework treats timeliness as a scarce policy lever and compares two instruments on a common scale: operational speed-ups that reduce delay versus monetary incentives that compensate unavoidable waiting. Because the burden of delay is heterogeneous, the essay shows how targeting decisions can be made explicitly equity-sensitive—prioritizing groups with higher delay costs (and, where desired, higher social welfare weights) when allocating limited resources. The simulation framework illustrates how the preferred policy mix shifts with operational costs, budget constraints, coverage targets, and equity objectives.

\paragraph{Overall implication.}
Across the three essays, a consistent conclusion emerges: vaccination timing is economically meaningful, behaviorally heterogeneous, and policy-actionable. Measuring delay sensitivity (Essay~1) explains why timeliness matters; monetizing waiting (Essay~2) yields a practical exchange rate between time and money; and integrating these inputs (Essay~3) provides a structured way to choose between speeding up access and compensating delay under real-world constraints. This synthesis positions timeliness not as an exogenous feature of rollouts, but as a design dimension that can be optimized to improve both efficiency and equity in future epidemic vaccination programmes.


\section{Policy implications}
\label{sec:dissertation_policy_implications}

This dissertation has several implications for vaccination programme design, especially when timeliness is constrained by supply, staffing, or appointment capacity. The core message is that \emph{waiting time is not merely a logistical inconvenience}: it is a measurable behavioral burden with heterogeneous welfare costs, and it can be managed with policy instruments that should be compared on a common scale.

\subsection{Treat timeliness as a policy lever, not an operational residual}
The findings imply that vaccination policy should explicitly include timeliness as a design dimension. Reducing delay can generate disproportionately large behavioral returns when present-oriented preferences make near-term waiting especially costly. Practically, this supports prioritizing operational investments that reduce \emph{early} waiting time—such as appointment throughput, staffing, queue management, and supply coordination—rather than focusing only on later-stage catch-up efforts. In planning terms, timeliness should be monitored and reported as a performance metric alongside coverage rates.

\subsection{Use MWTA as a decision tool for calibrating incentives and budgeting}
The MWTA estimates provide a policy-relevant ``exchange rate'' between months of delay and RMB compensation. This enables planners to answer concrete questions: \emph{if delays are unavoidable, how much incentive is required to maintain uptake?} and \emph{is it cheaper to buy time (speed up access) or buy acceptance (compensate waiting)?} By expressing delay burdens in monetary units, MWTA can be incorporated into programme budgeting, incentive calibration, and scenario planning under alternative delay forecasts.

\subsection{Design incentives as a second-best tool when acceleration is infeasible}
The results do not imply that incentives should always substitute for operational improvements. Instead, incentives are most defensible as a second-best tool when speed-ups are infeasible in the short run or when the marginal cost of reducing delay exceeds the behavioral value of timeliness. In such cases, front-loaded, transparent incentive schemes can partially offset the welfare loss from waiting and help sustain target uptake thresholds. However, incentive design should account for administrative feasibility and legitimacy concerns, particularly in low-trust contexts.

\subsection{Target policies because delay burdens are heterogeneous}
A central implication is that uniform timing policies can generate unequal welfare burdens. Groups with higher delay valuation—such as those with lower institutional trust or lower preventive engagement—bear larger welfare losses from the same waiting time. This motivates targeted approaches, including reserved appointment capacity, priority scheduling, mobile delivery, and tailored outreach, alongside differential incentive calibration. Targeting does not require abandoning universalism; rather, it supports using scarce marginal resources where they yield the largest welfare and equity gains.

\subsection{Embed equity objectives explicitly in rollout decision-making}
Essay~3 shows how equity considerations can be incorporated transparently using subgroup welfare weights or priority rules. This provides a structured way to balance average coverage goals against distributional objectives when resources are limited. In practice, programmes can pre-specify equity priorities (e.g., prioritizing disadvantaged or high-burden groups) and then evaluate whether proposed speed-up or compensation policies align with those priorities. This approach improves accountability by making distributional trade-offs explicit rather than implicit.

\subsection{Implications beyond vaccination}
Finally, the framework extends to other preventive interventions where benefits arrive in the future but costs are immediate—screening, chronic disease management, and adherence-support programmes. When system frictions generate waiting time, the same logic applies: quantify the behavioral cost of delay, compare operational investments to compensatory policies, and target scarce resources to reduce unequal delay burdens.

Overall, these implications suggest that future epidemic preparedness should integrate behavioral valuation of time into operational planning, combining acceleration and compensation in a way that is both cost-effective and equity-sensitive.
 targeting; equity

\section{Methodological contributions}
\label{sec:dissertation_method_contrib}

This dissertation makes four methodological contributions to the study of vaccination behavior and the economic evaluation of timeliness policies.

\subsection{Measuring vaccination timing preferences using a unified DCE backbone}
First, the dissertation implements a discrete choice experiment (DCE) designed to identify preferences over \emph{when} vaccination occurs, not only whether it occurs. By embedding waiting time as an explicit attribute alongside efficacy, side effects, origin, and monetary incentives, the design provides clean preference-based identification of delay sensitivity under controlled variation. Importantly, the same experimental backbone supports all three essays, enabling direct comparability across discounting measures, money-metric valuation, and policy simulations without changing the underlying preference elicitation environment.

\subsection{Mapping reduced-form delay sensitivity into behavioral discounting measures}
Second, Essay~1 contributes a transparent mapping from estimated delay sensitivity in a random utility model to reduced-form behavioral discounting measures. This approach preserves the interpretability of intertemporal preference heterogeneity while remaining empirically tractable over operationally relevant monthly horizons. It allows the dissertation to connect stated-preference evidence to established concepts in intertemporal choice—such as present-oriented valuation—without requiring strong structural assumptions that are difficult to identify in typical vaccine rollout settings.

\subsection{Monetizing waiting time via MWTA and integrating inference for ratio estimands}
Third, Essay~2 monetizes vaccination delay by deriving marginal willingness-to-accept (MWTA) as a ratio of coefficients, thereby producing a policy-actionable measure of the behavioral cost of delay (RMB per month). The dissertation complements standard analytical inference with robustness tools tailored to nonlinear ratio estimands, including clustered bootstrap procedures, trimmed-sample sensitivity tests, and alternative model specifications. This strengthens credibility and provides a replicable template for estimating and validating money-metric trade-offs from DCEs in other public health contexts.

\subsection{Bridging behavioral estimation and policy design through decision-support simulations}
Fourth, Essay~3 contributes a practical integration of behavioral estimates into implementable policy rules and decision-support simulations. Rather than stopping at coefficient interpretation, the dissertation translates estimated preferences into uptake predictions under delay--incentive bundles, iso-acceptance frontiers for target coverage thresholds, and comparative cost scenarios for operational speed-ups versus compensation policies. The framework is flexible: it can incorporate explicit budget constraints, group-specific unit costs of reducing delay, and equity-sensitive welfare weights, making the simulation outputs directly usable for programme planning.

Together, these methodological contributions advance a general approach for treating timeliness as a measurable, monetizable, and optimizable policy dimension in preventive health interventions.

% DCE for timing; mapping; iso-acceptance; decision-support

\section{Limitations}
\label{sec:dissertation_limitations}

Several limitations should be considered when interpreting the dissertation’s findings and when applying the proposed policy framework to other settings.

\subsection{Stated-preference context and hypothetical bias}
The empirical results are based on a discrete choice experiment, which elicits stated preferences rather than observing real vaccination behavior. Although the survey design incorporates realism and internal-validity checks, respondents may still overstate or understate their willingness to vaccinate under particular delay--incentive bundles. As a result, predicted uptake levels and money-metric valuations should be interpreted as \emph{behavioral preference indicators} rather than literal forecasts of realized uptake under every operational scenario.

\subsection{Limited delay horizon and functional-form assumptions}
The DCE focuses on operationally relevant monthly delays (0--6 months). This range is well suited to epidemic rollout planning but does not capture longer-run intertemporal trade-offs. The mapping from delay sensitivity into reduced-form discounting measures (Essay~1) also relies on parsimonious functional forms to summarize time preferences. These choices support interpretation, but they may not capture all features of intertemporal valuation, especially over longer horizons or under richer uncertainty processes.

\subsection{Context specificity and external validity}
The study was fielded in Wuhan, China, in an environment shaped by the salience of epidemic risk and local institutional conditions. The level of institutional trust, perceptions of vaccine origin, and baseline acceptance may differ in other regions or countries.

Importantly, the generalizability of these findings to international health systems depends on the underlying structure of vaccine delivery and payment. In the Wuhan context, vaccination was centralized and provided free of charge, making "waiting time" and "monetary incentives" the primary policy levers. In international systems with private insurance or multi-payer models, the behavioral cost of delay may interact with out-of-pocket costs and varying levels of healthcare access. Furthermore, the high level of institutional trust typically observed in centralized systems may not be present in more fragmented or pluralistic societies, potentially leading to even higher MWTA for delayed protection in those contexts. While the conceptual framework—monetizing delay and comparing it to operational costs—is general, the numerical MWTA values and subgroup rankings may not transfer one-for-one to settings with different baseline acceptance, risk perceptions, or insurance structures. Applying the framework elsewhere would ideally involve local preference elicitation or calibration using context-specific data.

\subsection{Operational cost inputs are modeled rather than directly estimated}
Essay~3 compares speed-ups and compensation using stylized operational cost assumptions (e.g., cost per person-month of delay reduction). These cost parameters are necessary for policy comparison but are not directly measured from administrative cost data in this dissertation. The simulation results should therefore be interpreted as \emph{threshold and scenario analyses}—identifying when one instrument dominates the other under plausible cost ranges—rather than definitive statements about any specific programme’s cost-effectiveness.

\subsection{Limited modeling of epidemic dynamics and spillovers}
The framework evaluates delay as an individual-level welfare and uptake determinant and does not explicitly model transmission dynamics, herd protection, or general equilibrium responses. In practice, accelerating vaccination can generate spillover benefits by reducing infections and mitigating peaks, which would strengthen the case for speed-ups beyond what individual preference-based valuation alone implies. Similarly, incentives may have informational or social effects not captured in the static choice setting.

\subsection{Equity weights require normative choices}
The equity-sensitive targeting framework makes distributional objectives explicit through subgroup welfare weights, but selecting these weights is ultimately normative. Different stakeholders may place different priority on vulnerability, disadvantage, or trust-related inequities. The dissertation’s contribution is to provide a transparent structure for incorporating such priorities; it does not claim a single ethically ``correct'' weighting scheme.

Despite these limitations, the dissertation provides a coherent measurement-to-policy pipeline for quantifying the behavioral cost of vaccination delay, translating it into money-metric terms, and using it to support policy design under scarcity and equity objectives.

% stated preference; external validity; dynamics; context

\section{Directions for future research}
\label{sec:dissertation_future_research}

The dissertation opens several avenues for future research on vaccination timing and, more broadly, on intertemporal health policy design.

\subsection{Linking stated preferences to realized behavior}
A priority is to validate and calibrate DCE-based delay valuations using revealed-preference data. Future work could link preference estimates to administrative vaccination records or field experiments that vary appointment timing and incentive design in real settings. Such validation would help quantify the magnitude of hypothetical bias (if any) and improve the predictive use of delay--cash simulations for programme planning.

\subsection{Embedding timing preferences in dynamic epidemic and operational models}
The current framework treats delay primarily as an individual-level welfare and uptake determinant, potentially understating the true social cost of waiting. Future research could integrate the estimated timing preferences into dynamic transmission models (e.g., SEIR-based compartmental models) and operational queueing models to quantify the full social value of accelerating access, including spillovers from reduced infections and avoided peaks. 

Currently, the gap between behavioral delay valuation and epidemiological externalities remains a significant area for development. A complete accounting would acknowledge that while MWTA measures individual welfare loss, it does not capture the "Social Shadow Price" of delay—the aggregated cost of potential infections prevented by faster protection. Combining behavioral MWTA with these epidemiological externalities would provide a more complete cost-benefit foundation for timeliness investments and offer a more robust case for capacity expansion over monetary compensation when outbreaks are active.

\subsection{Richer representations of uncertainty and trust}
The conceptual framework highlights uncertainty and institutional trust as key drivers of impatience toward delayed protection. Future studies could experimentally vary information quality, credibility cues, and uncertainty about vaccine effectiveness or delivery reliability to disentangle the mechanisms linking trust to delay valuation. This would improve causal interpretation and help tailor communication and delivery strategies in low-trust environments.

\subsection{Heterogeneity, targeting, and implementation constraints}
The dissertation documents strong heterogeneity in delay valuation and proposes equity-sensitive targeting rules. Future work could study the administrative feasibility, political acceptability, and ethical trade-offs of targeted speed-ups and targeted incentives, including potential stigma or perceptions of unfairness. Evaluating alternative targeting criteria—such as vulnerability, disadvantage, or behavioral responsiveness—could inform practical programme design.

\subsection{Cost measurement for operational speed-ups}
A key input for policy rules in Essay~3 is the cost of reducing delay. Future research should collect and harmonize micro-costing data on operational interventions that shorten waiting time (e.g., additional staffing, mobile clinics, extended hours, IT throughput, and outreach). Better cost evidence would allow the framework to deliver more precise recommendations and facilitate cost-effectiveness comparisons across delivery strategies.

\subsection{Generalizing the framework beyond vaccination}
Finally, the approach can be extended to other preventive and chronic-care settings where benefits arrive later but costs are immediate, and where system frictions create waiting time (e.g., screening, medication initiation, and follow-up care). Applying the same measurement-to-policy pipeline—estimating time valuation, monetizing delay, and deriving targeting rules—could help design more behaviorally informed and equity-sensitive health systems.

Together, these directions emphasize that understanding how people value \emph{time} in health decisions is not only a behavioral question but also a central input for practical policy design under scarcity and uncertainty.

\section{Closing remarks}
\label{sec:closing_remarks}

This dissertation studied vaccination timing as an intertemporal health decision under operational frictions. Using a discrete choice experiment fielded among adults in Wuhan, China, the three essays jointly quantified (i) how strongly individuals discount delayed protection, (ii) the monetary value they place on avoiding delay, and (iii) how these behavioral parameters can be translated into actionable and equity-sensitive policy rules.

Several themes emerge across the essays. First, timing preferences are consequential: even modest delays carry meaningful behavioral costs that can depress uptake and disproportionately burden particular groups. Second, monetizing delay via MWTA provides a policy-ready metric---a ``shadow price of waiting''---that links individual preferences to budgeting, incentive calibration, and operational planning. Third, optimal programme design is inherently a trade-off problem: when capacity expansions are costly or infeasible in the short run, well-calibrated compensation can sustain uptake; when speed-ups are affordable, reducing waiting time can deliver larger welfare gains than equivalent spending on incentives. Finally, heterogeneity is not a secondary detail but a central design input: the same delay or incentive policy can produce uneven welfare consequences across subgroups differentiated by vulnerability, health behaviors, and institutional trust.

The broader implication is that vaccination policy should treat \emph{time} as a scarce resource with measurable value. Incorporating behavioral delay costs into programme design can improve both efficiency and fairness, particularly during epidemic settings where timeliness directly shapes risk exposure. By providing empirically grounded measures of delay valuation and a transparent framework for choosing between speed-up investments and compensation, this dissertation offers a practical toolkit for designing vaccination programmes that are not only effective, but also responsive to how populations experience and value waiting.








\clearpage
% TODO: Add Essay 3 content.



\appendix

\chapter{Common DCE Design and Survey Instrument}
\label{app:common_dce}




\chapter{Supplementary Results for Essay 2}
\label{ch:supp_essay2}



This chapter reports additional validity and robustness evidence for Essay~2, with emphasis on inference for monetized delay valuation (MWTA) and the stability of the main preference estimates. Consistent with the dissertation’s shared DCE backbone, several general internal-validity diagnostics are documented in Chapter~\ref{ch:supp_essay1}; the analyses below focus on robustness checks that are particularly relevant for the monetization exercise in Essay~2.

\section{Validity checks and robustness summary}
\label{sec:validity_robustness}

\subsection{Bootstrap standard errors}
To assess estimation stability and sampling variability, I compute non-parametric bootstrap standard errors using 1,000 replications clustered at the respondent level. For each replication, the full choice dataset is resampled with replacement at the individual level and the multinomial logit (MNL) model is re-estimated. Table~\ref{tab:essay2_bootstrap_se} reports bootstrap standard errors and percentile-based confidence intervals. The resulting intervals closely match the analytical ones reported in the main text, indicating that the core coefficient estimates are not sensitive to sampling noise.

\begin{table}[htbp]
\centering
\caption{Bootstrap standard errors and 95\% percentile intervals (1,000 clustered replications)}
\label{tab:essay2_bootstrap_se}
\begin{threeparttable}
\begin{tabular}{lccc}
\toprule
Attribute & Mean coefficient & Bootstrap SE & 95\% CI \\
\midrule
Vaccine delay (per month) & -0.125 & 0.027 & [-0.177,\; -0.072] \\
Monetary incentives (per 100 RMB) & 0.313 & 0.018 & [0.278,\; 0.347] \\
Vaccine efficacy (per 10 pp) & 0.163 & 0.025 & [0.115,\; 0.212] \\
Side effects (per severity) & -0.082 & 0.033 & [-0.147,\; -0.017] \\
Imported origin (dummy) & 0.239 & 0.048 & [0.145,\; 0.333] \\
\bottomrule
\end{tabular}
\begin{tablenotes}
\footnotesize
\item Notes: Non-parametric bootstrap with 1,000 replications clustered at the respondent level. Confidence intervals are percentile-based.
\end{tablenotes}
\end{threeparttable}
\end{table}

\subsection{Trimmed-sample sensitivity tests}
To examine whether extreme preference values influence the results, I conduct trimmed-sample analyses by removing respondents in the upper and lower 5\% of the MWTA distribution and re-estimating the benchmark MNL model. As shown in Table~\ref{tab:essay2_trimmed}, the re-estimated coefficients differ by less than 3\% from the full-sample values, indicating that outliers are not driving the main behavioral patterns.

\begin{table}[htbp]
\centering
\caption{Trimmed-sample sensitivity analysis: comparison with full-sample estimates}
\label{tab:essay2_trimmed}
\begin{threeparttable}
\begin{tabular}{lccc}
\toprule
Attribute & Full sample & Trimmed sample & Change (\%) \\
\midrule
Vaccine delay (per month) & -0.125 & -0.122 & -2.4 \\
Monetary incentives (per 100 RMB) & 0.313 & 0.309 & -1.3 \\
Vaccine efficacy (per 10 pp) & 0.163 & 0.161 & -1.2 \\
Side effects (per severity) & -0.082 & -0.084 & +2.4 \\
Imported origin (dummy) & 0.239 & 0.236 & -1.3 \\
\bottomrule
\end{tabular}
\begin{tablenotes}
\footnotesize
\item Notes: The trimmed sample removes respondents in the upper and lower 5\% of the MWTA distribution. Percent change is computed relative to the full-sample estimate.
\end{tablenotes}
\end{threeparttable}
\end{table}

\subsection{Mixed logit re-estimation}
To evaluate sensitivity to the independence of irrelevant alternatives (IIA) assumption and to allow for unobserved heterogeneity, I estimate a mixed logit (MXL) model with random coefficients on delay and cash and fixed coefficients for the remaining attributes. This specification balances flexibility and empirical stability given the limited number of choice tasks per respondent. The estimated means of the random coefficients closely match the MNL estimates, and the implied MWTA values remain within the 95\% confidence intervals of the benchmark results. Several variance components are imprecisely estimated---a common feature of short DCE panels---but the overall ranking and magnitude of attribute effects are unchanged. These results indicate that the main behavioral conclusions are robust to relaxing the IIA assumption.

\subsection{Delay--monetary incentives interaction model}
To test whether monetary incentives compensate differently for short versus long delays, I augment the MNL model with an explicit interaction between delay and monetary incentives. Both variables are mean-centered prior to constructing the interaction term to reduce collinearity. The interaction coefficient is small and statistically insignificant ($\hat{\beta}_{int}=0.034$, SE = 0.250, $p=0.89$), and its inclusion does not materially alter the main delay or cash coefficients. This suggests that the marginal effectiveness of incentives does not systematically vary with waiting time within the design range, consistent with the approximately linear iso-acceptance contours reported in the main text.

\begin{table}[htbp]
\centering
\caption{Delay--monetary incentives interaction model (MNL, centered covariates)}
\label{tab:essay2_interaction}
\begin{threeparttable}
\begin{tabular}{lccc}
\toprule
Coefficient & Estimate & Std. error & $t$-value \\
\midrule
ASC (B vs.\ A) & -0.5864 & 3.5419 & -0.1656 \\
ASC (C vs.\ A) & -1.6224 & 5.5255 & -0.2936 \\
WaitTime$_c$ & -0.0746 & 0.8287 & -0.0900 \\
Cash$_c$ & 0.0008 & 0.6015 & 0.0014 \\
E70 & 0.8558 & 4.9993 & 0.1712 \\
E95 & 0.3405 & 3.9703 & 0.0858 \\
SE\_Moderate & 0.5523 & 4.0829 & 0.1353 \\
Imported & -0.6511 & 3.4941 & -0.1863 \\
WaitTime$_c \times$ Cash$_c$ & 0.0345 & 0.2496 & 0.1382 \\
\bottomrule
\end{tabular}
\begin{tablenotes}
\footnotesize
\item Notes: WaitTime$_c$ and Cash$_c$ denote mean-centered delay (months) and monetary incentives (per 100 RMB), respectively. Standard errors are clustered at the respondent level. The small and statistically insignificant interaction term indicates that the marginal utility of cash does not vary systematically with delay within the design range.
\end{tablenotes}
\end{threeparttable}
\end{table}

\subsection{Summary}
Across these robustness assessments, the DCE estimates exhibit high stability and theoretical coherence. Bootstrap and trimmed-sample analyses confirm that results are not driven by sampling variation or extreme observations. The mixed logit re-estimation and interaction test further indicate that the main behavioral conclusions are robust to relaxing IIA and to allowing for potential nonlinearity in delay--cash compensation within the experimental range. Collectively, these diagnostics support the credibility of the estimated MWTA values and the associated policy simulations in Essay~2.


\chapter{Supplementary Results for Essay 3}
\label{ch:supp_essay3}

\section{Simulation framework}
\label{sec:essay3_sim_framework}

This section describes how the empirical estimates from Essays~1--2 are translated into policy-relevant counterfactuals for Essay~3. The simulations are designed to be transparent and directly interpretable for programme planning: they map operational delays and monetary incentives into predicted uptake, and they quantify the resources required to achieve coverage targets under alternative targeting rules.

\subsection{Core inputs from Essays~1--2}
\label{subsec:essay3_inputs}

The simulation framework uses two sets of estimated objects:

\begin{enumerate}[leftmargin=*]
\item \textbf{Choice-model preference parameters.} The estimated coefficients on vaccination delay and monetary incentives, together with coefficients on other attributes (held fixed in baseline profiles), determine predicted uptake under any delay--incentive bundle.
\item \textbf{Money-metric time valuation.} Group-specific MWTA estimates (RMB per month) provide a welfare-consistent exchange rate between months saved and RMB transferred, enabling comparisons of speed-ups versus compensation on a common scale.
\end{enumerate}

These inputs are estimated from the common DCE backbone shared with Essays~1--2; therefore, the simulations do not introduce additional behavioral assumptions beyond those already embedded in the estimated preference parameters.

\subsection{Policy instruments and counterfactual bundles}
\label{subsec:essay3_instruments}

For each subgroup $g$, baseline rollout conditions are summarized by an expected delay $t_g$ (months) and an offered incentive $m_g$ (RMB). Policy choices are represented as changes to these two levers:

\begin{itemize}[leftmargin=*]
\item \textbf{Speed-up (operational acceleration):} reduces delay by $\Delta t_g \ge 0$, so that $t_g' = t_g - \Delta t_g$.
\item \textbf{Compensation (monetary incentives):} increases incentives by $\Delta m_g \ge 0$, so that $m_g' = m_g + \Delta m_g$.
\end{itemize}

A policy therefore corresponds to a bundle $(t_g',m_g')$ for each group. In the baseline analysis, delays are evaluated over 0--6 months and incentives over 0--800 RMB, matching operationally relevant planning ranges and the experimental support of the DCE.

\subsection{Predicting uptake}
\label{subsec:essay3_predict_uptake}

Predicted uptake is computed using the estimated discrete-choice model as the probability of selecting vaccination (relative to opting out) under a counterfactual bundle:
\begin{equation}
\widehat{P}_g(t_g',m_g') \;=\; \Pr(\text{vaccinate}\mid t_g',m_g',x_0),
\end{equation}
where $x_0$ denotes a fixed reference profile for non-policy attributes (e.g., efficacy, side-effect severity, and origin). Holding $x_0$ constant isolates the policy-relevant trade-off between \emph{timeliness} and \emph{compensation}. Results are presented both for (i) subgroup-specific uptake and (ii) an overall uptake measure obtained by population-weighted averaging across groups.

\subsection{Coverage targets and iso-acceptance frontiers}
\label{subsec:essay3_iso_acceptance}

To align the analysis with programme objectives, I evaluate policies relative to uptake targets $\bar{P}$ (e.g., 60\% and 80\%). For any given delay level, the minimum incentive required to achieve the target is defined implicitly by:
\begin{equation}
\widehat{P}_g(t, m^\ast_g(t;\bar{P})) = \bar{P}.
\end{equation}
The function $m^\ast_g(t;\bar{P})$ traces an \emph{iso-acceptance frontier}: combinations of delay and compensation that sustain a target level of predicted uptake. These frontiers provide a direct decision-support mapping from operational delay scenarios to incentive calibration.

\subsection{Cost accounting: speed-up costs versus incentive outlays}
\label{subsec:essay3_costs}

The simulations compare the two instruments using a common RMB-per-person cost metric.

\paragraph{Incentive costs.}
Total incentive outlays are computed mechanically as the incentive amount multiplied by the eligible population:
\begin{equation}
\text{Cost}^{M} = \sum_g N_g \cdot \Delta m_g,
\end{equation}
where $N_g$ is the size of subgroup $g$.

\paragraph{Speed-up costs.}
Operational acceleration is represented by a per-person unit cost per month of delay reduced, $c_g$ (RMB per month saved). The associated programme cost is:
\begin{equation}
\text{Cost}^{S} = \sum_g N_g \cdot c_g \cdot \Delta t_g,
\end{equation}
with $c_g$ allowed to vary across groups in sensitivity analyses to reflect heterogeneous delivery difficulty (e.g., urban versus rural outreach).

\subsection{Targeting rules and equity weights}
\label{subsec:essay3_targeting}

To incorporate equity objectives, the framework allows subgroup welfare weights $\omega_g$ that upweight benefits accruing to priority populations. Two families of targeting strategies are evaluated:

\begin{enumerate}[leftmargin=*]
\item \textbf{Uniform strategies:} the same speed-up or incentive schedule is applied broadly without differentiation.
\item \textbf{Equity-sensitive strategies:} speed-ups and/or incentives are preferentially allocated to groups with higher equity-weighted delay burden, proxied by higher $\omega_g \cdot \text{MWTA}_g$.
\end{enumerate}

In practice, this produces a ranked allocation rule: under a fixed budget, groups with larger $\omega_g \cdot \text{MWTA}_g$ and lower operational unit costs receive speed-ups first; groups for whom speed-ups are relatively expensive are more likely to receive compensation.

\subsection{Outputs reported}
\label{subsec:essay3_outputs}

For each policy scenario, the simulations report:
\begin{enumerate}[leftmargin=*]
\item \textbf{Predicted uptake} by subgroup and overall under each delay--incentive bundle.
\item \textbf{Minimum incentive requirements} to sustain uptake targets under given delay scenarios (iso-acceptance frontiers);
\item \textbf{Total programme costs} for speed-ups versus incentives under assumed unit cost parameters;
\item \textbf{Equity implications}, including changes in uptake gaps and equity-weighted welfare comparisons across targeting rules.
\end{enumerate}

\paragraph{Appendix note.}
Full implementation details (parameter inputs, grid resolution, and additional sensitivity analyses for unit costs and welfare weights) are provided in Chapter~\ref{ch:supp_essay3}.


\subsection{Baseline uptake prediction model}
\label{subsec:essay3_baseline_uptake}

The simulations require a baseline model that maps any policy bundle—defined by vaccination delay and monetary incentives—into predicted vaccination uptake. I use the estimated discrete-choice model from Essays~1--2 and interpret \emph{baseline uptake} as the predicted probability that an individual selects a vaccination alternative rather than opting out, evaluated at a fixed reference vaccine profile.

\paragraph{Utility specification.}
For respondent $i$ in subgroup $g$, the indirect utility from vaccination alternative $j$ is:
\begin{equation}
U_{ijg} \;=\; \alpha_{g} \;+\; \beta_{t,g}\,t_{ij} \;+\; \beta_{m,g}\,m_{ij}
\;+\; \beta_{x,g}'x_{ij} \;+\; \varepsilon_{ijg},
\label{eq:utility_baseline}
\end{equation}
where $t_{ij}$ is vaccination delay (months), $m_{ij}$ is the monetary incentive (RMB, scaled as in Essay~2), and $x_{ij}$ collects the remaining vaccine attributes (e.g., efficacy, side-effect severity, origin). The constant $\alpha_g$ captures baseline inclination toward vaccination (relative to the opt-out) and can vary across subgroups. The error term $\varepsilon_{ijg}$ follows the standard extreme-value distribution under the multinomial logit benchmark, with mixed logit extensions used for robustness.

\paragraph{Reference profile and baseline scenario.}
To isolate the policy levers of interest, predicted uptake is evaluated at a common reference vaccine profile $x_0$ (e.g., efficacy fixed at 70\%, side effects fixed at mild, and origin fixed as domestic), consistent with the reporting convention in Essay~2. The baseline uptake for subgroup $g$ under a bundle $(t,m)$ is therefore:
\begin{equation}
\widehat{P}_g(t,m) \;=\; \Pr(\text{vaccinate}\mid t,m,x_0)
\;=\;
\frac{\exp\!\left(\alpha_g + \beta_{t,g}t + \beta_{m,g}m + \beta_{x,g}'x_0\right)}
{1+\exp\!\left(\alpha_g + \beta_{t,g}t + \beta_{m,g}m + \beta_{x,g}'x_0\right)}.
\label{eq:uptake_logit}
\end{equation}
This expression is the predicted choice probability for vaccination when the outside option (no vaccination) is normalized to zero utility. For multi-alternative choice sets, the same logic applies by summing over the available vaccination alternatives and comparing to the opt-out; the simulation uses the estimated model to compute probabilities directly rather than relying on closed-form approximations when multiple vaccination profiles are included.

\paragraph{Population aggregation.}
Overall predicted uptake is computed as the population-weighted average across subgroups:
\begin{equation}
\widehat{P}(t,m) \;=\; \sum_{g} \pi_g\,\widehat{P}_g(t,m),
\label{eq:aggregate_uptake}
\end{equation}
where $\pi_g$ denotes the subgroup population share (from the sample or external benchmarks, depending on the scenario).

\paragraph{Interpretation.}
The baseline uptake model provides a transparent mapping from operational conditions (delay) and policy instruments (cash incentives) to predicted uptake. This mapping is used throughout the chapter to (i) compute uptake under counterfactual delay--incentive bundles, (ii) trace iso-acceptance frontiers for target uptake thresholds, and (iii) evaluate the relative cost-effectiveness of speed-ups versus compensation policies under alternative equity-weighting schemes.


\subsection{Operational cost assumptions}
\label{subsec:essay3_operational_costs}

A central comparison in Essay~3 is between (i) reducing waiting time through operational investments and (ii) compensating unavoidable delays through monetary incentives. To put these two instruments on a common budget scale, the simulations require explicit assumptions about the programme cost of achieving a given reduction in delay. Because detailed micro-costing data for epidemic rollouts are rarely available ex ante, I adopt a transparent reduced-form cost representation and examine sensitivity across plausible ranges.

\paragraph{Unit cost per month saved.}
Let $c_g$ denote the per-person operational cost (RMB) of reducing expected vaccination delay by one month for subgroup $g$. The total cost of a speed-up policy that reduces delay by $\Delta t_g$ months is modeled as:
\begin{equation}
\text{Cost}^{S} \;=\; \sum_{g} N_g \cdot c_g \cdot \Delta t_g,
\label{eq:speedup_cost}
\end{equation}
where $N_g$ is the size of subgroup $g$. This linear specification captures the idea that the marginal resources required to save time—additional appointment slots, staffing, mobile delivery capacity, logistics, and administrative throughput—scale with both the population targeted and the magnitude of delay reduction.

\paragraph{Heterogeneous delivery difficulty.}
In practice, operational costs may differ across groups due to geography, access barriers, or administrative complexity. The framework therefore allows $c_g$ to vary. The baseline simulations consider:
\begin{itemize}[leftmargin=*]
\item \textbf{Uniform cost:} $c_g=c$ for all $g$, providing a benchmark comparison of speed-ups versus incentives.
\item \textbf{Differential costs:} higher $c_g$ for subgroups that are harder to reach (e.g., rural populations) or require more intensive delivery modalities, to reflect realistic implementation constraints.
\end{itemize}

\paragraph{Diminishing returns and feasibility bounds.}
Although \eqref{eq:speedup_cost} is linear, large speed-ups may become increasingly expensive as programmes approach capacity limits. To reflect this in a parsimonious way, I impose feasibility constraints:
\begin{equation}
0 \le \Delta t_g \le t_g,
\end{equation}
so that delay cannot be reduced below zero, and I examine alternative cost schedules in sensitivity analyses, including convex costs where marginal cost increases with the magnitude of acceleration:
\begin{equation}
\text{Cost}^{S} \;=\; \sum_{g} N_g \cdot \left(c_{1,g}\Delta t_g + c_{2,g}\Delta t_g^{2}\right),
\label{eq:convex_cost}
\end{equation}
with $c_{2,g}\ge 0$. This extension captures diminishing returns to capacity expansion when the easiest efficiency gains are exhausted first.

\paragraph{Benchmarking and interpretation.}
The unit cost parameter $c_g$ should be interpreted as an all-in programme cost per person-month of delay reduction, encompassing the operational margin relevant for timeliness (e.g., incremental staffing, extended hours, additional delivery points, transportation and cold-chain resources, IT throughput, and queue management). Rather than claim a single ``true'' value, the analysis reports policy comparisons across a range of $c_g$ values. The key output is the threshold condition implied by the policy rules: speed-ups are preferred when the marginal operational cost of saving one month is below the subgroup-specific shadow price of waiting (MWTA), and compensation is preferred otherwise.

\paragraph{Appendix note.}
The baseline values used for $c_g$, alternative convex specifications, and additional sensitivity scenarios are documented in Chapter~\ref{ch:supp_essay3}.

\subsection{Budget scenarios}
\label{subsec:essay3_budget_scenarios}

Vaccination rollouts operate under explicit fiscal constraints and short-run capacity limits. To mirror this planning environment, the simulations evaluate policy performance under alternative budget envelopes and ask how resources should be allocated between operational speed-ups and monetary compensation. Because budget levels differ by programme scale and institutional context, I report results using a transparent per-person budgeting metric and then translate to total budgets using the targeted population size.

\paragraph{Budget constraint.}
Let $B$ denote the total programme budget available for timeliness-related interventions over the rollout window under consideration. Total spending includes both operational acceleration and incentive outlays:
\begin{equation}
\sum_{g} \Big[ N_g \cdot C^{S}_g(\Delta t_g) \;+\; N_g \cdot \Delta m_g \Big] \;\le\; B,
\label{eq:budget_constraint}
\end{equation}
where $C^{S}_g(\Delta t_g)$ is defined in Section~\ref{subsec:essay3_operational_costs} and $\Delta m_g$ denotes the incentive paid per vaccinated individual (or per eligible individual, depending on programme implementation). In the baseline specification, $C^{S}_g(\Delta t_g)=c_g \Delta t_g$.

\paragraph{Per-capita budget benchmarking.}
For interpretability, I parameterize budgets in RMB \emph{per eligible person}, $\bar{b}$, so that:
\begin{equation}
B = \bar{b}\cdot \sum_{g} N_g.
\end{equation}
This allows results to be communicated as: ``given $\bar{b}$ RMB per person, what coverage and equity outcomes are achievable?'' It also facilitates comparison across programme scales.

\paragraph{Scenario grid.}
The main analyses consider three budget tiers intended to represent low, moderate, and high fiscal capacity for timeliness policies:
\begin{itemize}[leftmargin=*]
\item \textbf{Tight budget:} limited ability to either expand capacity or offer meaningful incentives (small $\bar{b}$).
\item \textbf{Moderate budget:} sufficient resources for either targeted speed-ups or moderate incentives.
\item \textbf{Expanded budget:} ability to combine speed-ups with incentives and to pursue equity-sensitive targeting at larger scale.
\end{itemize}
Rather than fix a single numerical value for each tier, I report results over a grid of $\bar{b}$ values and highlight thresholds at which the optimal instrument mix shifts from compensation-dominant to speed-up-dominant strategies.

\paragraph{Budget objectives.}
Two decision objectives are evaluated under each budget scenario:
\begin{enumerate}[leftmargin=*]
\item \textbf{Maximize coverage (with equity tracking).} Choose $(\Delta t_g,\Delta m_g)$ to maximize predicted uptake, reporting distributional outcomes across subgroups.
\item \textbf{Maximize equity-weighted welfare.} Choose $(\Delta t_g,\Delta m_g)$ to maximize an equity-weighted objective using subgroup weights $\omega_g$, subject to \eqref{eq:budget_constraint}.
\end{enumerate}

\paragraph{Interpretation for planners.}
These budget scenarios are intended to translate the conceptual policy rules into actionable guidance. Under tight budgets, the framework clarifies whether modest, targeted acceleration produces larger returns than small incentive payments spread thinly. Under larger budgets, it identifies when combining targeted speed-ups with calibrated incentives becomes the most effective way to meet coverage targets while reducing inequities in the burden of waiting.

\paragraph{Appendix note.}
The full budget grid, alternative assumptions about whether incentives are paid conditional on vaccination versus eligibility, and corresponding sensitivity analyses are reported in Chapter~\ref{ch:supp_essay3}.


\section{Sensitivity analysis: operational cost variations}
\label{sec:supp_essay3_sensitivity}

To address uncertainty in programme cost parameters, Table~\ref{tab:supp_sensitivity} reports the sensitivity of the optimal policy mix (speed-up versus compensation) to variations in the unit cost of reducing delay ($c$). The analysis considers deviations of $\pm 20\%$ and $\pm 50\%$ from the baseline assumption of $c=50$ RMB per month saved per person.

\begin{table}[htbp]
\centering
\caption{Sensitivity of optimal policy instrument to operational unit costs ($MWTA \approx 45$ RMB/month)}
\label{tab:supp_sensitivity}
\begin{threeparttable}
\begin{tabular}{lccc}
\toprule
Scenario & Unit Cost ($c$) & Cost vs. MWTA & Dominant Strategy \\
\midrule
Low cost (-50\%) & 25 RMB & $c \ll MWTA$ & Speed-up (Universal) \\
Low cost (-20\%) & 40 RMB & $c < MWTA$ & Speed-up (Universal) \\
\textbf{Baseline} & \textbf{50 RMB} & $\mathbf{c \approx MWTA}$ & \textbf{Mixed / Targeted} \\
High cost (+20\%) & 60 RMB & $c > MWTA$ & Compensation (Universal) \\
High cost (+50\%) & 75 RMB & $c \gg MWTA$ & Compensation (Universal) \\
\bottomrule
\end{tabular}
\begin{tablenotes}
\footnotesize
\item Notes: Dominant strategy refers to the instrument that yields higher welfare per RMB spent for the average respondent. "Mixed / Targeted" implies that speed-ups are preferred for high-MWTA subgroups while compensation is preferred for others.
\end{tablenotes}
\end{threeparttable}
\end{table}

The results confirm that the "mixed" policy recommendation in the main text---where targeting is crucial---is specific to the regime where operational costs are comparable to behavioral valuation. If operational improvements are significantly cheaper ($<40$ RMB), accelerating access is robustly superior to offering incentives. Conversely, if acceleration is costly ($>60$ RMB), financial compensation becomes the clear second-best alternative.





% Begin the bibliography:
\begin{thesisbib}
\bibliography{merged_unique_references}
\end{thesisbib}











\end{document}

